% O Sistema Geodésico Brasileiro — Parte 2: a rede altimétrica. O zero (o mar e o marégrafo), o
% nivelamento, a história dos ajustamentos e o REALT-2018, e os dois zeros do Brasil (Imbituba e
% Santana). Continua sgb_sirgas2000.tex, cujo último slide promete a Parte 2 ("A Rede Altimétrica,
% o Nivelamento e o Geoide"). Quatro aulas de 50 minutos, uma por parte. Tudo começa em linguagem
% de ensino médio, em passos pequenos, com apêndice técnico no fim. Mesmo formato, mesma paleta
% e mesmas convenções de sgb_sirgas2000.tex.
%
% Enfoque COMPLEMENTAR ao do geoide: geoide, superfícies de nível, número geopotencial, tipos de
% altitude e IHRS/IHRF já estão em geoide_e_alturas.tex; Molodensky, altitude normal, quase-geoide
% e hgeoHNOR2020 já estão em quase_geoide.tex; gravidade e gravimetria, em gravimetria.tex. Aqui
% eles só são relembrados. O que é novo: o nível médio do mar e o marégrafo; o datum vertical como
% UM valor de NMM em UMA estação; o nivelamento geométrico e o erro que cresce com a raiz da
% distância; a rede de RRNN; a história dos ajustamentos até o REALT-2018; os dois zeros.
%
% Bilíngue (AGENTS.md): o mesmo arquivo gera as duas versões.
%   Português:  latexmk -pdf sgb_altimetria.tex      -> sgb_altimetria.pdf
%   English:    latexmk -pdf sgb_altimetria_en.tex   -> sgb_altimetria_en.pdf
%   Limpar auxiliares:  latexmk -c sgb_altimetria.tex sgb_altimetria_en.tex
% Todo texto visível passa por \tr{português}{english}. Números decimais passam por \num{}, que
% troca vírgula por ponto na versão em inglês. Números derivados e coordenadas das figuras:
% sgb_altimetria_numeros.py (rodar sem argumentos, ou com --figuras).
%
% Acompanha a aula o simulador nivelamento/ (nivelamento por mínimos quadrados: sigma em mm/raiz(km),
% fechamento, erro grosseiro, Monte Carlo) e continua a série sgb_sirgas2000.tex.
%
% Fontes primárias consultadas em 2026-10-21 (PDFs lidos, não apenas resumos de busca):
% IBGE (1983), R.PR 22/1983, Quadro VI, p. 30: seção 3 mm raiz(K), linha 4 mm raiz(K),
% circuito 0,5 mm/km, erro-padrão da linha após ajuste 2 mm raiz(K).
% IBGE (2011), Ajustamento simultâneo da RAAP: pp. 2, 8, 10-12.
% IBGE (2019), REALT-2018, 2.a ed.: p. 13 (data), pp. 35-36 (contagens e resíduos),
% pp. 40-44 (linhas individuais e Amapá). 66 882 RRNN = bloco principal, não total nacional.
% Siqueira, Lima e Luz (2018), SIRGAS, slides 3-4, 9, 20-21, 31-36.
% SIRGAS GT III: sistema vertical, IHRS e densificação regional do IHRF.
% Referências completas, páginas, alcance da revisão e ressalvas: sgb_altimetria_fontes.md.
% As estimativas de estudos sobre NMM, subsidência e diferença entre data permanecem
% qualificadas. A extrapolação de tendência até 2026 é ilustrativa.

\newif\ifEN
\ifdefined\aulaEN \ENtrue \fi

\documentclass[aspectratio=169,11pt]{beamer}

\usetheme[progressbar=frametitle, sectionpage=progressbar, numbering=fraction]{metropolis}
\metroset{block=fill}

\usepackage[T1]{fontenc}
\usepackage{lmodern}
\ifEN
  \usepackage[brazilian,english]{babel}
\else
  \usepackage[english,brazilian]{babel}
\fi
\usepackage{amsmath, amssymb, mathtools}
\usepackage{booktabs}
\usepackage{array}
\usepackage{siunitx}
\usepackage[most]{tcolorbox}
\usepackage{xspace}
\usepackage{tikz}
\usetikzlibrary{arrows.meta, calc, positioning, decorations.pathreplacing,
                decorations.pathmorphing, decorations.markings, patterns, babel}

% ------------------------------------------------------------------ idioma / language
\newcommand{\tr}[2]{\ifEN#2\else#1\fi}
\ifEN
  \sisetup{output-decimal-marker={.}, group-separator={\,}, group-minimum-digits=5,
           group-digits=integer}
\else
  \sisetup{output-decimal-marker={,}, group-separator={\,}, group-minimum-digits=5,
           group-digits=integer}
\fi

% coluna de tabela sem justificar (evita espaços largos em colunas estreitas)
\newcolumntype{L}[1]{>{\raggedright\arraybackslash}p{#1}}

% ------------------------------------------------------------------ paleta (a mesma das outras aulas)
\definecolor{mStone}{HTML}{1C1917}
\definecolor{mStoneMid}{HTML}{78716C}
\definecolor{mStoneLight}{HTML}{F5F5F4}
\definecolor{mTeal}{HTML}{0D9488}
\definecolor{mRose}{HTML}{E11D48}
\definecolor{mAmber}{HTML}{D97706}
\definecolor{mIndigo}{HTML}{4F46E5}
\definecolor{mSky}{HTML}{0284C7}

% As ideias que se repetem em toda figura desta aula
\colorlet{corSist}{mIndigo}    % o SISTEMA / a regra escrita — datum, definição, convenção
\colorlet{corFrame}{mTeal}     % o FRAME / o medido — RRNN, linhas de nivelamento, a rede
\colorlet{corTempo}{mAmber}    % o TEMPO — época, ano, ciclo, história
\colorlet{corErro}{mRose}      % a DIFERENÇA — erro, fechamento, deslocamento
\colorlet{corMar}{mSky}        % o MAR / o geoide — nível médio do mar, a água

\setbeamercolor{normal text}{fg=mStone, bg=white}
\setbeamercolor{palette primary}{fg=white, bg=mStone}
\setbeamercolor{frametitle}{fg=white, bg=mStone}
\setbeamercolor{alerted text}{fg=mRose}
\setbeamercolor{example text}{fg=mTeal}
\setbeamercolor{progress bar}{fg=mTeal, bg=mTeal!20}
\setbeamercolor{progress bar in head/foot}{fg=mTeal, bg=mTeal!20}
\setbeamercolor{progress bar in section page}{fg=mTeal, bg=mTeal!20}
\setbeamercolor{title separator}{fg=mTeal}
\setbeamercolor{block title}{fg=white, bg=mTeal}
\setbeamercolor{block body}{fg=mStone, bg=mTeal!8}

% ------------------------------------------------------------------ caixas
\newtcolorbox{ideia}[1][]{enhanced, colback=mTeal!7, colframe=mTeal, boxrule=0.6pt, arc=2pt,
  left=6pt, right=6pt, top=4pt, bottom=4pt, fonttitle=\bfseries\small, coltitle=white, #1}
\newtcolorbox{cuidado}[1][]{enhanced, colback=mAmber!8, colframe=mAmber, boxrule=0.6pt, arc=2pt,
  left=6pt, right=6pt, top=4pt, bottom=4pt, fonttitle=\bfseries\small, coltitle=white, #1}
\newtcolorbox{armadilha}[1][]{enhanced, colback=mRose!6, colframe=mRose, boxrule=0.6pt, arc=2pt,
  left=6pt, right=6pt, top=4pt, bottom=4pt, fonttitle=\bfseries\small, coltitle=white, #1}
\newtcolorbox{regra}[1][]{enhanced, colback=mIndigo!7, colframe=mIndigo, boxrule=0.6pt, arc=2pt,
  left=6pt, right=6pt, top=4pt, bottom=4pt, fonttitle=\bfseries\small, coltitle=white, #1}

% ------------------------------------------------------------------ estilos TikZ
\tikzset{
  seta/.style={-{Stealth[length=2.4mm]}, thick},
  seta2/.style={{Stealth[length=2mm]}-{Stealth[length=2mm]}},
  rotulo/.style={font=\scriptsize, text=mStone},
  caixa/.style={draw, rounded corners=3pt, thick, align=center, font=\scriptsize},
  mar/.style={corMar, very thick},
  rede/.style={corFrame!70, thick},
  marco/.pic={
    \draw[thick, draw=corFrame, fill=corFrame!25] (-0.1,0) -- (0.1,0) -- (0,0.18) -- cycle;
  },
  pessoa/.pic={
    \fill (0,0.40) circle (0.075);
    \draw[line width=0.9pt] (0,0.33) -- (0,0.13);
    \draw[line width=0.9pt] (0,0.13) -- (-0.08,0) (0,0.13) -- (0.08,0);
    \draw[line width=0.9pt] (-0.11,0.26) -- (0.11,0.26);
  },
  antena/.pic={
    \draw[thick, draw=corFrame] (0,0) -- (0,0.16);
    \draw[thick, draw=corFrame, fill=corFrame!25] (-0.11,0.16) arc (180:0:0.11) -- cycle;
  },
}

\newcommand{\bnum}[1]{\ifmmode\textcolor{mTeal}{\mathbf{#1}}\else\textbf{\textcolor{mTeal}{#1}}\fi}
\newcommand{\sist}[1]{\textcolor{corSist}{\textbf{#1}}}
\newcommand{\fram}[1]{\textcolor{corFrame}{\textbf{#1}}}
\newcommand{\tmpo}[1]{\textcolor{corTempo}{\textbf{#1}}}
\newcommand{\marr}[1]{\textcolor{corMar}{\textbf{#1}}}
\newcommand{\errr}[1]{\textcolor{corErro}{\textbf{#1}}}
\newcommand{\RN}{RN\xspace}
\newcommand{\RRNN}{RRNN\xspace}
\newcommand{\sK}{\ensuremath{\sqrt{K}}\xspace}
\newcommand{\mmsK}{\ensuremath{\,\mathrm{mm}\sqrt{K}}\xspace}
\newcommand{\mmkm}{\ensuremath{\,\mathrm{mm}/\sqrt{\mathrm{km}}}\xspace}
\newcommand{\asec}{\ensuremath{{}^{\prime\prime}}}
\newcommand{\gr}{\ensuremath{{}^{\circ}}}
\newcommand{\ud}[1]{\ensuremath{\,\mathrm{#1}}\xspace}
\newcommand{\geopu}{\ensuremath{\,\mathrm{m^2/s^2}}\xspace}
\newcommand{\fonteslide}[1]{\par\smallskip{\tiny\color{mStoneMid}#1}}

% ---- coordenadas das figuras: sgb_altimetria_numeros.py --figuras (série sintética: ilustrativa) ----
\newcommand{\mareDia}[1]{\draw[#1] plot[smooth] coordinates {(0.00,1.00) (0.07,1.16) (0.15,1.25) (0.22,1.27) (0.30,1.21) (0.38,1.09) (0.45,0.91) (0.53,0.68) (0.60,0.42) (0.67,0.15) (0.75,-0.12) (0.82,-0.38) (0.90,-0.59) (0.97,-0.76) (1.05,-0.87) (1.12,-0.91) (1.20,-0.89) (1.27,-0.80) (1.35,-0.65) (1.43,-0.46) (1.50,-0.24) (1.57,0.00) (1.65,0.24) (1.72,0.47) (1.80,0.66) (1.88,0.80) (1.95,0.89) (2.02,0.91) (2.10,0.87) (2.17,0.76) (2.25,0.60) (2.32,0.38) (2.40,0.13) (2.48,-0.15) (2.55,-0.42) (2.62,-0.68) (2.70,-0.91) (2.77,-1.09) (2.85,-1.22) (2.92,-1.27) (3.00,-1.26) (3.07,-1.17) (3.15,-1.01) (3.23,-0.79) (3.30,-0.53) (3.38,-0.23) (3.45,0.08) (3.52,0.38) (3.60,0.67) (3.67,0.92) (3.75,1.11) (3.82,1.25) (3.90,1.31) (3.97,1.30) (4.05,1.21) (4.12,1.06) (4.20,0.86) (4.27,0.61) (4.35,0.34) (4.42,0.05) (4.50,-0.22) (4.58,-0.47) (4.65,-0.68) (4.72,-0.84) (4.80,-0.93) (4.88,-0.96) (4.95,-0.92) (5.02,-0.82) (5.10,-0.66) (5.17,-0.46) (5.25,-0.23) (5.33,0.02) (5.40,0.26) (5.47,0.47) (5.55,0.66) (5.62,0.79) (5.70,0.86) (5.77,0.87) (5.85,0.81) (5.92,0.69) (6.00,0.52) (6.08,0.29) (6.15,0.04) (6.22,-0.22) (6.30,-0.49) (6.38,-0.73) (6.45,-0.94) (6.52,-1.10) (6.60,-1.20) (6.67,-1.23) (6.75,-1.18) (6.83,-1.07) (6.90,-0.89) (6.97,-0.65) (7.05,-0.38) (7.12,-0.07) (7.20,0.24)};}
\newcommand{\mareSizigia}[1]{\draw[#1] plot coordinates {(0.00,1.30) (0.02,0.94) (0.04,0.05) (0.06,-0.87) (0.07,-1.30) (0.09,-1.00) (0.11,-0.14) (0.13,0.79) (0.15,1.28) (0.17,1.06) (0.19,0.24) (0.21,-0.71) (0.22,-1.26) (0.24,-1.10) (0.26,-0.33) (0.28,0.63) (0.30,1.22) (0.32,1.14) (0.34,0.41) (0.36,-0.54) (0.38,-1.18) (0.39,-1.16) (0.41,-0.50) (0.43,0.44) (0.45,1.13) (0.47,1.18) (0.49,0.57) (0.51,-0.35) (0.53,-1.07) (0.54,-1.19) (0.56,-0.64) (0.58,0.25) (0.60,1.00) (0.62,1.18) (0.64,0.71) (0.66,-0.16) (0.67,-0.92) (0.69,-1.17) (0.71,-0.76) (0.73,0.06) (0.75,0.84) (0.77,1.14) (0.79,0.81) (0.81,0.03) (0.83,-0.76) (0.84,-1.11) (0.86,-0.85) (0.88,-0.12) (0.90,0.67) (0.92,1.07) (0.94,0.87) (0.96,0.20) (0.97,-0.57) (0.99,-1.02) (1.01,-0.89) (1.03,-0.28) (1.05,0.48) (1.07,0.96) (1.09,0.90) (1.11,0.35) (1.12,-0.38) (1.14,-0.89) (1.16,-0.90) (1.18,-0.41) (1.20,0.29) (1.22,0.81) (1.24,0.88) (1.26,0.47) (1.27,-0.19) (1.29,-0.73) (1.31,-0.86) (1.33,-0.52) (1.35,0.10) (1.37,0.65) (1.39,0.83) (1.41,0.56) (1.42,-0.01) (1.44,-0.56) (1.46,-0.79) (1.48,-0.58) (1.50,-0.07) (1.52,0.47) (1.54,0.74) (1.56,0.60) (1.57,0.15) (1.59,-0.37) (1.61,-0.68) (1.63,-0.61) (1.65,-0.22) (1.67,0.28) (1.69,0.61) (1.71,0.61) (1.72,0.29) (1.74,-0.18) (1.76,-0.54) (1.78,-0.60) (1.80,-0.34) (1.82,0.09) (1.84,0.46) (1.86,0.58) (1.88,0.39) (1.89,0.01) (1.91,-0.37) (1.93,-0.55) (1.95,-0.43) (1.97,-0.09) (1.99,0.29) (2.01,0.51) (2.02,0.46) (2.04,0.18) (2.06,-0.19) (2.08,-0.46) (2.10,-0.48) (2.12,-0.26) (2.14,0.10) (2.16,0.40) (2.17,0.49) (2.19,0.33) (2.21,-0.00) (2.23,-0.33) (2.25,-0.49) (2.27,-0.40) (2.29,-0.09) (2.31,0.26) (2.32,0.48) (2.34,0.46) (2.36,0.19) (2.38,-0.18) (2.40,-0.46) (2.42,-0.51) (2.44,-0.28) (2.46,0.10) (2.48,0.43) (2.49,0.55) (2.51,0.37) (2.53,-0.01) (2.55,-0.40) (2.57,-0.58) (2.59,-0.45) (2.61,-0.08) (2.62,0.35) (2.64,0.60) (2.66,0.53) (2.68,0.18) (2.70,-0.29) (2.72,-0.61) (2.74,-0.61) (2.76,-0.27) (2.77,0.23) (2.79,0.61) (2.81,0.67) (2.83,0.37) (2.85,-0.15) (2.87,-0.60) (2.89,-0.73) (2.91,-0.46) (2.93,0.08) (2.94,0.59) (2.96,0.79) (2.98,0.55) (3.00,0.01) (3.02,-0.56) (3.04,-0.83) (3.06,-0.64) (3.07,-0.10) (3.09,0.52) (3.11,0.86) (3.13,0.73) (3.15,0.19) (3.17,-0.47) (3.19,-0.88) (3.21,-0.81) (3.22,-0.28) (3.24,0.42) (3.26,0.90) (3.28,0.88) (3.30,0.38) (3.32,-0.35) (3.34,-0.90) (3.36,-0.95) (3.38,-0.47) (3.39,0.28) (3.41,0.89) (3.43,1.01) (3.45,0.57) (3.47,-0.20) (3.49,-0.88) (3.51,-1.07) (3.52,-0.66) (3.54,0.12) (3.56,0.85) (3.58,1.11) (3.60,0.75) (3.62,-0.03) (3.64,-0.81) (3.66,-1.14) (3.68,-0.84) (3.69,-0.06) (3.71,0.77) (3.73,1.17) (3.75,0.92) (3.77,0.15) (3.79,-0.71) (3.81,-1.18) (3.82,-0.99) (3.84,-0.25) (3.86,0.65) (3.88,1.19) (3.90,1.06) (3.92,0.34) (3.94,-0.58) (3.96,-1.18) (3.97,-1.12) (3.99,-0.44) (4.01,0.50) (4.03,1.16) (4.05,1.18) (4.07,0.53) (4.09,-0.42) (4.11,-1.14) (4.12,-1.22) (4.14,-0.62) (4.16,0.33) (4.18,1.10) (4.20,1.26) (4.22,0.71) (4.24,-0.24) (4.26,-1.06) (4.27,-1.28) (4.29,-0.79) (4.31,0.15) (4.33,1.01) (4.35,1.30) (4.37,0.87) (4.39,-0.05) (4.41,-0.94) (4.42,-1.30) (4.44,-0.94) (4.46,-0.04) (4.48,0.87) (4.50,1.30) (4.52,1.00) (4.54,0.14) (4.56,-0.80) (4.58,-1.29) (4.59,-1.05) (4.61,-0.23) (4.63,0.72) (4.65,1.26) (4.67,1.10) (4.69,0.32) (4.71,-0.63) (4.72,-1.23) (4.74,-1.13) (4.76,-0.41) (4.78,0.54) (4.80,1.18) (4.82,1.16) (4.84,0.49) (4.86,-0.45) (4.88,-1.13) (4.89,-1.18) (4.91,-0.57) (4.93,0.35) (4.95,1.07) (4.97,1.19) (4.99,0.64) (5.01,-0.26) (5.02,-1.00) (5.04,-1.18) (5.06,-0.70) (5.08,0.16) (5.10,0.93) (5.12,1.17) (5.14,0.76) (5.16,-0.07) (5.17,-0.85) (5.19,-1.15) (5.21,-0.81) (5.23,-0.02) (5.25,0.76) (5.27,1.11) (5.29,0.84) (5.31,0.11) (5.33,-0.67) (5.34,-1.07) (5.36,-0.87) (5.38,-0.20) (5.40,0.58) (5.42,1.02) (5.44,0.89) (5.46,0.27) (5.48,-0.48) (5.49,-0.96) (5.51,-0.90) (5.53,-0.34) (5.55,0.39) (5.57,0.89) (5.59,0.90) (5.61,0.41) (5.62,-0.29) (5.64,-0.82) (5.66,-0.89) (5.68,-0.47) (5.70,0.20) (5.72,0.74) (5.74,0.86) (5.76,0.51) (5.77,-0.11) (5.79,-0.65) (5.81,-0.83) (5.83,-0.55) (5.85,0.02) (5.87,0.57) (5.89,0.79) (5.91,0.58) (5.92,0.07) (5.94,-0.47) (5.96,-0.74) (5.98,-0.60) (6.00,-0.14) (6.02,0.38) (6.04,0.68) (6.06,0.61) (6.07,0.22) (6.09,-0.28) (6.11,-0.62) (6.13,-0.61) (6.15,-0.28) (6.17,0.19) (6.19,0.54) (6.21,0.60) (6.23,0.34) (6.24,-0.09) (6.26,-0.46) (6.28,-0.58) (6.30,-0.39) (6.32,0.00) (6.34,0.38) (6.36,0.55) (6.38,0.43) (6.39,0.09) (6.41,-0.29) (6.43,-0.51) (6.45,-0.46) (6.47,-0.17) (6.49,0.20) (6.51,0.46) (6.52,0.48) (6.54,0.25) (6.56,-0.10) (6.58,-0.40) (6.60,-0.49) (6.62,-0.33) (6.64,0.01) (6.66,0.34) (6.67,0.49) (6.69,0.39) (6.71,0.09) (6.73,-0.27) (6.75,-0.48) (6.77,-0.45) (6.79,-0.18) (6.81,0.19) (6.82,0.47) (6.84,0.50) (6.86,0.27) (6.88,-0.10) (6.90,-0.44) (6.92,-0.54) (6.94,-0.36) (6.96,0.02) (6.98,0.40) (6.99,0.58) (7.01,0.45) (7.03,0.08) (7.05,-0.35) (7.07,-0.60) (7.09,-0.53) (7.11,-0.17) (7.12,0.29) (7.14,0.61) (7.16,0.60) (7.18,0.26) (7.20,-0.23) (7.22,-0.61) (7.24,-0.67) (7.26,-0.36) (7.27,0.16) (7.29,0.61) (7.31,0.73) (7.33,0.46) (7.35,-0.08) (7.37,-0.59) (7.39,-0.78) (7.41,-0.55) (7.42,-0.00) (7.44,0.56) (7.46,0.83) (7.48,0.64) (7.50,0.09) (7.52,-0.52) (7.54,-0.86) (7.56,-0.72) (7.57,-0.18) (7.59,0.48) (7.61,0.88) (7.63,0.81) (7.65,0.28) (7.67,-0.42) (7.69,-0.90) (7.71,-0.88) (7.73,-0.37) (7.74,0.36) (7.76,0.90) (7.78,0.95) (7.80,0.47) (7.82,-0.29) (7.84,-0.89) (7.86,-1.01) (7.88,-0.56) (7.89,0.21) (7.91,0.88) (7.93,1.06) (7.95,0.65) (7.97,-0.13) (7.99,-0.85) (8.01,-1.11) (8.03,-0.75) (8.04,0.04) (8.06,0.81) (8.08,1.14) (8.10,0.83) (8.12,0.05) (8.14,-0.77) (8.16,-1.17) (8.17,-0.91) (8.19,-0.14) (8.21,0.71) (8.23,1.18) (8.25,0.99) (8.27,0.24) (8.29,-0.65) (8.31,-1.19) (8.32,-1.06) (8.34,-0.34) (8.36,0.58) (8.38,1.18) (8.40,1.12) (8.42,0.43) (8.44,-0.51) (8.46,-1.17) (8.47,-1.17) (8.49,-0.52) (8.51,0.42) (8.53,1.14) (8.55,1.22) (8.57,0.61) (8.59,-0.34) (8.61,-1.11) (8.62,-1.26) (8.64,-0.70) (8.66,0.25) (8.68,1.06) (8.70,1.28) (8.72,0.78) (8.74,-0.15) (8.76,-1.01) (8.78,-1.30) (8.79,-0.86) (8.81,0.06) (8.83,0.95) (8.85,1.30) (8.87,0.93) (8.89,0.04) (8.91,-0.88) (8.92,-1.30) (8.94,-0.99) (8.96,-0.13) (8.98,0.80) (9.00,1.29)};}
\newcommand{\mareEnvSup}[1]{\draw[#1] plot[smooth] coordinates {(0.00,1.30) (0.07,1.30) (0.15,1.30) (0.22,1.29) (0.30,1.28) (0.38,1.27) (0.45,1.25) (0.53,1.23) (0.60,1.21) (0.67,1.18) (0.75,1.15) (0.83,1.12) (0.90,1.09) (0.97,1.05) (1.05,1.02) (1.12,0.98) (1.20,0.94) (1.27,0.90) (1.35,0.85) (1.42,0.81) (1.50,0.77) (1.57,0.72) (1.65,0.68) (1.72,0.64) (1.80,0.61) (1.88,0.57) (1.95,0.54) (2.02,0.52) (2.10,0.50) (2.17,0.50) (2.25,0.50) (2.32,0.50) (2.40,0.52) (2.48,0.54) (2.55,0.57) (2.62,0.60) (2.70,0.64) (2.77,0.68) (2.85,0.72) (2.93,0.76) (3.00,0.81) (3.07,0.85) (3.15,0.89) (3.22,0.93) (3.30,0.98) (3.38,1.01) (3.45,1.05) (3.52,1.09) (3.60,1.12) (3.68,1.15) (3.75,1.18) (3.82,1.20) (3.90,1.23) (3.97,1.25) (4.05,1.26) (4.12,1.28) (4.20,1.29) (4.27,1.30) (4.35,1.30) (4.42,1.30) (4.50,1.30) (4.58,1.30) (4.65,1.29) (4.72,1.28) (4.80,1.27) (4.88,1.25) (4.95,1.23) (5.02,1.21) (5.10,1.18) (5.17,1.15) (5.25,1.12) (5.33,1.09) (5.40,1.06) (5.48,1.02) (5.55,0.98) (5.62,0.94) (5.70,0.90) (5.77,0.86) (5.85,0.81) (5.92,0.77) (6.00,0.73) (6.07,0.69) (6.15,0.65) (6.23,0.61) (6.30,0.57) (6.38,0.55) (6.45,0.52) (6.52,0.51) (6.60,0.50) (6.67,0.50) (6.75,0.50) (6.82,0.52) (6.90,0.54) (6.98,0.57) (7.05,0.60) (7.12,0.64) (7.20,0.68) (7.27,0.72) (7.35,0.76) (7.42,0.81) (7.50,0.85) (7.57,0.89) (7.65,0.93) (7.73,0.97) (7.80,1.01) (7.88,1.05) (7.95,1.08) (8.03,1.12) (8.10,1.15) (8.17,1.18) (8.25,1.20) (8.32,1.23) (8.40,1.25) (8.47,1.26) (8.55,1.28) (8.62,1.29) (8.70,1.30) (8.78,1.30) (8.85,1.30) (8.92,1.30) (9.00,1.30)};}
\newcommand{\mareEnvInf}[1]{\draw[#1] plot[smooth] coordinates {(0.00,-1.30) (0.07,-1.30) (0.15,-1.30) (0.22,-1.29) (0.30,-1.28) (0.38,-1.27) (0.45,-1.25) (0.53,-1.23) (0.60,-1.21) (0.67,-1.18) (0.75,-1.15) (0.83,-1.12) (0.90,-1.09) (0.97,-1.05) (1.05,-1.02) (1.12,-0.98) (1.20,-0.94) (1.27,-0.90) (1.35,-0.85) (1.42,-0.81) (1.50,-0.77) (1.57,-0.72) (1.65,-0.68) (1.72,-0.64) (1.80,-0.61) (1.88,-0.57) (1.95,-0.54) (2.02,-0.52) (2.10,-0.50) (2.17,-0.50) (2.25,-0.50) (2.32,-0.50) (2.40,-0.52) (2.48,-0.54) (2.55,-0.57) (2.62,-0.60) (2.70,-0.64) (2.77,-0.68) (2.85,-0.72) (2.93,-0.76) (3.00,-0.81) (3.07,-0.85) (3.15,-0.89) (3.22,-0.93) (3.30,-0.98) (3.38,-1.01) (3.45,-1.05) (3.52,-1.09) (3.60,-1.12) (3.68,-1.15) (3.75,-1.18) (3.82,-1.20) (3.90,-1.23) (3.97,-1.25) (4.05,-1.26) (4.12,-1.28) (4.20,-1.29) (4.27,-1.30) (4.35,-1.30) (4.42,-1.30) (4.50,-1.30) (4.58,-1.30) (4.65,-1.29) (4.72,-1.28) (4.80,-1.27) (4.88,-1.25) (4.95,-1.23) (5.02,-1.21) (5.10,-1.18) (5.17,-1.15) (5.25,-1.12) (5.33,-1.09) (5.40,-1.06) (5.48,-1.02) (5.55,-0.98) (5.62,-0.94) (5.70,-0.90) (5.77,-0.86) (5.85,-0.81) (5.92,-0.77) (6.00,-0.73) (6.07,-0.69) (6.15,-0.65) (6.23,-0.61) (6.30,-0.57) (6.38,-0.55) (6.45,-0.52) (6.52,-0.51) (6.60,-0.50) (6.67,-0.50) (6.75,-0.50) (6.82,-0.52) (6.90,-0.54) (6.98,-0.57) (7.05,-0.60) (7.12,-0.64) (7.20,-0.68) (7.27,-0.72) (7.35,-0.76) (7.42,-0.81) (7.50,-0.85) (7.57,-0.89) (7.65,-0.93) (7.73,-0.97) (7.80,-1.01) (7.88,-1.05) (7.95,-1.08) (8.03,-1.12) (8.10,-1.15) (8.17,-1.18) (8.25,-1.20) (8.32,-1.23) (8.40,-1.25) (8.47,-1.26) (8.55,-1.28) (8.62,-1.29) (8.70,-1.30) (8.78,-1.30) (8.85,-1.30) (8.92,-1.30) (9.00,-1.30)};}
\newcommand{\declNodal}[1]{\draw[#1] plot[smooth] coordinates {(0.00,0.61) (0.15,0.80) (0.30,1.02) (0.45,1.26) (0.60,1.50) (0.75,1.74) (0.90,1.98) (1.05,2.20) (1.20,2.41) (1.35,2.58) (1.50,2.73) (1.65,2.85) (1.80,2.92) (1.95,2.96) (2.10,2.96) (2.25,2.91) (2.40,2.83) (2.55,2.71) (2.70,2.55) (2.85,2.37) (3.00,2.16) (3.15,1.93) (3.30,1.69) (3.45,1.45) (3.60,1.21) (3.75,0.98) (3.90,0.76) (4.05,0.57) (4.20,0.40) (4.35,0.26) (4.50,0.16) (4.65,0.10) (4.80,0.08) (4.95,0.10) (5.10,0.16) (5.25,0.26) (5.40,0.40) (5.55,0.56) (5.70,0.76) (5.85,0.97) (6.00,1.20) (6.15,1.45) (6.30,1.69) (6.45,1.93) (6.60,2.15) (6.75,2.36) (6.90,2.55) (7.05,2.70) (7.20,2.82) (7.35,2.91) (7.50,2.96) (7.65,2.96) (7.80,2.93) (7.95,2.85) (8.10,2.74) (8.25,2.59) (8.40,2.41) (8.55,2.21) (8.70,1.98) (8.85,1.75) (9.00,1.51)};}
\newcommand{\eNmes}{8.3}
\newcommand{\eNano}{3.1}
\newcommand{\eNdois}{2.5}
\newcommand{\eNnove}{1.4}
\newcommand{\eNdezenove}{0.7}
\newcommand{\passeioA}[1]{\draw[#1] plot coordinates {(0.00,0.00) (0.07,-0.02) (0.14,-0.04) (0.21,-0.05) (0.28,0.03) (0.35,0.02) (0.42,-0.16) (0.49,-0.12) (0.56,-0.15) (0.63,-0.18) (0.70,-0.17) (0.77,-0.14) (0.84,0.00) (0.91,0.08) (0.98,0.09) (1.05,0.00) (1.12,-0.12) (1.19,-0.09) (1.26,0.07) (1.33,0.07) (1.40,0.06) (1.47,0.13) (1.54,-0.05) (1.61,-0.09) (1.68,-0.03) (1.75,0.08) (1.82,0.05) (1.89,0.09) (1.96,0.12) (2.03,0.22) (2.10,0.08) (2.17,0.15) (2.24,-0.03) (2.31,-0.34) (2.38,-0.42) (2.45,-0.53) (2.52,-0.42) (2.59,-0.34) (2.66,-0.49) (2.73,-0.39) (2.80,-0.51) (2.87,-0.52) (2.94,-0.55) (3.01,-0.54) (3.08,-0.44) (3.15,-0.36) (3.22,-0.32) (3.29,-0.24) (3.36,-0.19) (3.43,-0.26) (3.50,-0.35) (3.57,-0.40) (3.64,-0.34) (3.71,-0.37) (3.78,-0.09) (3.85,-0.19) (3.92,-0.32) (3.99,-0.23) (4.06,-0.06) (4.13,-0.00) (4.20,0.10) (4.27,0.27) (4.34,0.26) (4.41,0.09) (4.48,0.02) (4.55,0.14) (4.62,-0.03) (4.69,-0.03) (4.76,0.00) (4.83,-0.04) (4.90,0.05) (4.97,0.12) (5.04,0.40) (5.11,0.47) (5.18,0.40) (5.25,0.33) (5.32,0.23) (5.39,0.35) (5.46,0.28) (5.53,0.27) (5.60,0.36) (5.67,0.27) (5.74,0.24) (5.81,0.02) (5.88,-0.11) (5.95,-0.18) (6.02,-0.13) (6.09,0.01) (6.16,0.01) (6.23,0.04) (6.30,0.06) (6.37,0.19) (6.44,0.30) (6.51,0.33) (6.58,0.21) (6.65,0.32) (6.72,0.36) (6.79,0.51) (6.86,0.51) (6.93,0.74) (7.00,0.70)};}
\newcommand{\passeioB}[1]{\draw[#1] plot coordinates {(0.00,0.00) (0.07,0.19) (0.14,0.20) (0.21,0.14) (0.28,0.01) (0.35,-0.01) (0.42,0.16) (0.49,0.26) (0.56,0.34) (0.63,0.06) (0.70,0.14) (0.77,0.21) (0.84,0.14) (0.91,0.07) (0.98,0.07) (1.05,0.27) (1.12,0.15) (1.19,0.10) (1.26,0.26) (1.33,0.21) (1.40,0.16) (1.47,0.17) (1.54,0.02) (1.61,0.05) (1.68,-0.09) (1.75,0.01) (1.82,0.01) (1.89,0.29) (1.96,0.32) (2.03,0.48) (2.10,0.33) (2.17,0.31) (2.24,0.35) (2.31,0.56) (2.38,0.36) (2.45,0.48) (2.52,0.55) (2.59,0.73) (2.66,0.82) (2.73,0.82) (2.80,0.76) (2.87,0.61) (2.94,0.64) (3.01,0.61) (3.08,0.86) (3.15,0.78) (3.22,0.82) (3.29,0.63) (3.36,0.58) (3.43,0.62) (3.50,0.71) (3.57,0.89) (3.64,0.88) (3.71,0.75) (3.78,0.80) (3.85,0.87) (3.92,0.93) (3.99,0.84) (4.06,0.98) (4.13,0.99) (4.20,1.07) (4.27,1.22) (4.34,1.30) (4.41,1.33) (4.48,1.59) (4.55,1.62) (4.62,1.58) (4.69,1.60) (4.76,1.78) (4.83,1.79) (4.90,1.85) (4.97,2.00) (5.04,1.93) (5.11,1.73) (5.18,1.76) (5.25,1.79) (5.32,1.72) (5.39,1.82) (5.46,1.89) (5.53,1.78) (5.60,1.84) (5.67,1.81) (5.74,1.64) (5.81,1.69) (5.88,1.69) (5.95,1.60) (6.02,1.66) (6.09,1.72) (6.16,1.65) (6.23,1.70) (6.30,1.82) (6.37,1.73) (6.44,1.78) (6.51,1.78) (6.58,2.05) (6.65,1.83) (6.72,1.91) (6.79,1.87) (6.86,1.86) (6.93,2.09) (7.00,2.09)};}
\newcommand{\passeioC}[1]{\draw[#1] plot coordinates {(0.00,0.00) (0.07,0.27) (0.14,0.22) (0.21,0.26) (0.28,0.20) (0.35,0.12) (0.42,0.12) (0.49,0.09) (0.56,0.11) (0.63,-0.14) (0.70,0.10) (0.77,0.08) (0.84,0.29) (0.91,0.17) (0.98,0.21) (1.05,0.58) (1.12,0.48) (1.19,0.30) (1.26,0.24) (1.33,0.29) (1.40,0.38) (1.47,0.53) (1.54,0.50) (1.61,0.31) (1.68,0.25) (1.75,0.40) (1.82,0.46) (1.89,0.22) (1.96,0.22) (2.03,0.39) (2.10,0.64) (2.17,0.71) (2.24,0.75) (2.31,0.60) (2.38,0.50) (2.45,0.50) (2.52,0.56) (2.59,0.63) (2.66,0.58) (2.73,0.44) (2.80,0.34) (2.87,0.21) (2.94,0.28) (3.01,0.01) (3.08,-0.03) (3.15,0.02) (3.22,0.20) (3.29,0.21) (3.36,0.33) (3.43,0.28) (3.50,0.19) (3.57,0.11) (3.64,0.30) (3.71,0.42) (3.78,0.47) (3.85,0.87) (3.92,0.87) (3.99,0.94) (4.06,0.98) (4.13,0.95) (4.20,1.23) (4.27,1.41) (4.34,1.25) (4.41,1.20) (4.48,1.25) (4.55,1.34) (4.62,1.18) (4.69,0.91) (4.76,0.68) (4.83,0.67) (4.90,0.66) (4.97,0.70) (5.04,0.61) (5.11,0.46) (5.18,0.22) (5.25,0.26) (5.32,0.30) (5.39,0.42) (5.46,0.52) (5.53,0.49) (5.60,0.66) (5.67,0.64) (5.74,0.56) (5.81,0.50) (5.88,0.43) (5.95,0.17) (6.02,0.19) (6.09,0.22) (6.16,0.17) (6.23,0.09) (6.30,0.13) (6.37,0.34) (6.44,0.34) (6.51,0.28) (6.58,0.21) (6.65,0.20) (6.72,0.05) (6.79,0.04) (6.86,0.04) (6.93,0.27) (7.00,0.38)};}
\newcommand{\passeioD}[1]{\draw[#1] plot coordinates {(0.00,0.00) (0.07,0.12) (0.14,0.04) (0.21,0.12) (0.28,-0.02) (0.35,0.02) (0.42,0.07) (0.49,-0.02) (0.56,0.22) (0.63,0.29) (0.70,0.06) (0.77,0.12) (0.84,0.08) (0.91,0.08) (0.98,0.13) (1.05,0.18) (1.12,-0.07) (1.19,-0.21) (1.26,-0.11) (1.33,0.04) (1.40,0.27) (1.47,0.49) (1.54,0.24) (1.61,0.33) (1.68,0.10) (1.75,0.30) (1.82,0.33) (1.89,0.19) (1.96,0.43) (2.03,0.51) (2.10,0.28) (2.17,0.32) (2.24,0.23) (2.31,0.24) (2.38,0.19) (2.45,0.35) (2.52,0.18) (2.59,0.22) (2.66,0.44) (2.73,0.34) (2.80,0.36) (2.87,0.39) (2.94,0.30) (3.01,0.38) (3.08,0.39) (3.15,0.28) (3.22,0.49) (3.29,0.59) (3.36,0.57) (3.43,0.49) (3.50,0.40) (3.57,0.54) (3.64,0.53) (3.71,0.58) (3.78,0.56) (3.85,0.63) (3.92,0.62) (3.99,0.71) (4.06,0.69) (4.13,0.79) (4.20,0.87) (4.27,0.86) (4.34,0.75) (4.41,0.64) (4.48,0.55) (4.55,0.43) (4.62,0.54) (4.69,0.53) (4.76,0.72) (4.83,0.99) (4.90,1.00) (4.97,0.81) (5.04,0.76) (5.11,0.76) (5.18,0.58) (5.25,0.55) (5.32,0.58) (5.39,0.53) (5.46,0.43) (5.53,0.34) (5.60,0.56) (5.67,0.54) (5.74,0.44) (5.81,0.39) (5.88,0.51) (5.95,0.43) (6.02,0.48) (6.09,0.57) (6.16,0.66) (6.23,0.83) (6.30,0.76) (6.37,0.90) (6.44,0.80) (6.51,0.76) (6.58,0.90) (6.65,1.00) (6.72,1.00) (6.79,0.86) (6.86,0.92) (6.93,0.85) (7.00,0.74)};}
\newcommand{\passeioE}[1]{\draw[#1] plot coordinates {(0.00,0.00) (0.07,-0.08) (0.14,-0.05) (0.21,-0.30) (0.28,-0.26) (0.35,-0.24) (0.42,-0.31) (0.49,-0.18) (0.56,-0.11) (0.63,-0.18) (0.70,-0.26) (0.77,-0.18) (0.84,-0.36) (0.91,-0.40) (0.98,-0.48) (1.05,-0.47) (1.12,-0.44) (1.19,-0.44) (1.26,-0.30) (1.33,-0.28) (1.40,-0.31) (1.47,-0.46) (1.54,-0.37) (1.61,-0.31) (1.68,-0.14) (1.75,-0.23) (1.82,-0.22) (1.89,-0.23) (1.96,-0.24) (2.03,-0.24) (2.10,-0.44) (2.17,-0.35) (2.24,-0.27) (2.31,-0.13) (2.38,0.14) (2.45,0.10) (2.52,0.10) (2.59,0.10) (2.66,0.33) (2.73,0.12) (2.80,0.18) (2.87,0.00) (2.94,-0.31) (3.01,-0.55) (3.08,-0.71) (3.15,-0.59) (3.22,-0.69) (3.29,-0.72) (3.36,-0.86) (3.43,-0.78) (3.50,-0.92) (3.57,-0.76) (3.64,-0.88) (3.71,-0.97) (3.78,-0.95) (3.85,-0.95) (3.92,-0.78) (3.99,-0.77) (4.06,-0.91) (4.13,-0.98) (4.20,-1.05) (4.27,-0.96) (4.34,-0.91) (4.41,-0.93) (4.48,-0.82) (4.55,-0.89) (4.62,-0.90) (4.69,-0.88) (4.76,-0.70) (4.83,-0.71) (4.90,-0.75) (4.97,-0.88) (5.04,-0.79) (5.11,-0.87) (5.18,-0.82) (5.25,-1.05) (5.32,-0.93) (5.39,-1.09) (5.46,-1.18) (5.53,-1.05) (5.60,-1.21) (5.67,-1.22) (5.74,-1.21) (5.81,-1.33) (5.88,-1.17) (5.95,-0.99) (6.02,-1.23) (6.09,-1.40) (6.16,-1.31) (6.23,-1.46) (6.30,-1.51) (6.37,-1.49) (6.44,-1.41) (6.51,-1.44) (6.58,-1.43) (6.65,-1.52) (6.72,-1.36) (6.79,-1.42) (6.86,-1.38) (6.93,-1.44) (7.00,-1.43)};}
\newcommand{\passeioF}[1]{\draw[#1] plot coordinates {(0.00,0.00) (0.07,-0.02) (0.14,-0.04) (0.21,-0.16) (0.28,-0.07) (0.35,0.10) (0.42,0.13) (0.49,0.22) (0.56,0.32) (0.63,0.43) (0.70,0.47) (0.77,0.49) (0.84,0.60) (0.91,0.61) (0.98,0.49) (1.05,0.49) (1.12,0.27) (1.19,0.27) (1.26,0.32) (1.33,0.46) (1.40,0.48) (1.47,0.59) (1.54,0.59) (1.61,0.47) (1.68,0.30) (1.75,0.18) (1.82,0.02) (1.89,0.10) (1.96,0.07) (2.03,0.14) (2.10,0.15) (2.17,0.27) (2.24,0.40) (2.31,0.32) (2.38,0.29) (2.45,0.23) (2.52,0.03) (2.59,0.11) (2.66,0.26) (2.73,0.26) (2.80,0.20) (2.87,0.33) (2.94,0.05) (3.01,0.07) (3.08,0.08) (3.15,0.02) (3.22,0.18) (3.29,0.30) (3.36,0.27) (3.43,0.26) (3.50,0.21) (3.57,0.15) (3.64,-0.01) (3.71,0.16) (3.78,0.28) (3.85,0.38) (3.92,0.25) (3.99,0.50) (4.06,0.73) (4.13,0.76) (4.20,0.65) (4.27,0.54) (4.34,0.59) (4.41,0.48) (4.48,0.48) (4.55,0.61) (4.62,0.43) (4.69,0.67) (4.76,0.86) (4.83,0.71) (4.90,0.54) (4.97,0.51) (5.04,0.38) (5.11,0.35) (5.18,0.04) (5.25,0.01) (5.32,0.23) (5.39,0.10) (5.46,0.06) (5.53,-0.12) (5.60,-0.04) (5.67,-0.07) (5.74,0.17) (5.81,0.19) (5.88,0.15) (5.95,0.21) (6.02,0.15) (6.09,-0.01) (6.16,0.03) (6.23,0.01) (6.30,0.03) (6.37,0.07) (6.44,0.01) (6.51,0.06) (6.58,0.27) (6.65,0.14) (6.72,0.35) (6.79,0.32) (6.86,0.22) (6.93,0.52) (7.00,0.53)};}
\newcommand{\costaBR}[1]{\draw[#1] plot[smooth, tension=0.35] coordinates {(0.54,6.14) (0.79,5.71) (1.00,5.33) (1.17,5.34) (1.65,5.04) (2.04,4.98) (2.55,4.84) (3.02,4.61) (3.12,4.30) (3.11,4.15) (2.98,3.89) (2.55,3.36) (2.47,3.07) (2.43,2.60) (2.27,2.19) (2.01,1.78) (1.83,1.77) (1.34,1.61) (1.00,1.36) (1.00,1.02) (0.81,0.75) (0.45,0.32) (0.24,0.05)};}
\newcommand{\ptTorres}{0.81,0.75}
\newcommand{\ptImbituba}{0.98,0.92}
\newcommand{\ptSantana}{0.59,5.43}
\newcommand{\ptBelem}{1.00,5.21}
\newcommand{\ptArraialdoCabo}{2.01,1.76}
\newcommand{\ptSalvador}{2.55,3.36}
\newcommand{\ptFortaleza}{2.55,4.84}


\title{\tr{O Sistema Geodésico Brasileiro --- Parte 2: a rede altimétrica}{The Brazilian
Geodetic System --- Part 2: the vertical network}}
\subtitle{\tr{Do nível do mar ao REALT: o zero, o nivelamento, os ajustamentos e os dois data
--- quatro aulas}{From sea level to REALT: the zero, the leveling, the adjustments and the two
datums --- four classes}}
\author{}
\institute{\tr{Simulador de Nivelamento (mínimos quadrados) e de Sistemas de
Coordenadas}{Leveling (least squares) and Coordinate Systems Simulators}}
\date{}

\begin{document}

\maketitle

% ==================================================================
\begin{frame}{\tr{O que vamos ver}{What we will see}}
  \setbeamertemplate{section in toc}[sections numbered]
  \small
  \begin{columns}[T]
    \begin{column}{0.56\textwidth}
      \tableofcontents
    \end{column}
    \begin{column}{0.40\textwidth}
      \begin{ideia}[title={\tr{Como usar}{How to use}}]
        \small\tr{Quatro aulas de 50 minutos. Cada \emph{Sua vez} tem a resposta no clique
        seguinte. As fórmulas e as contas completas ficam no \bnum{apêndice}.}{Four 50-minute
        classes. Each \emph{Your turn} has the answer on the next click. Formulas and complete
        computations are in the \bnum{appendix}.}
      \end{ideia}
    \end{column}
  \end{columns}
\end{frame}

\begin{frame}{\tr{Esta é a Parte 2}{This is Part 2}}
  \begin{columns}[T]
    \begin{column}{0.47\textwidth}
      \begin{regra}[title={\tr{Parte 1 (a série anterior)}{Part 1 (the previous series)}}]
        \small\tr{A \textbf{posição}: SIRGAS2000, como o Brasil sabe \emph{onde} fica cada ponto
        --- latitude, longitude e altura \emph{sobre o elipsoide}.}{\textbf{Position}:
        SIRGAS2000, how Brazil knows \emph{where} each point is --- latitude, longitude and
        height \emph{above the ellipsoid}.}
      \end{regra}
    \end{column}
    \begin{column}{0.49\textwidth}
      \begin{ideia}[title={\tr{Parte 2 (esta série)}{Part 2 (this series)}}]
        \small\tr{A \bnum{altitude}: a distância \emph{acima do mar}. Onde fica o zero, como se
        leva o zero até a sua porta, e por que a mesma pedra já teve várias altitudes.}{The
        \bnum{altitude}: the distance \emph{above the sea}. Where the zero is, how the zero is
        carried to your door, and why the same marker has had several altitudes.}
      \end{ideia}
    \end{column}
  \end{columns}
  \bigskip
  \centering
  \tr{A Parte 2 também \bnum{começa do zero}. E ela \bnum{não repete} as aulas do geoide: o que
  elas já ensinaram só é relembrado.}{Part 2 also \bnum{starts from scratch}. And it
  \bnum{does not repeat} the geoid classes: what they already taught is only recalled.}
\end{frame}

\begin{frame}{\tr{Quatro aulas}{Four classes}}
  \begin{columns}[T]
    \begin{column}{0.47\textwidth}
      \begin{regra}[title={\tr{Aula 1 · Onde fica o zero?}{Class 1 · Where is the zero?}}]
        \small\tr{O nível médio do mar, o marégrafo e os \emph{zeros} do Brasil: Torres,
        Imbituba e Santana.}{Mean sea level, the tide gauge and Brazil's \emph{zeros}: Torres,
        Imbituba and Santana.}
      \end{regra}
      \vspace{2pt}
      \begin{regra}[title={\tr{Aula 3 · Ajustamentos e REALT}{Class 3 · Adjustments, REALT}}]
        \small\tr{Por que a mesma RN já teve várias altitudes; os ajustamentos até o
        REALT-2018.}{Why one benchmark has had several altitudes; the adjustments up to
        REALT-2018.}
      \end{regra}
    \end{column}
    \begin{column}{0.47\textwidth}
      \begin{regra}[title={\tr{Aula 2 · Lance a lance}{Class 2 · Leg by leg}}]
        \small\tr{O nivelamento geométrico, o erro que cresce com $\sqrt{K}$ e a rede de
        referências de nível.}{Geometric leveling, the error that grows with $\sqrt{K}$, and the
        network of benchmarks.}
      \end{regra}
      \vspace{2pt}
      \begin{regra}[title={\tr{Aula 4 · O zero hoje e amanhã}{Class 4 · The zero, now and next}}]
        \small\tr{Marégrafos com GNSS, os dois data que não concordam e o uso prático.}{Tide
        gauges with GNSS, the two datums that disagree, and practical use.}
      \end{regra}
    \end{column}
  \end{columns}
  \vspace{2pt}
  \centering\footnotesize
  \tr{Cada aula tem 50 minutos, começa relembrando a anterior e termina com um resumo.
  Vamos \bnum{devagar}: uma ideia por slide.}{Each class is 50 minutes long, starts by
  recalling the previous one and ends with a summary. We will go \bnum{slowly}: one idea per
  slide.}
\end{frame}
% ==================================================================
\begin{frame}[standout]
  \tr{Aula 1}{Class 1}\\[6pt]
  {\normalsize \tr{Onde fica o zero?}{Where is the zero?}}\\[4pt]
  {\scriptsize \tr{aula 1 de 4}{class 1 of 4}}
\end{frame}

\section{\tr{Onde paramos}{Where we left off}}

\begin{frame}{\tr{As perguntas de hoje}{Today's questions}}
  \begin{columns}[T]
    \begin{column}{0.58\textwidth}
      \begin{enumerate}
        \setlength{\itemsep}{8pt}
        \item \tr{Por que a altitude precisa de um ``zero'', e quem escolhe onde ele fica?}{Why
              does altitude need a ``zero'', and who chooses where it is?}
        \item \tr{Como se mede o nível do mar --- e durante quanto tempo?}{How is sea level
              measured --- and for how long?}
        \item \tr{Qual foi o zero do Brasil, e por que hoje existem dois?}{What has been
              Brazil's zero, and why are there two today?}
        \item \tr{Esse zero é ``o mar'' de verdade?}{Is that zero really ``the sea''?}
      \end{enumerate}
    \end{column}
    \begin{column}{0.38\textwidth}
      \begin{ideia}[title={\tr{Combinado}{Deal}}]
        \tr{No fim da aula, você vai saber responder às quatro --- com as suas
        palavras.}{By the end of the class you will be able to answer all four --- in your own
        words.}
        \medskip

        \tr{Elas voltam no último slide.}{They come back on the last slide.}
      \end{ideia}
    \end{column}
  \end{columns}
\end{frame}

\begin{frame}{\tr{O que a Parte 1 deixou}{What Part 1 left us}}
  \begin{columns}[T]
    \begin{column}{0.50\textwidth}
      \small
      \tr{A Parte 1 terminou com uma \bnum{posição completa}, medida por satélite no
      SIRGAS2000:}{Part 1 ended with a \bnum{complete position}, measured by satellite in
      SIRGAS2000:}
      \begin{enumerate}
        \setlength{\itemsep}{3pt}
        \item \tr{a \textbf{latitude} $\varphi$;}{the \textbf{latitude} $\varphi$;}
        \item \tr{a \textbf{longitude} $\lambda$;}{the \textbf{longitude} $\lambda$;}
        \item \tr{a \textbf{altura} $h$ --- \emph{sobre o elipsoide}.}{the \textbf{height} $h$ ---
              \emph{above the ellipsoid}.}
      \end{enumerate}
      \medskip
      \begin{cuidado}[title={\tr{Falta uma coisa}{Something is missing}}]
        \footnotesize\tr{O elipsoide é uma superfície \emph{matemática}. Quem constrói uma estrada, uma
        represa ou um esgoto quer saber \bnum{para onde a água corre}.}{The ellipsoid is a
        \emph{mathematical} surface. Whoever builds a road, a dam or a sewer wants to know
        \bnum{where the water flows}.}
      \end{cuidado}
    \end{column}
    \begin{column}{0.46\textwidth}
      \centering
      \begin{tikzpicture}[x=1cm,y=1cm,font=\scriptsize]
        \draw[corSist, very thick] (-3.0,0.2) .. controls (-1.5,-0.55) and (1.5,-0.55) .. (3.0,0.2);
        \node[corSist, below] at (0,-0.5) {\tr{elipsoide (GRS80)}{ellipsoid (GRS80)}};
        \coordinate (E) at (1.2,-0.30);
        \coordinate (P) at (1.2,2.2);
        \draw[corFrame, thick, dashed] (E) -- (P);
        \draw[corFrame, thick, <->] ($(E)+(-0.35,0)$) -- node[left, text=corFrame] {$h$} ($(P)+(-0.35,0)$);
        \fill[corFrame] (P) circle (0.07);
        \node[right, text=corFrame] at ($(P)+(0.1,0.1)$) {$P\;(\varphi,\lambda,h)$};
        \draw[thick, mStoneMid] (-3.0,1.5) .. controls (-1.5,2.4) and (0,1.4) .. (1.2,2.2)
              .. controls (1.8,2.6) and (2.4,2.3) .. (3.0,1.9);
        \node[mStoneMid, above] at (-1.6,2.25) {\tr{terreno}{ground}};
      \end{tikzpicture}\\[2pt]
      {\scriptsize\tr{$h$ é medido ao longo da normal ao elipsoide.}{$h$ is measured along the
      normal to the ellipsoid.}}
    \end{column}
  \end{columns}
\end{frame}

\begin{frame}{\tr{Lembrete: a água não corre pelo elipsoide}{Reminder: water does not flow along
the ellipsoid}}
  \begin{columns}[T]
    \begin{column}{0.58\textwidth}
      \centering
      \begin{tikzpicture}[x=0.68cm,y=0.68cm,font=\scriptsize]
        % elipsoide
        \draw[corSist, very thick] (0,0) -- (10,0);
        \node[corSist, above left] at (10.3,0.02) {\tr{elipsoide}{ellipsoid}};
        % geoide: y = 3.5 - 0.25 x  (0,5 unidade por metro)
        \draw[corMar, very thick] (0,3.5) -- (10,1.0);
        \node[corMar, above left] at (10.3,1.0) {\tr{geoide}{geoid}};
        % terreno
        \draw[thick, mStoneMid] plot[smooth] coordinates {(0,5.4) (1,6.2) (2,6.0) (3,5.5) (4,5.7)
              (5,5.4) (6,5.6) (7,5.3) (8,5.0) (9,4.6) (10,4.4)};
        \coordinate (A) at (2,6.0);
        \coordinate (B) at (8,5.0);
        \fill[corFrame] (A) circle (0.11) node[above, text=corFrame] {$A$};
        \fill[corFrame] (B) circle (0.11) node[above, text=corFrame] {$B$};
        \draw[corFrame, dashed] (A) -- (2,0) (B) -- (8,0);
        \draw[<->, corErro, thick] (1.6,0) -- node[left, text=corErro] {$h_A$} (1.6,6.0);
        \draw[<->, corErro, thick] (7.6,0) -- node[left, text=corErro] {$h_B$} (7.6,5.0);
        % superficie de nivel por B (1 m = 0,5 unidade acima de A)
        \draw[corMar, thick, dashed] (8,5.0) -- (2,6.5);
        \draw[<->, corTempo, thick] (2.45,6.0) -- (2.45,6.43);
        \node[corTempo, right] at (2.6,6.95) {\tr{$\approx1$~m acima de $A$}{$\approx1$~m above $A$}};
        \node[corMar, above, inner sep=1pt] at (5.4,5.85) {\tr{superfície de nível por $B$}{level surface through $B$}};
      \end{tikzpicture}
    \end{column}
    \begin{column}{0.40\textwidth}
      \small
      \begin{tabular}{@{}lcc@{}}
        \toprule
         & $A$ & $B$ \\
        \midrule
        $h$ (\tr{m}{m}) & \num{12.0} & \num{10.0} \\
        $N$ (\tr{m}{m}) & \num{6.0} & \num{3.0} \\
        $H=h-N$ (m) & \bnum{\num{6.0}} & \bnum{\num{7.0}} \\
        \bottomrule
      \end{tabular}
      \medskip

      \tr{Só com $h$, $A$ parece \emph{mais alto}. Mas a água corre de $B$ para $A$: $B$ é mais
      alto em $H$.}{Looking only at $h$, $A$ seems \emph{higher}. But water flows from $B$ to
      $A$: $B$ is higher in $H$.}
      \smallskip

      {\scriptsize\tr{Visto nas aulas do geoide. Figura fora de escala; números didáticos.}{Seen
      in the geoid classes. Figure not to scale; didactic numbers.}}
    \end{column}
  \end{columns}
\end{frame}

\begin{frame}{\tr{O que já vimos, e o que é novo aqui}{What we have seen, and what is new here}}
  \begin{columns}[T]
    \begin{column}{0.47\textwidth}
      \begin{regra}[title={\tr{Já visto nas aulas do geoide}{Already seen in the geoid
        classes}}]
        \small
        \begin{itemize}
          \setlength{\itemsep}{2pt}
          \item \tr{geoide e superfícies de nível;}{geoid and level surfaces;}
          \item \tr{número geopotencial $C$ e tipos de altitude;}{geopotential number $C$ and
                kinds of height;}
          \item \tr{altitude normal $H^{*}=C/\bar\gamma$, quase-geoide;}{normal height
                $H^{*}=C/\bar\gamma$, quasi-geoid;}
          \item \tr{IHRS e IHRF; o modelo hgeoHNOR2020.}{IHRS and IHRF; the hgeoHNOR2020
                model.}
        \end{itemize}
        \tr{Aqui, só \emph{relembrado}.}{Here, only \emph{recalled}.}
      \end{regra}
    \end{column}
    \begin{column}{0.49\textwidth}
      \begin{ideia}[title={\tr{Novo na Parte 2}{New in Part 2}}]
        \small
        \begin{itemize}
          \setlength{\itemsep}{2pt}
          \item \tr{o \bnum{mar} e o \bnum{marégrafo}: o nível médio;}{the \bnum{sea} and the
                \bnum{tide gauge}: mean sea level;}
          \item \tr{o \bnum{datum vertical}: um zero escolhido;}{the \bnum{vertical datum}: a
                chosen zero;}
          \item \tr{o \bnum{nivelamento} e o erro que cresce com $\sqrt{K}$;}{\bnum{leveling} and
                the error that grows with $\sqrt{K}$;}
          \item \tr{a \bnum{rede} de referências de nível;}{the \bnum{network} of benchmarks;}
          \item \tr{os \bnum{ajustamentos} até o \bnum{REALT-2018};}{the \bnum{adjustments} up to
                \bnum{REALT-2018};}
          \item \tr{os \bnum{dois zeros} do Brasil.}{Brazil's \bnum{two zeros}.}
        \end{itemize}
      \end{ideia}
    \end{column}
  \end{columns}
\end{frame}

\begin{frame}{\tr{O fio da aula}{The thread of the class}}
  \centering
  \begin{tikzpicture}[x=1cm,y=1cm,font=\small]
    \node[caixa, draw=corMar, fill=corMar!8, text width=3.7cm, minimum height=2.0cm,
          font=\small] (a) at (0,0)
      {\textbf{\tr{1.\ O mar como régua}{1.\ The sea as a ruler}}\\[2pt]
       \scriptsize\tr{nível médio, marégrafo, marés, quanto tempo de média}{mean level, tide
       gauge, tides, how long to average}};
    \node[caixa, draw=corTempo, fill=corTempo!8, text width=3.7cm, minimum height=2.0cm,
          font=\small] (b) at (4.7,0)
      {\textbf{\tr{2.\ Os zeros do Brasil}{2.\ Brazil's zeros}}\\[2pt]
       \scriptsize\tr{Torres, Imbituba, Santana}{Torres, Imbituba, Santana}};
    \node[caixa, draw=corSist, fill=corSist!8, text width=3.7cm, minimum height=2.0cm,
          font=\small] (c) at (9.4,0)
      {\textbf{\tr{3.\ O zero não é o geoide}{3.\ The zero is not the geoid}}\\[2pt]
       \scriptsize\tr{datum = um NMM, níveis de redução, o mar sobe}{datum = one MSL, reduction
       levels, the sea rises}};
    \draw[seta] (a) -- (b);
    \draw[seta] (b) -- (c);
    \node[rotulo, below=0.15cm of a] {\tr{o que é o ``mar''?}{what is ``the sea''?}};
    \node[rotulo, below=0.15cm of b] {\tr{onde e quando se escolheu?}{where and when was it
      chosen?}};
    \node[rotulo, below=0.15cm of c] {\tr{e hoje, o que isso significa?}{and today, what does
      it mean?}};
  \end{tikzpicture}
  \bigskip

  \tr{Ao final, uma ponte para a \bnum{Aula 2}: como levar o zero do litoral até a sua
  porta.}{At the end, a bridge to \bnum{Class 2}: how to carry the zero from the coast to your
  door.}
\end{frame}

% ==================================================================
\section{\tr{O mar como régua}{The sea as a ruler}}

\begin{frame}{\tr{Uma régua precisa de um zero}{A ruler needs a zero}}
  \begin{columns}[T]
    \begin{column}{0.50\textwidth}
      \centering
      \begin{tikzpicture}[x=0.92cm,y=0.92cm,font=\scriptsize]
        % terreno: P lê 34 na régua A (0,11 cm por unidade) e 29 na régua B (zero 5 unidades acima)
        \draw[thick, mStoneMid] (1.2,0) .. controls (2.0,0.5) and (2.4,3.0) .. (3.6,3.74)
              .. controls (4.4,4.1) and (4.8,2.2) .. (5.4,1.4);
        \fill[corFrame] (3.6,3.74) circle (0.08) node[above, text=corFrame] {$P$};
        % régua A
        \draw[thick, corSist] (0,0) -- (0,4.6);
        \foreach \y/\l in {0/0,1.1/10,2.2/20,3.3/30,4.4/40} {
          \draw[corSist] (0,\y) -- (0.18,\y) node[right, text=corSist] {\l};
        }
        \node[corSist, below] at (0,-0.05) {\tr{zero A}{zero A}};
        \draw[corSist, dashed] (0,3.74) -- (3.5,3.74);
        \node[corSist, above] at (1.5,3.78) {\textbf{34}};
        % régua B: zero 5 unidades mais acima (0,55 cm)
        \draw[thick, corTempo] (6.4,0.55) -- (6.4,4.6);
        \foreach \y/\l in {0.55/0,1.65/10,2.75/20,3.85/30} {
          \draw[corTempo] (6.4,\y) -- (6.22,\y) node[left, text=corTempo] {\l};
        }
        \node[corTempo, below] at (6.4,0.5) {\tr{zero B}{zero B}};
        \draw[corTempo, dashed] (3.7,3.74) -- (6.4,3.74);
        \node[corTempo, above] at (5.1,3.78) {\textbf{29}};
      \end{tikzpicture}
    \end{column}
    \begin{column}{0.46\textwidth}
      \small
      \tr{\textbf{Altitude} é a leitura de uma ``régua'' vertical que começa num \bnum{zero}.}{\textbf{Altitude} is the reading of a vertical ``ruler'' that starts at a \bnum{zero}.}
      \medskip

      \tr{O mesmo ponto $P$ lê \textcolor{corSist}{\textbf{34}} numa régua e
      \textcolor{corTempo}{\textbf{29}} em outra, só porque o zero foi posto em outro
      lugar.}{The same point $P$ reads \textcolor{corSist}{\textbf{34}} on one ruler and
      \textcolor{corTempo}{\textbf{29}} on another, only because the zero was put somewhere
      else.}
      \medskip
      \begin{ideia}[title={\tr{Moral}{Moral}}]
        \small\tr{Mudar o zero soma o \emph{mesmo número} em todas as leituras. Para um país
        inteiro, todo mundo precisa usar \bnum{o mesmo zero}. Esse zero tem nome:
        \bnum{datum vertical}.}{Changing the zero adds the \emph{same number} to every reading.
        For a whole country, everyone must use \bnum{the same zero}. That zero has a name:
        \bnum{vertical datum}.}
      \end{ideia}
    \end{column}
  \end{columns}
\end{frame}

\begin{frame}{\tr{O mar nunca fica parado}{The sea never stands still}}
  \begin{columns}[T]
    \begin{column}{0.58\textwidth}
      \footnotesize
      \renewcommand{\arraystretch}{1.5}
      \begin{tabular}{@{}L{2.4cm}L{1.5cm}L{3.3cm}@{}}
        \toprule
        \textbf{\tr{O que move a água}{What moves the water}} &
        \textbf{\tr{Escala}{Scale}} & \textbf{\tr{Efeito típico}{Typical effect}} \\
        \midrule
        \tr{ondas}{waves} & \tr{segundos}{seconds} & \tr{sobe e desce o tempo todo}{up and down
        all the time} \\
        \tr{marés}{tides} & \tr{horas}{hours} & \tr{duas altas por dia}{two highs a day} \\
        \tr{vento e pressão}{wind and pressure} & \tr{dias}{days} & \tr{empilha ou afasta a
        água}{piles up or pushes the water away} \\
        \tr{estações do ano}{seasons} & \tr{meses}{months} & \tr{água quente ocupa mais
        volume}{warm water takes more volume} \\
        \tr{nodos da Lua}{Moon's nodes} & \tr{18,6 anos}{18.6 years} & \tr{muda um pouco a
        maré}{changes the tide a little} \\
        \bottomrule
      \end{tabular}
    \end{column}
    \begin{column}{0.38\textwidth}
      \begin{cuidado}[title={\tr{Qual é ``o'' nível do mar?}{Which is ``the'' sea level?}}]
        \small\tr{Se você medir de manhã, de tarde, na lua cheia ou numa tempestade, o número
        muda. Não existe \emph{uma} altura do mar.}{If you measure in the morning, in the
        afternoon, at full moon or in a storm, the number changes. There is no single height of
        the sea.}
      \end{cuidado}
      \medskip
      \begin{ideia}[title={\tr{Saída}{Way out}}]
        \small\tr{Tirar a \bnum{média}: deixar o que sobe e desce se cancelar.}{Take the
        \bnum{average}: let what goes up and down cancel out.}
      \end{ideia}
    \end{column}
  \end{columns}
\end{frame}

\begin{frame}{\tr{O nível médio do mar}{Mean sea level}}
  \begin{columns}[T]
    \begin{column}{0.60\textwidth}
      \centering
      \begin{tikzpicture}[x=0.84cm,y=0.9cm,font=\scriptsize]
        \draw[->, thick] (-0.1,0) -- (7.6,0);
        \node[below] at (3.6,-0.4) {\tr{tempo (h)}{time (h)}};
        \draw[->, thick] (0,-1.8) -- (0,1.9) node[above] {\tr{altura da água}{water height}};
        \mareDia{corMar, very thick}
        \draw[corErro, very thick, dashed] (0,0) -- (7.4,0);
        \node[corErro, right] at (7.35,0.0) {NMM};
        \foreach \x/\l in {0/0,1.8/12,3.6/24,5.4/36,7.2/48} {
          \draw (\x,0.06) -- (\x,-0.06) node[below] {\l};
        }
      \end{tikzpicture}\\[2pt]
      {\scriptsize\tr{Curva ilustrativa de dois dias (não é um porto real).}{Illustrative
      two-day curve (not a real port).}}
    \end{column}
    \begin{column}{0.36\textwidth}
      \small
      \tr{Anota-se a altura da água \bnum{de hora em hora}, por muito tempo, e tira-se a
      média:}{The water height is recorded \bnum{every hour}, for a long time, and averaged:}
      \[
        \mathrm{NMM}=\frac{1}{n}\sum_{i=1}^{n}h_i
      \]
      \tr{A média apaga ondas, marés e vento. O que sobra é o \bnum{nível médio do mar} (em
      inglês, \emph{MSL}).}{The average erases waves, tides and wind. What remains is the
      \bnum{mean sea level} (\emph{MSL}).}
    \end{column}
  \end{columns}
\end{frame}

\begin{frame}{\tr{Como se mede: o marégrafo}{How it is measured: the tide gauge}}
  \begin{columns}[T]
    \begin{column}{0.56\textwidth}
      \centering
      \begin{tikzpicture}[x=0.86cm,y=0.86cm,font=\scriptsize]
        % terra à esquerda
        \fill[mStoneMid!22] (0,-2.0) rectangle (2.2,0.9);
        \draw[mStoneMid, thick] (2.2,-2.0) -- (2.2,0.9) -- (0,0.9);
        % mar
        \fill[corMar!14] (2.2,-2.0) rectangle (8.4,0);
        \draw[corMar, very thick, decorate, decoration={snake, amplitude=0.7pt, segment length=8pt}]
              (2.2,0) -- (4.3,0);
        \draw[corMar, very thick, decorate, decoration={snake, amplitude=0.7pt, segment length=8pt}]
              (5.4,0) -- (8.4,0);
        \draw[corMar, very thick] (4.3,0) -- (5.4,0);
        % poço de tranquilização
        \draw[thick] (4.3,-1.6) -- (4.3,3.0);
        \draw[thick] (5.4,-1.6) -- (5.4,-1.25) (5.4,-0.95) -- (5.4,3.0);
        \draw[thick] (4.3,-1.6) -- (5.4,-1.6);
        % boia, fio, roldana, registro
        \draw[thick, fill=white] (4.62,-0.16) rectangle (5.08,0.10);
        \draw (4.85,0.10) -- (4.85,2.3);
        \draw[thick, fill=white] (4.85,2.5) circle (0.2);
        \draw (5.05,2.5) -- (6.6,2.5);
        \draw[thick, fill=white, rounded corners=2pt] (6.6,2.0) rectangle (8.2,3.0);
        \node at (7.4,2.5) {\tr{registro}{recorder}};
        % RN e ligação por nivelamento
        \pic[scale=1.5] at (0.9,0.9) {marco};
        \node[corFrame, above] at (0.9,1.35) {RN};
        \draw[corFrame, thick, dashed, -{Stealth[length=2mm]}] (1.15,1.35) -- (1.15,3.4) -- (4.3,3.4);
        \node[corFrame, above, align=center] at (2.9,3.4) {\tr{nivelamento}{leveling}};
        % rótulos
        \node[above right, align=left] at (5.5,3.15) {\tr{poço de}{stilling}\\[-1pt]\tr{tranquilização}{well}};
        \node[left] at (4.25,0.0) {\tr{boia}{float}};
        \node[corErro, right, align=left] at (5.5,-1.1) {\tr{furo pequeno}{small hole}};
      \end{tikzpicture}
    \end{column}
    \begin{column}{0.42\textwidth}
      \small
      \begin{itemize}
        \setlength{\itemsep}{4pt}
        \item \tr{O cano tem um \bnum{furo pequeno}: a água entra devagar e o nível lá dentro
              sobe e desce \emph{sem as ondas}.}{The pipe has a \bnum{small hole}: the water
              enters slowly and the level inside goes up and down \emph{without the waves}.}
        \item \tr{Uma \bnum{boia} acompanha o nível; um fio leva o movimento ao
              \bnum{registro} (antes em papel, hoje em arquivo).}{A \bnum{float} follows the
              level; a wire carries the motion to the \bnum{recorder} (once on paper, now in a
              file).}
        \item \tr{O zero do aparelho é amarrado a uma \bnum{RN} em terra, por
              \textbf{nivelamento}.}{The instrument's zero is tied to a \bnum{benchmark} (RN) on
              land, by \textbf{leveling}.}
      \end{itemize}
    \end{column}
  \end{columns}
\end{frame}

\begin{frame}{\tr{Marégrafos de hoje}{Today's tide gauges}}
  \footnotesize
  \renewcommand{\arraystretch}{1.25}
  \begin{tabular}{@{}L{2.3cm}L{4.0cm}L{3.0cm}L{3.4cm}@{}}
    \toprule
    \textbf{\tr{Tipo}{Type}} & \textbf{\tr{Como mede}{How it measures}} &
    \textbf{\tr{Ponto forte}{Strength}} & \textbf{\tr{Cuidado}{Care}} \\
    \midrule
    \tr{Régua de maré}{Tide staff} & \tr{alguém lê a altura da água}{someone reads the water
    height} & \tr{simples e barato}{simple, cheap} & \tr{depende do observador}{depends on the
    observer} \\
    \tr{Boia em poço}{Float in a well} & \tr{boia, fio e registro}{float, wire, recorder} &
    \tr{filtra as ondas}{filters the waves} & \tr{o furo pode entupir}{the hole can clog} \\
    \tr{Pressão}{Pressure} & \tr{sensor no fundo pesa a coluna de água}{bottom sensor weighs the
    water column} & \tr{fácil de instalar}{easy to install} & \tr{corrigir pressão do ar e
    densidade}{correct for air pressure, density} \\
    \tr{Radar}{Radar} & \tr{mede a distância até a superfície}{measures the distance to the
    surface} & \tr{nada dentro da água}{nothing in the water} & \tr{ondas e chuva forte}{waves,
    heavy rain} \\
    \bottomrule
  \end{tabular}
  \smallskip

  \begin{ideia}[title={\tr{Na geodésia}{In geodesy}}]
    \footnotesize\tr{Uma estação para geodésia junta o sensor de nível, sensores de tempo
    (vento, pressão do ar) e uma \bnum{antena GNSS}. A próxima ideia explica por quê.}{A
    geodetic station adds weather sensors and a \bnum{GNSS antenna} to the level sensor. The
    next idea explains why.}
  \end{ideia}
\end{frame}

\begin{frame}{\tr{O chão também se mexe}{The ground moves too}}
  \begin{columns}[T]
    \begin{column}{0.52\textwidth}
      \centering
      \begin{tikzpicture}[x=0.88cm,y=0.88cm,font=\scriptsize]
        \fill[mStoneMid!22] (0,-2.0) rectangle (3.0,0.8);
        \draw[mStoneMid, thick] (3.0,-2.0) -- (3.0,0.8) -- (0,0.8);
        \fill[corMar!14] (3.0,-2.0) rectangle (7.2,0);
        \draw[corMar, very thick, decorate, decoration={snake, amplitude=0.7pt, segment length=8pt}]
              (3.0,0) -- (7.2,0);
        \draw[thick] (1.5,0.8) -- (1.5,2.0);
        \pic[scale=1.6] at (1.5,2.0) {antena};
        \node[corFrame, right] at (1.8,2.45) {GNSS};
        \draw[thick] (4.3,-1.3) -- (4.3,1.5) (4.9,-1.3) -- (4.9,1.5);
        \draw[thick, fill=white] (4.42,-0.14) rectangle (4.78,0.09);
        \node[right] at (5.0,1.25) {\tr{marégrafo}{tide gauge}};
        % setas: mar sobe, chão desce
        \draw[seta, corMar] (6.6,-0.1) -- (6.6,1.0);
        \node[corMar, left, align=right] at (6.55,0.55) {\tr{mar}{sea}\\$+\num{2}$};
        \draw[seta, corTempo] (0.55,0.45) -- (0.55,-0.75);
        \node[corTempo, right, align=left] at (0.6,-0.45) {\tr{chão}{ground}\\$-\num{3}$};
        \node[mStoneMid, below] at (3.6,-2.0) {\tr{mm/ano (números didáticos)}{mm/year (didactic numbers)}};
      \end{tikzpicture}
    \end{column}
    \begin{column}{0.44\textwidth}
      \footnotesize
      \tr{O marégrafo está \emph{preso ao chão}. Ele mede o mar \bnum{em relação ao
      chão}.}{The tide gauge is \emph{fixed to the ground}. It measures the sea \bnum{relative
      to the ground}.}
      \medskip

      \tr{Se o chão afunda, o marégrafo ``vê'' o mar subir \emph{mais}: no exemplo, o mar sobe
      \num{2} e o chão desce \num{3}, então o aparelho registra \bnum{+\num{5}}
      mm/ano.}{If the ground sinks, the gauge ``sees'' the sea rise \emph{more}: in the example
      the sea rises \num{2} and the ground drops \num{3}, so the instrument records
      \bnum{+\num{5}} mm/year.}
      \medskip
      \begin{ideia}[title={\tr{Por isso o GNSS}{That is why GNSS}}]
        \small\tr{Uma antena GNSS fixa ao lado mede o movimento vertical do chão. Assim se
        separa \bnum{o que é mar} do que é \bnum{chão}.}{A fixed GNSS antenna next to it
        measures the vertical motion of the ground. That separates \bnum{what is sea} from what
        is \bnum{ground}.}
      \end{ideia}
    \end{column}
  \end{columns}
\end{frame}

\begin{frame}{\tr{Por que a Lua manda nas marés}{Why the Moon rules the tides}}
  \begin{columns}[T]
    \begin{column}{0.54\textwidth}
      \centering
      \begin{tikzpicture}[x=0.70cm,y=0.80cm,font=\scriptsize]
        \fill[corMar!20] (0,0) ellipse (1.75 and 1.25);
        \draw[corMar, thick] (0,0) ellipse (1.75 and 1.25);
        \fill[mStoneMid!40] (0,0) circle (1.15);
        \node at (0,0) {\tr{Terra}{Earth}};
        \fill[mStoneMid] (5.6,0) circle (0.30);
        \node[below] at (5.6,-0.35) {\tr{Lua}{Moon}};
        % atração: longa (perto), média (centro), curta (longe)
        \draw[seta, corErro] (1.75,1.9) -- (4.45,1.9);
        \draw[seta, corErro] (0,1.9) -- (1.65,1.9);
        \draw[seta, corErro] (-1.75,1.9) -- (-0.35,1.9);
        \node[corErro, above] at (3.1,1.95) {\tr{mais forte}{stronger}};
        \node[corErro, above] at (0.85,1.95) {\tr{média}{mean}};
        \node[corErro, above] at (-1.05,1.95) {\tr{mais fraca}{weaker}};
        \node[corMar, below, align=center] at (0,-1.35) {\tr{a diferença estica a água dos
              \emph{dois} lados}{the difference stretches the water on \emph{both} sides}};
      \end{tikzpicture}
    \end{column}
    \begin{column}{0.42\textwidth}
      \small
      \tr{A Lua atrai a água do lado de cá um pouco \emph{mais} do que atrai a Terra, e a água do
      lado de lá um pouco \emph{menos}. Essa \bnum{diferença} estica o oceano em \bnum{dois
      altos}: um sob a Lua, outro no lado oposto.}{The Moon pulls the water on this side slightly
      \emph{more} than it pulls the Earth, and the water on the far side slightly \emph{less}.
      That \bnum{difference} stretches the ocean into \bnum{two highs}: one under the Moon,
      one on the opposite side.}
      \medskip

      \tr{A Terra gira sob esses dois altos; como a Lua também anda na órbita, o intervalo
      entre duas altas é de \bnum{\num{12.42}\,h} (e não 12).}{The Earth turns under these two
      highs; since the Moon also moves along its orbit, the interval between two high waters is
      \bnum{\num{12.42}\,h} (not 12).}
      \medskip

      {\footnotesize\tr{O Sol faz o mesmo, mais fraco (cerca de 46\,\% da Lua).}{The Sun does the
      same, but weaker (about 46\,\% of the Moon).}}
    \end{column}
  \end{columns}
\end{frame}

\begin{frame}{\tr{Marés vivas e mortas}{Spring and neap tides}}
  \begin{columns}[T]
    \begin{column}{0.40\textwidth}
      \centering
      \begin{tikzpicture}[x=0.9cm,y=0.9cm,font=\scriptsize]
        % sizígia
        \fill[corTempo] (0,1.5) circle (0.30);
        \node[above] at (0,1.85) {\tr{Sol}{Sun}};
        \fill[corMar!25] (2.2,1.5) ellipse (0.50 and 0.28);
        \draw[corMar] (2.2,1.5) ellipse (0.50 and 0.28);
        \fill[mStoneMid!40] (2.2,1.5) circle (0.2);
        \fill[mStoneMid] (3.25,1.5) circle (0.12);
        \node[right] at (3.4,1.5) {\tr{Lua}{Moon}};
        \node[corMar, below, align=center] at (1.6,1.05) {\textbf{\tr{sizígia}{spring}}};
        % quadratura
        \fill[corTempo] (0,-0.5) circle (0.30);
        \fill[corMar!25] (2.2,-0.5) ellipse (0.28 and 0.40);
        \draw[corMar] (2.2,-0.5) ellipse (0.28 and 0.40);
        \fill[mStoneMid!40] (2.2,-0.5) circle (0.2);
        \fill[mStoneMid] (2.2,0.35) circle (0.12);
        \node[right] at (2.4,0.35) {\tr{Lua}{Moon}};
        \node[corMar, below, align=center] at (1.6,-1.15) {\textbf{\tr{quadratura}{neap}}};
      \end{tikzpicture}
    \end{column}
    \begin{column}{0.56\textwidth}
      \centering
      \begin{tikzpicture}[x=0.72cm,y=0.80cm,font=\scriptsize]
        \draw[->] (-0.1,0) -- (9.5,0);
        \node[below] at (9.0,-0.05) {\tr{dias}{days}};
        \draw[->] (0,-1.7) -- (0,1.8) node[above] {\tr{altura}{height}};
        \mareSizigia{corMar!70}
        \mareEnvSup{corErro, thick, dashed}
        \mareEnvInf{corErro, thick, dashed}
        \foreach \x/\l in {0/0,1.5/5,3.0/10,4.5/15,6.0/20,7.5/25,9.0/30} {
          \draw (\x,0.05) -- (\x,-0.05) node[below, font=\tiny] {\l};
        }
        \draw[<->, corTempo, thick] (0.0,-1.55) -- node[below, text=corTempo] {\num{14.8}\,\tr{dias}{days}} (4.43,-1.55);
      \end{tikzpicture}
    \end{column}
  \end{columns}
  \smallskip

  \small
  \tr{Com Sol e Lua \bnum{alinhados} (lua nova ou cheia), as duas forças se somam: maré
  \bnum{viva} (sizígia). Em \bnum{ângulo reto} (quartos), uma atrapalha a outra: maré
  \bnum{morta} (quadratura). O vai-e-vem se repete a cada \bnum{\num{14.8}\,dias}.}{With Sun and
  Moon \bnum{aligned} (new or full moon), the two forces add up: \bnum{spring} tide. At
  \bnum{right angles} (quarters), one hinders the other: \bnum{neap} tide. The back-and-forth
  repeats every \bnum{\num{14.8}\,days}.}
\end{frame}

\begin{frame}{\tr{Ritmos mais longos}{Longer rhythms}}
  \begin{columns}[T]
    \begin{column}{0.62\textwidth}
      \footnotesize
      \renewcommand{\arraystretch}{1.5}
      \begin{tabular}{@{}L{1.6cm}L{3.3cm}L{3.0cm}@{}}
        \toprule
        \textbf{\tr{Período}{Period}} & \textbf{\tr{Quem manda}{What drives it}} &
        \textbf{\tr{O que muda}{What changes}} \\
        \midrule
        \num{12.42}\,h & \tr{a Lua dá a volta na Terra}{the Moon goes around the Earth} &
        \tr{duas altas por dia}{two high waters a day} \\
        \num{14.8}\,\tr{dias}{days} & \tr{Sol e Lua alinhados ou em ângulo reto}{Sun and Moon
        aligned or at right angles} & \tr{marés vivas e mortas}{spring and neap tides} \\
        \tr{1 ano}{1 year} & \tr{estações: calor, vento, chuva}{seasons: heat, wind, rain} &
        \tr{o NMM do verão não é o do inverno}{summer MSL is not winter MSL} \\
        \num{18.6}\,\tr{anos}{years} & \tr{a órbita da Lua ``gira'' no espaço}{the Moon's orbit
        ``turns'' in space} & \tr{muda um pouco a maré e o NMM}{changes the tide and MSL a
        little} \\
        \bottomrule
      \end{tabular}
    \end{column}
    \begin{column}{0.34\textwidth}
      \begin{cuidado}[title={\tr{Moral}{Moral}}]
        \small\tr{Um mês de dados pega só o ritmo curto. Para o nível médio \emph{de verdade},
        é preciso observar por um ciclo \bnum{inteiro} --- o mais longo é o de \bnum{18,6
        anos}.}{One month of data only catches the short rhythm. For the \emph{real} mean level
        you need to observe a \bnum{whole} cycle --- the longest is the \bnum{18.6-year}
        one.}
      \end{cuidado}
    \end{column}
  \end{columns}
\end{frame}

\begin{frame}{\tr{O ciclo de 18,6 anos}{The 18.6-year cycle}}
  \begin{columns}[T]
    \begin{column}{0.52\textwidth}
      \centering
      \begin{tikzpicture}[x=0.60cm,y=0.95cm,font=\scriptsize]
        \draw[->] (-0.1,0) -- (9.5,0);
        \draw[->] (0,-0.1) -- (0,3.5) node[above, align=center] {\tr{declinação máxima}{max.\ declination}\\\tr{da Lua (\gr)}{of the Moon (\gr)}};
        \declNodal{corSist, very thick}
        \foreach \x/\l in {0/2000,3/2010,6/2020,9/2030} {
          \draw (\x,0.05) -- (\x,-0.05) node[below] {\l};
        }
        \foreach \y/\l in {0.56/20,1.96/25} {
          \draw (0.05,\y) -- (-0.05,\y) node[left] {\l};
        }
        \draw[corErro, dashed] (0,2.97) -- (9.0,2.97);
        \node[corErro, left] at (0,2.97) {\num{28.6}};
        \draw[corFrame, dashed] (0,0.08) -- (9.0,0.08);
        \node[corFrame, left] at (-0.35,0.08) {\num{18.3}};
        \fill[corErro] (2.01,2.97) circle (0.07) node[above right] {$\approx$2006};
        \fill[corFrame] (4.8,0.08) circle (0.07) node[above right] {$\approx$2016};
        \fill[corErro] (7.59,2.97) circle (0.07) node[above left] {$\approx$2025};
      \end{tikzpicture}
    \end{column}
    \begin{column}{0.44\textwidth}
      \small
      \tr{A órbita da Lua é inclinada \bnum{\num{5.1}\gr} em relação ao plano da órbita da
      Terra. Essa inclinação ``aponta'' para um lado e vai girando: uma volta a cada \bnum{\num{18.6}
      anos}.}{The Moon's orbit is tilted \bnum{\num{5.1}\gr} relative to the plane of the
      Earth's orbit. That tilt ``points'' to one side and slowly turns: one full turn every
      \bnum{\num{18.6} years}.}
      \medskip

      \tr{Por isso a \emph{declinação máxima} da Lua (o quanto ela se afasta do equador
      celeste) oscila entre \num{18.3}\gr\ e \num{28.6}\gr.}{That is why the Moon's
      \emph{maximum declination} (how far it strays from the celestial equator) swings between
      \num{18.3}\gr\ and \num{28.6}\gr.}
      \medskip
      \begin{ideia}[title={\tr{Você já viu esse ciclo}{You have seen this cycle}}]
        \footnotesize\tr{É o mesmo período da \bnum{nutação} principal, na série sobre
        referenciais celestes.}{It is the same period as the main \bnum{nutation}, in the
        series on celestial reference systems.}
      \end{ideia}
    \end{column}
  \end{columns}
\end{frame}

\newcommand{\barrasNMM}[1]{%
  \begin{tikzpicture}[x=0.9cm,y=0.9cm,font=\scriptsize]
    \foreach \i/\v/\c in {0/\eNmes/corErro, 1/\eNano/corTempo, 2/\eNdois/corTempo,
                          3/\eNnove/corFrame, 4/\eNdezenove/corMar} {
      \pgfmathsetmacro{\w}{0.5*\v}
      \fill[\c!70] (0,{-0.75*\i}) rectangle (\w,{-0.75*\i-0.5});
      \node[right] at (\w,{-0.75*\i-0.25}) {\num{\v}\,cm};
    }
    \node[left] at (0,-0.25) {\tr{1 mês}{1 month}};
    \node[left] at (0,-1.0) {\tr{1 ano}{1 year}\ifnum#1=1 \ (Torres)\fi};
    \node[left] at (0,-1.75) {\tr{2 anos}{2 years}\ifnum#1=1 \ (Santana)\fi};
    \node[left] at (0,-2.5) {\tr{9 anos}{9 years}\ifnum#1=1 \ (Imbituba)\fi};
    \node[left] at (0,-3.25) {\tr{19 anos}{19 years}\ifnum#1=1 \ \tr{(padrão)}{(standard)}\fi};
    \draw[thick] (0,0.05) -- (0,-3.55);
  \end{tikzpicture}}

\begin{frame}{\tr{Quanto tempo de média?}{How long to average?}}
  \begin{columns}[T]
    \begin{column}{0.56\textwidth}
      \centering
      \barrasNMM{0}\\[3pt]
      {\scriptsize\tr{Erro típico do nível médio estimado, por tempo de observação. Simulação
      \emph{ilustrativa} (série sintética, não medida).}{Typical error of the estimated mean
      level, by observation time. \emph{Illustrative} simulation (synthetic series, not
      measured).}}
    \end{column}
    \begin{column}{0.40\textwidth}
      \footnotesize
      \begin{itemize}
        \setlength{\itemsep}{2pt}
        \item \tr{\textbf{1 mês}: erro de vários cm.}{\textbf{1 month}: error of several cm.}
        \item \tr{\textbf{1 ano}: ainda alguns cm.}{\textbf{1 year}: still a few cm.}
        \item \tr{\textbf{9 anos} (quase um ciclo nodal): cerca de 1\,cm.}{\textbf{9 years}
              (almost a nodal cycle): about 1\,cm.}
        \item \tr{\textbf{19 anos}: poucos mm.}{\textbf{19 years}: a few mm.}
      \end{itemize}
      \smallskip
      \begin{ideia}[title={\tr{Prática}{Practice}}]
        \scriptsize\tr{Os serviços de maré usam séries de \bnum{19 anos} (p.\,ex., a época de
        datum da NOAA): o ciclo nodal em anos inteiros.}{Tide services use \bnum{19-year} series
        (e.g., NOAA's datum epoch): the nodal cycle in whole years.}
      \end{ideia}
    \end{column}
  \end{columns}
\end{frame}

\begin{frame}{\tr{Sua vez}{Your turn}}
  \begin{columns}[T]
    \begin{column}{0.50\textwidth}
      \begin{regra}[title={\tr{O ritmo da maré}{The rhythm of the tide}}]
        \small
        \begin{enumerate}
          \item[(a)] \tr{Hoje a maré alta foi às 06:10. A próxima deve ser a que horas? E a
                primeira de amanhã?}{Today the high tide was at 06:10. About what time should
                the next one be? And tomorrow's first?}
          \item[(b)] \tr{Em \num{18.6} anos, mais ou menos quantas marés altas ocorrem? (Use um
                intervalo de \num{12.42}\,h entre elas.)}{In \num{18.6} years, roughly how many
                high tides occur? (Use an interval of \num{12.42}\,h between them.)}
        \end{enumerate}
      \end{regra}
    \end{column}
    \begin{column}{0.46\textwidth}
      \onslide<2->{
      \begin{ideia}[title={\tr{Respostas}{Answers}}]
        \small\raggedright
        (a) $\num{12.42}\,\mathrm{h}=12\,\mathrm{h}\,25\,\mathrm{min}$: \tr{a próxima alta é às}{the
        next high is at} \bnum{18:35}. \tr{Uma alta por dia atrasa $\approx$ 50\,min:
        amanhã a primeira é perto de}{Each day's high comes about 50\,min later: tomorrow's
        first is near} \bnum{07:00}.
        \medskip

        (b) $\dfrac{\num{18.6}\times\num{365.25}\times24}{\num{12.42}}\approx\bnum{\num{13100}}$.
        \smallskip

        \tr{São \emph{muitas} marés: cada uma é um dado do nível.}{That is a \emph{lot} of
        tides: each one is a data point of the level.}
      \end{ideia}
      }
    \end{column}
  \end{columns}
\end{frame}

% ==================================================================
\section{\tr{Os zeros do Brasil}{Brazil's zeros}}

% faixa do tempo: janelas de observação do NMM de cada datum (1 ano = 0,185 cm)
\newcommand{\faixaZeros}[1]{%
  \begin{tikzpicture}[x=0.185cm,y=0.5cm,font=\scriptsize]
    \draw[thick, ->] (-0.5,0) -- (71,0);
    \foreach \a in {1920,1930,1940,1950,1960,1970,1980} {
      \draw ({\a-1915},0.1) -- ({\a-1915},-0.1) node[below] {\a};
    }
    % Torres: fev/1919 a fev/1920
    \fill[corErro!75] ({1919.1-1915},0.3) rectangle ({1920.1-1915},0.75);
    \node[corErro, right] at ({1920.3-1915},0.52) {Torres (1919--20)};
    \ifnum#1>1
      \draw[corFrame, thick] ({1945.8-1915},0) -- ({1945.8-1915},2.2);
      \node[corFrame, above, font=\tiny] at ({1945.8-1915},2.2) {\tr{1945: início do nivelamento}{1945: leveling starts}};
      \draw[corFrame, thick, dashed] ({1946.9-1915},0) -- ({1946.9-1915},0.9);
      \node[corFrame, right, font=\tiny] at ({1947.1-1915},0.45) {\tr{dez/1946: ligação a Torres}{Dec 1946: tie to Torres}};
    \fi
    \ifnum#1>2
      % Imbituba: 1949 a 1957
      \fill[corFrame!75] ({1949-1915},0.95) rectangle ({1958-1915},1.4);
      \node[corFrame, right] at ({1960-1915},1.18) {Imbituba (1949--57)};
      % um ciclo nodal, para comparar
      \draw[decorate, decoration={brace, amplitude=3pt}] ({1915.2-1915},1.95) -- ({1933.8-1915},1.95);
      \node[above, font=\tiny] at ({1924.5-1915},2.1) {\num{18.6}\,\tr{anos}{years}: \tr{um ciclo nodal}{a nodal cycle}};
    \fi
    \ifnum#1>3
      \fill[corTempo!30] ({1958-1915},-0.1) rectangle ({1959.5-1915},2.2);
      \node[corTempo, above right, font=\tiny] at ({1958.0-1915},2.2) {\tr{1958 / 1959: datas nos relatórios}{1958 / 1959: report dates}};
    \fi
    \ifnum#1>4
      \fill[corSist!75] ({1957-1915},1.6) rectangle ({1959-1915},2.05);
      \node[corSist, right] at ({1960-1915},1.83) {Santana (1957--58)};
      \draw[corSist, thick, dashed] ({1980-1915},0) -- ({1980-1915},2.2);
      \node[corSist, above, font=\tiny] at ({1980-1915},2.2) {\tr{1980: adoção de Santana}{1980: Santana adopted}};
    \fi
  \end{tikzpicture}}

\begin{frame}{\tr{O primeiro zero: Torres (1919--1920)}{The first zero: Torres (1919--1920)}}
  \begin{columns}[T]
    \begin{column}{0.54\textwidth}
      \small
      \tr{Antes de existir uma rede de nivelamento, já havia um zero \bnum{provisório}.}{Before
      there was a leveling network, there already was a \bnum{provisional} zero.}
      \medskip

      \tr{A Comissão da Carta Geral do Brasil instalou um marégrafo em \textbf{Torres} (RS). O
      nível médio de \textbf{fevereiro de 1919 a fevereiro de 1920} --- cerca de \bnum{um
      ano} --- virou o zero.}{The Comissão da Carta Geral do Brasil installed a tide gauge at
      \textbf{Torres} (RS). The mean level from \textbf{February 1919 to February 1920} ---
      about \bnum{one year} --- became the zero.}
      \medskip

      \tr{Era provisório de propósito. Comparações com os marégrafos de Laguna e Florianópolis
      levaram a suspeitar de um erro de \bnum{um decímetro ou mais}.}{It was provisional on
      purpose. Comparisons with the Laguna and Florianópolis gauges led to the suspicion of an
      error of \bnum{a decimeter or more}.}
    \end{column}
    \begin{column}{0.42\textwidth}
      \begin{cuidado}[title={\tr{Lembra das barras?}{Remember the bars?}}]
        \small\tr{Com \emph{um ano} de dados, a simulação da aula deu um erro típico de alguns
        centímetros. A suspeita histórica (\num{10}\,cm ou mais) é da \emph{mesma ordem}.}{With
        \emph{one year} of data, the simulation in this class gave a typical error of a few
        centimeters. The historical suspicion (\num{10}\,cm or more) is of the \emph{same
        order}.}
      \end{cuidado}
    \end{column}
  \end{columns}
  \vfill
  \centering
  \faixaZeros{1}
\end{frame}

\begin{frame}{\tr{1945: começa o nivelamento de precisão}{1945: precision leveling begins}}
  \begin{columns}[T]
    \begin{column}{0.54\textwidth}
      \footnotesize
      \begin{itemize}
        \setlength{\itemsep}{4pt}
        \item \tr{Em \bnum{13 de outubro de 1945}, a Seção de Nivelamento do \textbf{Conselho
              Nacional de Geografia} (CNG, do IBGE) começa o nivelamento de precisão (data
              segundo H.\ Beserra, na \emph{RBG}).}{On \bnum{13 October 1945}, the Leveling
              Section of the \textbf{Conselho Nacional de Geografia} (CNG, of IBGE) starts
              precision leveling (date per H.\ Beserra, in the \emph{RBG}).}
        \item \tr{A rede nasce em \textbf{Urussanga} (SC) e cresce \emph{ao longo das
              estradas}.}{The network is born at \textbf{Urussanga} (SC) and grows \emph{along
              the roads}.}
        \item \tr{Só em \bnum{dezembro de 1946} a rede é ligada ao marégrafo de Torres: nos
              primeiros 14 meses, as altitudes ainda não tinham \emph{zero}.}{Only in
              \bnum{December 1946} is the network tied to the Torres gauge: for the first 14
              months the heights still had no \emph{zero}.}
      \end{itemize}
    \end{column}
    \begin{column}{0.42\textwidth}
      \begin{ideia}[title={\tr{Uma pergunta}{A question}}]
        \small\tr{Torres fica no litoral do Rio Grande do Sul. Como o zero chega a
        Urussanga, a \emph{centenas de quilômetros}, sem nenhum mar no caminho?}{Torres is on
        the coast of Rio Grande do Sul. How does the zero reach Urussanga, \emph{hundreds of
        kilometers} away, with no sea on the way?}
        \medskip

        \tr{\bnum{Por nivelamento}: é a Aula 2.}{\bnum{By leveling}: that is Class 2.}
      \end{ideia}
    \end{column}
  \end{columns}
  \vfill
  \centering
  \faixaZeros{2}
\end{frame}

\begin{frame}{\tr{Imbituba: um novo marégrafo (1949--1957)}{Imbituba: a new tide gauge (1949--1957)}}
  \begin{columns}[T]
    \begin{column}{0.54\textwidth}
      \small
      \begin{itemize}
        \setlength{\itemsep}{4pt}
        \item \tr{No \textbf{porto de Imbituba} (SC), há registros de marégrafo desde \num{1948}
              (as séries do PSMSL vão de 1948 a 1968).}{At the \textbf{port of Imbituba}
              (SC), there are tide gauge records since \num{1948} (the PSMSL series run from
              1948 to 1968).}
        \item \tr{O \bnum{datum} foi definido com as observações de \bnum{1949 a 1957}:
              \bnum{9 anos}.}{The \bnum{datum} was defined with the observations of \bnum{1949
              to 1957}: \bnum{9 years}.}
        \item \tr{No REALT-2018, a \textbf{RN 4X} é a referência do datum de Imbituba.}{In
              REALT-2018, \textbf{RN 4X} is the reference of the Imbituba datum.}
      \end{itemize}
    \end{column}
    \begin{column}{0.42\textwidth}
      \begin{ideia}[title={\tr{Nove anos é quase metade do ciclo}{Nine years is almost half the
        cycle}}]
        \small\tr{$9/\num{18.6}\approx\num{0.48}$: a janela cobre quase \emph{metade} do ciclo
        nodal. É bem melhor que o ano de Torres --- mas não são os 19 anos do padrão de
        hoje.}{$9/\num{18.6}\approx\num{0.48}$: the window covers almost \emph{half} of the
        nodal cycle. It is much better than Torres' one year --- but not today's 19-year
        standard.}
      \end{ideia}
    \end{column}
  \end{columns}
  \vfill
  \centering
  \faixaZeros{3}
\end{frame}

\begin{frame}{\tr{Torres e Imbituba: duas datas documentadas}{Torres and Imbituba: two documented dates}}
  \begin{columns}[T]
    \begin{column}{0.54\textwidth}
      \footnotesize
      \tr{O relatório do \bnum{IBGE de 2011}, p. 8, situa a substituição de Torres por
      Imbituba em \bnum{1958}, com cerca de \num{30000}\,km de linhas.}{
      The \bnum{2011 IBGE report}, p. 8, dates the replacement of Torres by Imbituba
      to \bnum{1958}, with about \num{30000}\,km of lines.}
      \par\medskip
      \tr{O relatório do \bnum{REALT, 2.ª ed. (2019)}, p. 13, descreve a adoção
      de Imbituba em \bnum{1959}, quando terminou o terceiro ajustamento.
      Os ajustes de 1948 e 1952 usaram Torres; a partir de 1959, Imbituba.}{
      The \bnum{REALT report, 2nd ed. (2019)}, p. 13, describes Imbituba's adoption
      in \bnum{1959}, when the third adjustment ended. The 1948 and 1952 adjustments
      used Torres; those from 1959 onward used Imbituba.}
    \end{column}
    \begin{column}{0.42\textwidth}
      \begin{cuidado}[title={\tr{Uma ressalva histórica}{A historical qualification}}]
        \footnotesize
        \tr{As duas publicações oficiais dão anos distintos. A cronologia registra
        ambos, sem supor que descrevam exatamente o mesmo ato administrativo.}{
        The two official publications give different years. The timeline records both
        without assuming they describe exactly the same administrative act.}
      \end{cuidado}
    \end{column}
  \end{columns}
  \vfill
  \centering
  \faixaZeros{4}
\end{frame}

\begin{frame}{\tr{O Amazonas não se atravessa}{The Amazon cannot be crossed}}
  \begin{columns}[T]
    \begin{column}{0.52\textwidth}
      \centering
      \begin{tikzpicture}[x=0.9cm,y=0.9cm,font=\scriptsize]
        % rio / estuário
        \fill[corMar!22] (0,2.0) -- (7.6,1.6) -- (7.6,3.3) -- (0,3.8) -- cycle;
        \node[corMar, font=\tiny, anchor=west, text width=2.6cm] at (0.15,2.85) {\tr{foz do Amazonas e ilha de Marajó}{Amazon mouth and Marajó island}};
        % redes
        \foreach \x/\y in {0.6/4.6, 1.5/5.0, 2.4/4.5, 1.1/4.1, 2.0/4.0} {
          \fill[corSist] (\x,\y) circle (0.07);
        }
        \draw[rede, corSist] (0.6,4.6) -- (1.5,5.0) -- (2.4,4.5) -- (2.0,4.0) -- (1.1,4.1) -- cycle;
        \foreach \x/\y in {5.0/0.6, 5.9/1.0, 6.8/0.5, 5.4/0.1, 6.4/0.0} {
          \fill[corFrame] (\x,\y) circle (0.07);
        }
        \draw[rede] (5.0,0.6) -- (5.9,1.0) -- (6.8,0.5) -- (6.4,0.0) -- (5.4,0.1) -- cycle;
        \node[corSist, above] at (1.5,5.1) {Amapá (Santana)};
        \node[corFrame, below] at (5.9,-0.1) {\tr{resto do país}{rest of the country}};
        \draw[corErro, thick, <->, dashed] (2.4,4.4) -- node[right=3pt, text=corErro, align=left]
              {$\approx$\num{330}\,km} (5.0,0.8);
      \end{tikzpicture}
    \end{column}
    \begin{column}{0.44\textwidth}
      \small
      \tr{O nivelamento geométrico de alta precisão anda \emph{em terra}, com nível e mira. Ele
      \bnum{não atravessa} a foz do Amazonas e a ilha de Marajó.}{High-precision geometric
      leveling walks \emph{on land}, with level and rod. It \bnum{cannot cross} the Amazon
      mouth and Marajó island.}
      \medskip

      \tr{Por isso a pequena rede do \bnum{Amapá} nunca foi ligada à do resto do país. Ela
      ganhou um \bnum{zero próprio}: o nível médio do mar no \textbf{porto de Santana}, de
      \bnum{1957--58}, adotado em \bnum{1980}. As RRNN antigas foram implantadas
      em \bnum{1980--81}.}{That is why the small \bnum{Amapá} network was never tied to
      that of the rest of the country. It got its \bnum{own zero}: mean sea level at the
      \textbf{port of Santana} in \bnum{1957--58}, adopted in \bnum{1980}.
      The old benchmarks were established in \bnum{1980--81}.}
      \fonteslide{IBGE (2019), \tr{pp. 13 e 44}{pp. 13 and 44}.}
      \medskip

      {\scriptsize\tr{As RRNN mais próximas das duas redes ficam a cerca de \num{330}\,km.}{The
      nearest benchmarks of the two networks are about \num{330}\,km apart.}}
    \end{column}
  \end{columns}
\end{frame}

\begin{frame}{\tr{Os dois zeros do Brasil}{Brazil's two zeros}}
  \begin{columns}[T]
    \begin{column}{0.27\textwidth}
      \centering
      \begin{tikzpicture}[x=0.70cm,y=0.70cm,font=\scriptsize]
        \costaBR{corMar, thick}
        \coordinate (Sa) at (\ptSantana);
        \coordinate (Be) at (\ptBelem);
        \coordinate (Im) at (\ptImbituba);
        \coordinate (To) at (\ptTorres);
        \fill[corSist] (Sa) circle (0.09);
        \fill[corFrame] (Im) circle (0.09);
        \fill[corErro] (To) circle (0.09);
        \node[corSist, left=0.05cm] at (Sa) {Santana};
        \node[corFrame, right=0.25cm] at (Im) {Imbituba};
        \node[corErro, below right=0.1cm] at (To) {Torres};
        \node[mStoneMid, right=0.05cm, font=\tiny] at (Be) {Belém};
        \fill[mStoneMid] (Be) circle (0.04);
      \end{tikzpicture}\\[1pt]
      {\tiny\tr{esquema: costa ligando cidades}{sketch: coast linking cities}}
    \end{column}
    \begin{column}{0.70\textwidth}
      \small
      \renewcommand{\arraystretch}{1.5}
      \begin{tabular}{@{}L{2.5cm}L{3.1cm}L{2.9cm}@{}}
        \toprule
         & \textcolor{corFrame}{\textbf{Imbituba} (SC)} & \textcolor{corSist}{\textbf{Santana}
        (AP)} \\
        \midrule
        \tr{NMM observado}{MSL observed} & \bnum{1949--1957} \tr{(9 anos)}{(9 years)} &
        \bnum{1957--1958} \tr{(2 anos)}{(2 years)} \\
        \tr{RN do datum}{Datum RN} & \textbf{4X} & \textbf{9329T} \\
        \tr{Onde vale}{Where it applies} & \tr{quase todo o país}{almost the whole country} &
        \tr{o Amapá}{Amapá} \\
        \tr{Adotado}{Adopted} & 1958 / 1959\textsuperscript{a} & 1980 \\
        \bottomrule
      \end{tabular}
      \bigskip

      \footnotesize\tr{\textsuperscript{a}IBGE (2011), p. 8: 1958; IBGE (2019), p. 13: 1959.
      As altitudes apoiam-se em dois valores de NMM; a relação entre eles volta na Aula 4.}{
      \textsuperscript{a}IBGE (2011), p. 8: 1958; IBGE (2019), p. 13: 1959.
      The heights rest on two MSL values; their relation returns in Class 4.}
    \end{column}
  \end{columns}
\end{frame}

\begin{frame}{\tr{Sua vez}{Your turn}}
  \begin{columns}[T]
    \begin{column}{0.50\textwidth}
      \begin{regra}[title={\tr{Três janelas, três zeros}{Three windows, three zeros}}]
        \small
        \begin{enumerate}
          \item[(a)] \tr{Que fração do ciclo nodal (\num{18.6} anos) cada janela cobre:
                Torres (1 ano), Santana (2 anos), Imbituba (9 anos)?}{What fraction of the
                nodal cycle (\num{18.6} years) does each window cover: Torres (1 year), Santana
                (2 years), Imbituba (9 years)?}
          \item[(b)] \tr{Na simulação, o erro do NMM é de \num{3.1}\,cm com 1 ano e de
                \num{1.4}\,cm com 9 anos. Quantas vezes o erro de Torres seria maior que o de
                Imbituba?}{In the simulation the MSL error is \num{3.1}\,cm with 1 year and
                \num{1.4}\,cm with 9 years. How many times larger would Torres' error be than
                Imbituba's?}
        \end{enumerate}
      \end{regra}
    \end{column}
    \begin{column}{0.46\textwidth}
      \onslide<2->{
      \begin{ideia}[title={\tr{Respostas}{Answers}}]
        \small\raggedright
        (a) $1/\num{18.6}\approx\bnum{\num{0.05}}$; \ $2/\num{18.6}\approx\bnum{\num{0.11}}$; \
        $9/\num{18.6}\approx\bnum{\num{0.48}}$.
        \medskip

        (b) $\num{3.1}/\num{1.4}\approx\bnum{\num{2.2}}$ \tr{vezes.}{times.}
        \medskip

        \tr{\alert{Cuidado}: as barras são de uma simulação. O que importa é a \emph{ideia}:
        quanto mais ciclo coberto, mais estável é o zero.}{\alert{Careful}: the bars come from a
        simulation. What matters is the \emph{idea}: the more of the cycle covered, the more
        stable the zero.}
      \end{ideia}
      }
    \end{column}
  \end{columns}
\end{frame}

% ==================================================================
\section{\tr{O zero não é o geoide}{The zero is not the geoid}}

\begin{frame}{\tr{O nível médio do mar não é uma superfície de nível}{Mean sea level is not a
level surface}}
  \begin{columns}[T]
    \begin{column}{0.55\textwidth}
      \centering
      \begin{tikzpicture}[x=0.82cm,y=0.95cm,font=\scriptsize]
        \draw[corSist, very thick] (0,1.5) -- (9,1.5);
        \node[corSist, below left] at (9,1.5) {\tr{geoide (nível)}{geoid (level)}};
        \draw[corMar, very thick] plot[smooth] coordinates {(0,1.9) (1,2.35) (2,1.8) (3,1.1)
              (4,0.95) (5,1.5) (6,2.3) (7,2.7) (8,2.4) (9,1.9)};
        \node[corMar, above] at (4.0,0.95) {};
        \node[corMar, below] at (4.0,0.92) {\tr{nível médio do mar}{mean sea level}};
        \draw[corFrame, dashed] (1.5,0.3) -- (1.5,2.2);
        \draw[corErro, thick, <->] (1.5,1.5) -- (1.5,2.1);
        \node[corFrame, below] at (1.5,0.3) {\tr{estação $A$}{station $A$}};
        \node[corErro, left] at (1.45,1.85) {$\Delta_A$};
        \draw[corFrame, dashed] (7.0,0.3) -- (7.0,2.9);
        \draw[corErro, thick, <->] (7.0,1.5) -- (7.0,2.7);
        \node[corFrame, below] at (7.0,0.3) {\tr{estação $B$}{station $B$}};
        \node[corErro, left] at (6.95,2.1) {$\Delta_B$};
      \end{tikzpicture}\\[1pt]
      {\scriptsize\tr{Esquema exagerado, fora de escala.}{Exaggerated sketch, not to scale.}}
    \end{column}
    \begin{column}{0.41\textwidth}
      \footnotesize
      \tr{O \textbf{geoide} é uma superfície de nível: a água \emph{parada} ficaria sobre
      ela.}{The \textbf{geoid} is a level surface: \emph{still} water would sit on it.}
      \medskip

      \tr{O mar de verdade não fica parado. \bnum{Correntes, ventos, pressão, temperatura e
      salinidade} o empurram para cima ou para baixo, por décimos de metro.}{The real sea does
      not stand still. \bnum{Currents, winds, pressure, temperature and salinity} push it up or
      down by tenths of a meter.}
      \medskip

      \begin{cuidado}[title={\tr{Consequência}{Consequence}}]
        \footnotesize\tr{O NMM de cada estação fica um pouco \emph{acima} ou \emph{abaixo} do geoide,
        e a diferença ($\Delta_A$, $\Delta_B$) \bnum{muda de lugar para lugar}.}{The MSL of each
        station lies a little \emph{above} or \emph{below} the geoid, and the difference
        ($\Delta_A$, $\Delta_B$) \bnum{changes from place to place}.}
      \end{cuidado}
    \end{column}
  \end{columns}
\end{frame}

\begin{frame}{\tr{Datum vertical: um valor, numa estação, numa época}{Vertical datum: one value,
at one station, at one epoch}}
  \centering
  \begin{tikzpicture}[x=1cm,y=1cm,font=\small]
    \node[caixa, draw=corFrame, fill=corFrame!8, text width=3.6cm, minimum height=2.2cm] (a) at (0,0)
      {\textbf{\tr{ONDE}{WHERE}}\\[3pt]\scriptsize\tr{uma estação: o marégrafo de
       Imbituba (ou o de Santana)}{one station: the Imbituba tide gauge (or Santana's)}};
    \node[caixa, draw=corTempo, fill=corTempo!8, text width=3.6cm, minimum height=2.2cm] (b) at (4.7,0)
      {\textbf{\tr{QUANDO}{WHEN}}\\[3pt]\scriptsize\tr{uma janela de anos: 1949--1957
       (Santana: 1957--1958)}{a window of years: 1949--1957 (Santana: 1957--1958)}};
    \node[caixa, draw=corSist, fill=corSist!8, text width=3.6cm, minimum height=2.2cm] (c) at (9.4,0)
      {\textbf{\tr{COMO}{HOW}}\\[3pt]\scriptsize\tr{a média do nível do mar nessa janela
       \emph{vale altitude zero}, por convenção}{the mean sea level in that window
       \emph{equals zero altitude}, by convention}};
  \end{tikzpicture}
  \bigskip

  \begin{regra}[title={\tr{Definição}{Definition}}]
    \small\tr{Um \bnum{datum vertical} é \emph{um valor do nível médio do mar}, medido
    \emph{numa estação} durante \emph{uma janela de tempo}, a que se atribui altitude zero. O
    resto do país não tem marégrafo: o zero chega até ele por \bnum{nivelamento} --- Aula
    2.}{A \bnum{vertical datum} is \emph{one value of mean sea level}, measured \emph{at one
    station} during \emph{one time window}, to which zero altitude is assigned. The rest of the
    country has no tide gauge: the zero reaches it by \bnum{leveling} --- Class 2.}
  \end{regra}
\end{frame}

\begin{frame}{\tr{Outros zeros: nível de redução e zero hidrográfico}{Other zeros: reduction level
and hydrographic zero}}
  \begin{columns}[T]
    \begin{column}{0.50\textwidth}
      \centering
      \begin{tikzpicture}[x=0.56cm,y=0.95cm,font=\scriptsize]
        \draw[->] (-0.1,0) -- (7.7,0);
        \draw[->] (0,-1.9) -- (0,1.9);
        \mareDia{corMar, thick}
        % nível mais alto e mais baixo da curva
        \draw[mStoneMid, dashed] (0,1.27) -- (7.4,1.27);
        \node[mStoneMid, right] at (7.4,1.27) {\tr{preamar}{high water}};
        \draw[corMar, dashed, thick] (0,0) -- (7.4,0);
        \node[corMar, right] at (7.4,0) {\tr{NMM local}{local MSL}};
        \draw[corSist, very thick] (0,-0.18) -- (7.4,-0.18);
        \node[corSist, right] at (7.4,-0.55) {\tr{datum do SGB}{SGB datum}};
        \draw[corTempo, very thick] (0,-1.3) -- (7.4,-1.3);
        \node[corTempo, right, align=left] at (7.4,-1.3) {\tr{zero hidrográfico}{hydrographic zero}};
      \end{tikzpicture}
    \end{column}
    \begin{column}{0.44\textwidth}
      \footnotesize
      \tr{\textbf{Carta náutica} (Marinha): a profundidade é contada a partir do \bnum{nível de
      redução} (zero hidrográfico), um plano perto do nível \emph{mais baixo} da maré. Assim o
      navegante nunca encontra menos água do que a carta diz.}{\textbf{Nautical chart} (Navy):
      depth is counted from the \bnum{reduction level} (hydrographic zero), a plane near the
      \emph{lowest} tide. This way the navigator never finds less water than the chart
      says.}
      \medskip

      \tr{\textbf{Altitude} (IBGE): contada a partir do \bnum{datum} do SGB, um nível
      \emph{médio}.}{\textbf{Altitude} (IBGE): counted from the SGB \bnum{datum}, a
      \emph{mean} level.}
      \medskip

      \begin{cuidado}[title={\tr{São zeros diferentes}{They are different zeros}}]
        \scriptsize\tr{Em cada porto, a relação entre eles é diferente, e o IBGE a estuda e
        publica (rede RMPG, Aula 4). A IHO recomenda, nas cartas, o nível mais baixo da maré
        astronômica, calculado com $\geq\num{18.6}$ anos.}{In each port the relation between
        them is different, and IBGE studies and publishes it (RMPG network, Class 4). The IHO
        recommends, for charts, the lowest astronomical tide, computed with $\geq\num{18.6}$
        years.}
      \end{cuidado}
    \end{column}
  \end{columns}
\end{frame}

\begin{frame}{\tr{O mar sobe --- e a altitude não muda}{The sea rises --- and the altitude does
not change}}
  \begin{columns}[T]
    \begin{column}{0.54\textwidth}
      \centering
      \begin{tikzpicture}[x=0.078cm,y=0.17cm,font=\scriptsize]
        % eixo: x = ano - 1945; y = cm
        \draw[->] (0,0) -- (86,0);
        \draw[->] (0,0) -- (0,21) node[above] {cm};
        \foreach \a in {1950,1970,1990,2010,2030} {
          \draw ({\a-1945},0.4) -- ({\a-1945},-0.4) node[below] {\a};
        }
        \foreach \y in {5,10,15,20} {
          \draw (0.6,\y) -- (-0.6,\y) node[left] {\y};
        }
        % janela do datum: 1949-1957
        \draw[corFrame, line width=3pt] ({1949-1945},0) -- ({1958-1945},0);
        % faixa de incerteza e linha de tendência
        \fill[corMar!25] ({1953.5-1945},0) -- ({2026-1945},13.3) -- ({2026-1945},19.1) -- cycle;
        \draw[corMar, very thick] ({1953.5-1945},0) -- ({2026-1945},16.2);
        \draw[corSist, dashed, thick] (0,0) -- (84,0);
        \node[corFrame, above] at ({1953.5-1945},0.8) {\tr{datum}{datum}};
        \node[corMar, left] at ({2026-1945},18.2) {\tr{$+\num{16}$ cm}{$+\num{16}$ cm}};
      \end{tikzpicture}\\[1pt]
      {\scriptsize\tr{NMM em relação ao zero do datum de Imbituba, com a taxa
      $\num{2.24}\pm\num{0.4}$\,mm/ano (estudo da UFPR, 2017).}{MSL relative to the Imbituba datum
      zero, with the rate $\num{2.24}\pm\num{0.4}$\,mm/year (UFPR study, 2017).}}
    \end{column}
    \begin{column}{0.42\textwidth}
      \small
      \tr{O datum usa a janela de 1949--57. Se a taxa do estudo fosse constante,
      a projeção até 2026 seria de cerca de \bnum{\num{16}\,cm}.
      É uma extrapolação ilustrativa.}{The datum uses the 1949--57 window. If the study's
      rate were constant, the projection to 2026 would be about \bnum{\num{16}\,cm}.
      This is an illustrative extrapolation.}
      \medskip

      \tr{A \textbf{altitude de uma RN} não muda por isso: ela é ``tantos metros acima do
      nível médio de \emph{1949--57}'', e continua sendo.}{The \textbf{altitude of a benchmark}
      does not change because of it: it is ``so many meters above the mean level of
      \emph{1949--57}'', and it still is.}
      \medskip

      {\scriptsize\tr{Mesmo estudo: o GNSS do marégrafo indicou subsidência de cerca de
      \num{3}\,mm/ano (campanhas curtas). Valores de um estudo, não oficiais.}{Same study: the
      gauge's GNSS indicated subsidence of about \num{3}\,mm/year (short campaigns). Values from
      one study, not official.}}
    \end{column}
  \end{columns}
\end{frame}

\begin{frame}{\tr{Sua vez}{Your turn}}
  \begin{columns}[T]
    \begin{column}{0.50\textwidth}
      \begin{regra}[title={\tr{O mar de hoje}{Today's sea}}]
        \small
        \begin{enumerate}
          \item[(a)] \tr{O NMM sobe \num{2.24}\,mm/ano. De 1953,5 (meio da janela) até 2026
                passam \num{72.5} anos. Quantos centímetros subiu?}{MSL rises \num{2.24}\,mm/year.
                From 1953.5 (middle of the window) to 2026, \num{72.5} years pass. How many
                centimeters has it risen?}
          \item[(b)] \tr{Uma RN no litoral tem altitude \num{7.000}\,m no datum. A que altura
                ela está acima do mar \emph{de hoje}? (Ignore maré, ondas e o movimento do
                chão.)}{A benchmark on the coast has altitude \num{7.000}\,m in the datum. How
                high is it above \emph{today's} sea? (Ignore tides, waves and ground
                motion.)}
        \end{enumerate}
      \end{regra}
    \end{column}
    \begin{column}{0.46\textwidth}
      \onslide<2->{
      \begin{ideia}[title={\tr{Respostas}{Answers}}]
        \small\raggedright
        (a) $\num{2.24}\times\num{72.5}=\num{162}$\,mm $\approx\bnum{\num{16}\,cm}$ ($\pm\num{3}$\,cm
        \tr{pela incerteza da taxa}{from the rate's uncertainty}).
        \medskip

        (b) $\num{7.000}-\num{0.162}\approx\bnum{\num{6.84}\,m}$.
        \medskip

        \tr{A altitude oficial é \emph{a mesma}; só a distância ao mar de hoje mudou.}{The
        official altitude is \emph{the same}; only the distance to today's sea changed.}
      \end{ideia}
      }
    \end{column}
  \end{columns}
\end{frame}

\begin{frame}{\tr{Pare e pense}{Stop and think}}
  \begin{columns}[T]
    \begin{column}{0.46\textwidth}
      \begin{armadilha}[title={\tr{Uma ideia tentadora}{A tempting idea}}]
        \small\tr{``Se o mar sobe, é só \emph{atualizar o zero} todo ano. Assim o zero
        acompanha o mar.''}{``If the sea rises, just \emph{update the zero} every year. The zero
        then follows the sea.''}
        \medskip

        \tr{Por que o Brasil \emph{não} faz isso?}{Why does Brazil \emph{not} do that?}
      \end{armadilha}
    \end{column}
    \begin{column}{0.50\textwidth}
      \onslide<2->{
      \begin{ideia}[title={\tr{Porque}{Because}}]
        \footnotesize
        \begin{itemize}
          \setlength{\itemsep}{2pt}
          \item \tr{\bnum{todas} as altitudes (mapas, projetos, leis) dependem do zero;}{\bnum{all}
                altitudes (maps, projects, laws) depend on the zero;}
          \item \tr{o datum é uma \emph{convenção} fixada numa época;}{the datum is a
                \emph{convention} fixed at an epoch;}
          \item \tr{o que se faz é \bnum{medir a diferença} entre o datum e o mar de hoje (rede
                RMPG, Aula 4) e, no futuro, ligar tudo a um zero \emph{global}.}{what is done is
                to \bnum{measure the difference} between the datum and today's sea (RMPG
                network, Class 4) and, in the future, tie everything to a \emph{global}
                zero.}
        \end{itemize}
      \end{ideia}
      }
    \end{column}
  \end{columns}
\end{frame}

% ==================================================================
\section{\tr{Fechamento}{Wrap-up}}

\begin{frame}{\tr{As quatro perguntas, respondidas}{The four questions, answered}}
  \small
  \begin{enumerate}
    \setlength{\itemsep}{5pt}
    \item \textbf{\tr{Por que a altitude precisa de um zero?}{Why does altitude need a
          zero?}} \tr{Porque é a leitura de uma régua, e toda régua começa em algum lugar.
          O zero é uma \bnum{convenção}: o nível médio do mar numa estação, numa
          época.}{Because it is the reading of a ruler, and every ruler starts somewhere. The
          zero is a \bnum{convention}: the mean sea level at a station, at an epoch.}
    \item \textbf{\tr{Como se mede o nível do mar --- e por quanto tempo?}{How is sea level
          measured --- and for how long?}} \tr{Com um \bnum{marégrafo}, tirando a média de
          muitas leituras. Um ano deixa alguns centímetros de erro; o padrão é \bnum{19 anos},
          o ciclo nodal da Lua.}{With a \bnum{tide gauge}, averaging many readings. One year
          leaves a few centimeters of error; the standard is \bnum{19 years}, the Moon's nodal
          cycle.}
    \item \textbf{\tr{Qual foi o zero do Brasil, e por que hoje há dois?}{What has been
          Brazil's zero, and why are there two today?}} \tr{Torres (1919--20, provisório),
          depois \bnum{Imbituba} (1949--57, adotado em 1958--59). O Amazonas não se atravessa:
          o Amapá usa \bnum{Santana} (1957--58).}{Torres (1919--20, provisional), then
          \bnum{Imbituba} (1949--57, adopted in 1958--59). The Amazon cannot be crossed: Amapá
          uses \bnum{Santana} (1957--58).}
    \item \textbf{\tr{Esse zero é ``o mar'' de verdade?}{Is that zero really ``the
          sea''?}} \tr{Não: é o mar \emph{daquela janela}. O mar sobe, o NMM local não é uma
          superfície de nível, e o zero das cartas náuticas é outro.}{No: it is the sea \emph{of
          that window}. The sea rises, local MSL is not a level surface, and the nautical charts'
          zero is another one.}
  \end{enumerate}
\end{frame}

\begin{frame}{\tr{Para levar para casa}{Take home}}
  \begin{columns}[T]
    \begin{column}{0.54\textwidth}
      \begin{regra}[title={\tr{Quatro ideias}{Four ideas}}]
        \small
        \begin{enumerate}
          \setlength{\itemsep}{3pt}
          \item \tr{Altitude é a leitura de uma régua; o zero é uma \bnum{escolha}.}{Altitude is
                the reading of a ruler; the zero is a \bnum{choice}.}
          \item \tr{Nível médio do mar = \bnum{média no tempo}; o padrão é 19 anos.}{Mean sea
                level = \bnum{average over time}; the standard is 19 years.}
          \item \tr{Datum = \bnum{um valor, numa estação, numa época}; o Brasil tem dois.}{Datum
                = \bnum{one value, at one station, at one epoch}; Brazil has two.}
          \item \tr{O mar sobe; a altitude da RN \bnum{não muda}.}{The sea rises; the benchmark's
                altitude \bnum{does not change}.}
        \end{enumerate}
      \end{regra}
    \end{column}
    \begin{column}{0.42\textwidth}
      \begin{ideia}[title={\tr{Ponte para a Aula 2}{Bridge to Class 2}}]
        \small\tr{O zero está em Imbituba. Mas a sua obra está a \emph{centenas de
        quilômetros} do mar. \bnum{Como se leva o zero até lá?}}{The zero is at Imbituba. But
        your construction site is \emph{hundreds of kilometers} from the sea. \bnum{How is the
        zero carried there?}}
        \medskip

        \tr{Resposta: com um nível, uma mira e muita paciência --- o \bnum{nivelamento}.}{Answer:
        with a level, a rod and lots of patience --- \bnum{leveling}.}
      \end{ideia}
    \end{column}
  \end{columns}
\end{frame}

% ==================================================================
\section{\tr{Aula 2: Lance a lance}{Class 2: Leg by leg}}

\begin{frame}{\tr{As perguntas de hoje}{Today's questions}}
  \begin{columns}[T]
    \begin{column}{0.58\textwidth}
      \begin{enumerate}
        \setlength{\itemsep}{8pt}
        \item \tr{Como se leva o zero do mar até um ponto a centenas de quilômetros?}{How is the
              zero carried from the sea to a point hundreds of kilometers away?}
        \item \tr{De onde vem o erro do nivelamento --- e por que ele cresce com a raiz da
              distância?}{Where does leveling error come from --- and why does it grow with the
              square root of the distance?}
        \item \tr{Como é a rede de referências de nível que o Brasil construiu?}{What does the
              network of benchmarks that Brazil built look like?}
        \item \tr{Por que um circuito medido \emph{sem erro} não fecha em zero?}{Why does a
              loop measured \emph{without error} not close at zero?}
      \end{enumerate}
    \end{column}
    \begin{column}{0.38\textwidth}
      \begin{ideia}[title={\tr{Combinado}{Deal}}]
        \tr{No fim da aula, você vai saber responder às quatro --- com as suas
        palavras.}{By the end of the class you will be able to answer all four --- in your own
        words.}
        \medskip

        \tr{Elas voltam no último slide.}{They come back on the last slide.}
      \end{ideia}
    \end{column}
  \end{columns}
\end{frame}

\begin{frame}{\tr{O que a aula anterior deixou}{What the previous class left us}}
  \begin{columns}[T]
    \begin{column}{0.52\textwidth}
      \small
      \tr{Na Aula 1, três ideias:}{In Class 1, three ideas:}
      \begin{enumerate}
        \setlength{\itemsep}{3pt}
        \item \tr{altitude é a leitura de uma régua; o \bnum{zero} é uma escolha;}{altitude is
              the reading of a ruler; the \bnum{zero} is a choice;}
        \item \tr{o zero do Brasil é o \bnum{nível médio do mar} num marégrafo, numa janela de
              anos: Imbituba (1949--57) e Santana (1957--58);}{Brazil's zero is the \bnum{mean
              sea level} at a tide gauge, in a window of years: Imbituba (1949--57) and Santana
              (1957--58);}
        \item \tr{o mar sobe, mas a altitude de uma RN \bnum{não muda}.}{the sea rises, but the
              altitude of a benchmark \bnum{does not change}.}
      \end{enumerate}
    \end{column}
    \begin{column}{0.44\textwidth}
      \centering
      \begin{tikzpicture}[x=0.95cm,y=0.95cm,font=\scriptsize]
        \fill[corMar!18] (0,0) rectangle (1.3,2.6);
        \node[corMar, rotate=90] at (0.65,1.3) {\tr{mar}{sea}};
        \draw[thick, mStoneMid] (1.3,0.5) .. controls (2.2,0.7) and (3.2,1.6) .. (4.2,1.5)
              .. controls (5.0,1.4) and (5.4,2.0) .. (6.0,2.2);
        \fill[corFrame] (1.45,0.55) circle (0.08);
        \node[corFrame, above] at (1.6,0.65) {\tr{marégrafo}{tide gauge}};
        \fill[corErro] (5.8,2.15) circle (0.09);
        \node[corErro, above] at (5.7,2.3) {\tr{a sua obra}{your site}};
        \draw[corErro, thick, dashed, -{Stealth[length=2mm]}] (1.55,0.45) .. controls (3,0.3) and (4.5,0.9) .. (5.75,2.0);
        \node[corErro, below right, align=left] at (3.0,0.35) {\tr{\textbf{600 km}: como}{\textbf{600 km}: how}\\\tr{o zero chega lá?}{does the zero get there?}};
      \end{tikzpicture}
    \end{column}
  \end{columns}
\end{frame}

\begin{frame}{\tr{O fio da aula}{The thread of the class}}
  \centering
  \begin{tikzpicture}[x=1cm,y=1cm,font=\small]
    \node[caixa, draw=corFrame, fill=corFrame!8, text width=2.8cm, minimum height=2.1cm] (a) at (0,0)
      {\textbf{\tr{1.\ Um lance}{1.\ One setup}}\\[2pt]
       \scriptsize\tr{nível, mira, ré e vante}{level, rod, backsight and foresight}};
    \node[caixa, draw=corErro, fill=corErro!8, text width=2.8cm, minimum height=2.1cm] (b) at (3.5,0)
      {\textbf{\tr{2.\ O erro que cresce}{2.\ The error that grows}}\\[2pt]
       \scriptsize\tr{$\sqrt{K}$ contra $K$; a tolerância}{$\sqrt{K}$ versus $K$; the tolerance}};
    \node[caixa, draw=corSist, fill=corSist!8, text width=2.8cm, minimum height=2.1cm] (c) at (7.0,0)
      {\textbf{\tr{3.\ A rede}{3.\ The network}}\\[2pt]
       \scriptsize\tr{RRNN, linhas, circuitos, fechamento}{benchmarks, lines, loops,
       misclosure}};
    \node[caixa, draw=corTempo, fill=corTempo!8, text width=2.8cm, minimum height=2.1cm] (d) at (10.5,0)
      {\textbf{\tr{4.\ O que o nível não vê}{4.\ What the level does not see}}\\[2pt]
       \scriptsize\tr{a gravidade; o gancho para a Aula 3}{gravity; the hook to Class 3}};
    \draw[seta] (a) -- (b);
    \draw[seta] (b) -- (c);
    \draw[seta] (c) -- (d);
  \end{tikzpicture}
\end{frame}

% ------------------------------------------------------------------
\begin{frame}{\tr{O nível e a mira}{The level and the rod}}
  \begin{columns}[T]
    \begin{column}{0.64\textwidth}
      \centering
      \begin{tikzpicture}[x=0.78cm,y=0.92cm,font=\scriptsize]
        % terreno
        \draw[thick, mStoneMid] plot[smooth] coordinates {(0,0.30) (1.3,0.35) (3,0.6) (4.5,0.75)
              (6,0.95) (7.7,1.167) (9,1.2) (10,1.1)};
        % miras A (ré) e B (vante)
        \draw[thick, fill=white] (1.2,0.35) rectangle (1.4,3.1);
        \draw[thick, fill=white] (7.6,1.167) rectangle (7.8,3.917);
        \foreach \t in {0.6,0.9,1.2,1.5,1.8,2.1,2.4,2.7} {
          \draw[mStoneMid] (1.2,0.35+\t) -- (1.3,0.35+\t);
          \draw[mStoneMid] (7.6,1.167+\t) -- (7.7,1.167+\t);
        }
        % linha de visada
        \draw[corErro, thick, dashed] (1.3,2.573) -- (7.7,2.573);
        % tripé e nível
        \draw[thick] (4.5,0.75) -- (4.15,2.2) (4.5,0.75) -- (4.85,2.2) (4.5,0.75) -- (4.5,2.2);
        \draw[thick, fill=corSist!25] (4.2,2.2) rectangle (4.8,2.45);
        \draw[thick, fill=corSist!45] (4.15,2.45) rectangle (4.85,2.68);
        \node[corSist, below] at (4.5,0.7) {\tr{nível}{level}};
        % leituras
        \draw[<->, corFrame, thick] (0.95,0.35) -- node[left, text=corFrame] {\textbf{\tr{ré}{back}}} (0.95,2.573);
        \draw[<->, corFrame, thick] (8.05,1.167) -- node[right, text=corFrame] {\textbf{\tr{vante}{fore}}} (8.05,2.573);
        % Delta h
        \draw[dotted] (1.4,0.35) -- (9.3,0.35);
        \draw[dotted] (7.8,1.167) -- (9.3,1.167);
        \draw[<->, corTempo, thick] (9.1,0.35) -- node[right, text=corTempo] {$\Delta h$} (9.1,1.167);
        \node[below] at (1.3,0.3) {$A$};
        \node[below] at (7.7,1.12) {$B$};
      \end{tikzpicture}
    \end{column}
    \begin{column}{0.32\textwidth}
      \footnotesize
      \tr{O \textbf{nível} é um telescópio que, no tripé, fica sempre \bnum{na horizontal}: a
      linha de visada desenha um plano de altura constante.}{The \textbf{level} is a telescope
      that, on its tripod, always stays \bnum{horizontal}: the line of sight draws a plane of
      constant height.}
      \medskip

      \tr{Nas \textbf{miras} lemos a que distância do chão essa linha passa: \bnum{ré} (para o
      ponto de altitude conhecida) e \bnum{vante} (para o ponto novo).}{On the \textbf{rods} we
      read how far above the ground that line passes: \bnum{backsight} (to the point of known
      altitude) and \bnum{foresight} (to the new point).}
      \medskip

      \begin{regra}[title={\tr{Desnível}{Height difference}}]
        \centering$\Delta h=\text{\tr{ré}{back}}-\text{\tr{vante}{fore}}$
      \end{regra}
    \end{column}
  \end{columns}
\end{frame}

\begin{frame}{\tr{Um lance, passo a passo}{One setup, step by step}}
  \begin{columns}[T]
    \begin{column}{0.50\textwidth}
      \small
      \begin{enumerate}
        \setlength{\itemsep}{3pt}
        \item \tr{Aponta-se para a mira em $A$ (altitude \num{12.000}\,m):
              \textbf{ré} $=\num{1.482}$\,m.}{Sight the rod at $A$ (altitude
              \num{12.000}\,m): \textbf{backsight} $=\num{1.482}$\,m.}
        \item \tr{Gira-se para a mira em $B$: \textbf{vante} $=\num{0.937}$\,m.}{Turn to the rod at
              $B$: \textbf{foresight} $=\num{0.937}$\,m.}
        \item \tr{$\Delta h=\num{1.482}-\num{0.937}=\bnum{+\num{0.545}}$\,m.}{$\Delta h=\num{1.482}-\num{0.937}=\bnum{+\num{0.545}}$\,m.}
        \item \tr{$H_B=\num{12.000}+\num{0.545}=\bnum{\num{12.545}}$\,m.}{$H_B=\num{12.000}+\num{0.545}=\bnum{\num{12.545}}$\,m.}
      \end{enumerate}
    \end{column}
    \begin{column}{0.46\textwidth}
      \begin{ideia}[title={\tr{Por que ré $-$ vante?}{Why back $-$ fore?}}]
        \footnotesize\tr{A linha de visada é a mesma nas duas leituras. Quanto \emph{maior} a leitura
        na mira, \emph{mais baixo} está o ponto. Se a ré (em $A$) é maior que a vante (em $B$),
        então $B$ está \emph{mais alto}.}{The line of sight is the same for both readings. The
        \emph{larger} the rod reading, the \emph{lower} the point. If the backsight (at $A$) is
        larger than the foresight (at $B$), then $B$ is \emph{higher}.}
      \end{ideia}
      \medskip

      \begin{cuidado}[title={\tr{Cuidado com o sinal}{Mind the sign}}]
        \footnotesize\tr{Trocar ré e vante troca o sinal do desnível. Por isso a caderneta anota
        \emph{sempre} qual é qual.}{Swapping backsight and foresight swaps the sign of the
        height difference. That is why the field book \emph{always} records which is
        which.}
      \end{cuidado}
    \end{column}
  \end{columns}
\end{frame}

\begin{frame}{\tr{Lance após lance}{Setup after setup}}
  \centering
  \begin{tikzpicture}[x=1cm,y=1cm,font=\small]
    \coordinate (P0) at (0,0.0);
    \coordinate (P1) at (2.3,0.5);
    \coordinate (P2) at (4.6,0.2);
    \coordinate (P3) at (6.9,0.9);
    \coordinate (P4) at (9.2,1.4);
    \draw[thick, mStoneMid] (P0) -- (P1) -- (P2) -- (P3) -- (P4);
    \foreach \p/\l in {P0/RN\,$A$, P1/$1$, P2/$2$, P3/$3$, P4/RN\,$B$} {
      \fill[corFrame] (\p) circle (0.1);
      \node[below=3pt, text=corFrame] at (\p) {\l};
    }
    \foreach \a/\b/\l in {P0/P1/{\Delta h_1}, P1/P2/{\Delta h_2}, P2/P3/{\Delta h_3}, P3/P4/{\Delta h_4}} {
      \draw[seta, corTempo] ($(\a)+(0,1.0)$) -- node[above, text=corTempo] {$\l$} ($(\b)+(0,1.0)$);
    }
  \end{tikzpicture}
  \bigskip

  \begin{regra}[title={\tr{O zero caminha}{The zero walks}}]
    \centering\small
    $H_B=H_A+\Delta h_1+\Delta h_2+\Delta h_3+\Delta h_4=H_A+\sum_i\Delta h_i$
    \smallskip

    \tr{Cada lance tem poucas dezenas de metros. Um nivelamento de \num{1000}\,km são
    \emph{milhares} de lances encadeados --- cada um carrega um pequeno erro.}{Each setup is a
    few tens of meters. A leveling line of \num{1000}\,km is \emph{thousands} of chained setups
    --- each one carries a small error.}
  \end{regra}
\end{frame}

\begin{frame}{\tr{Por que visadas curtas e iguais?}{Why short, equal sights?}}
  \begin{columns}[T]
    \begin{column}{0.54\textwidth}
      \small
      \tr{Três coisas atrapalham, e as três crescem com a distância da visada:}{Three things get
      in the way, and all three grow with the sight distance:}
      \begin{enumerate}
        \setlength{\itemsep}{3pt}
        \item \tr{a \textbf{curvatura da Terra}: a linha horizontal ``sobe'' em relação ao
              chão, $c=s^{2}/2R$;}{the \textbf{curvature of the Earth}: the horizontal line
              ``rises'' relative to the ground, $c=s^{2}/2R$;}
        \item \tr{a \textbf{refração} do ar encurva o raio de luz e come \num{13}\,\% disso;}{the
              \textbf{refraction} of the air bends the light ray and eats \num{13}\,\% of
              that;}
        \item \tr{a \textbf{colimação}: o eixo do nível pode estar levemente inclinado.}{the
              \textbf{collimation}: the level's axis may be slightly tilted.}
      \end{enumerate}
    \end{column}
    \begin{column}{0.42\textwidth}
      \renewcommand{\arraystretch}{1.4}
      \small
      \begin{tabular}{@{}lrr@{}}
        \toprule
        \tr{visada}{sight} & \tr{curvatura}{curvature} & \tr{com refração}{with refraction} \\
        \midrule
        \num{30}\,m & \num{0.07}\,mm & \num{0.06}\,mm \\
        \num{50}\,m & \num{0.20}\,mm & \num{0.17}\,mm \\
        \num{100}\,m & \num{0.79}\,mm & \num{0.68}\,mm \\
        \bottomrule
      \end{tabular}
      \bigskip

      \begin{ideia}[title={\tr{Receita}{Recipe}}]
        \footnotesize\tr{Visadas \bnum{curtas} (erros pequenos) e \bnum{iguais} (o que sobra se
        cancela): o nível vai no \emph{meio}.}{\bnum{Short} sights (small errors) and \bnum{equal}
        ones (what is left cancels): the level goes in the \emph{middle}.}
      \end{ideia}
    \end{column}
  \end{columns}
\end{frame}

\begin{frame}{\tr{Visadas iguais: o erro se cancela}{Equal sights: the error cancels}}
  \begin{columns}[T]
    \begin{column}{0.64\textwidth}
      \centering
      \begin{tikzpicture}[x=0.55cm,y=0.9cm,font=\scriptsize]
        % painel 1: distâncias iguais
        \node[above, font=\small] at (3.6,2.1) {\tr{distâncias \bnum{iguais}}{\bnum{equal} distances}};
        \draw[mStoneMid] (0,0) -- (7.2,0);
        \draw[thick] (0,0) -- (0,1.9) (7.2,0) -- (7.2,1.9);
        \draw[dashed] (0,1.0) -- (7.2,1.0);
        \draw[corErro, thick] (3.6,1.0) -- (0,1.0+0.30) (3.6,1.0) -- (7.2,1.0+0.30);
        \fill[corSist] (3.6,1.0) circle (0.1);
        \draw[<->, corErro] (-0.35,1.0) -- node[left, text=corErro] {$e$} (-0.35,1.30);
        \draw[<->, corErro] (7.55,1.0) -- node[right, text=corErro] {$e$} (7.55,1.30);
        % painel 2: distâncias desiguais
        \node[above, font=\small] at (3.6,-0.9) {\tr{distâncias \bnum{desiguais}}{\bnum{unequal} distances}};
        \begin{scope}[shift={(0,-3.2)}]
          \draw[mStoneMid] (0,0) -- (7.2,0);
          \draw[thick] (0,0) -- (0,1.9) (7.2,0) -- (7.2,1.9);
          \draw[dashed] (0,1.0) -- (7.2,1.0);
          \draw[corErro, thick] (1.8,1.0) -- (0,1.0+0.15) (1.8,1.0) -- (7.2,1.0+0.45);
          \fill[corSist] (1.8,1.0) circle (0.1);
          \draw[<->, corErro] (-0.35,1.0) -- node[left, text=corErro] {$e_1$} (-0.35,1.15);
          \draw[<->, corErro] (7.55,1.0) -- node[right, text=corErro] {$e_2$} (7.55,1.45);
        \end{scope}
      \end{tikzpicture}
    \end{column}
    \begin{column}{0.32\textwidth}
      \footnotesize
      \tr{O eixo inclinado de $10\asec$ (ou a curvatura) faz as duas leituras ficarem
      \emph{altas demais}, em proporção à distância.}{An axis tilted by $10\asec$ (or the
      curvature) makes both readings come out \emph{too high}, in proportion to the
      distance.}
      \medskip

      \tr{\bnum{Iguais}: $e$ nas duas; em $\text{ré}-\text{vante}$ o erro some.}{\bnum{Equal}:
      $e$ on both; in $\text{back}-\text{fore}$ the error vanishes.}
      \medskip

      \tr{\bnum{Desiguais}: sobra $e_2-e_1$. Com $10\asec$ e \num{20}\,m de diferença, é
      \bnum{\num{0.97}\,mm por lance} --- em 100 lances, quase \bnum{10\,cm}.}{\bnum{Unequal}:
      $e_2-e_1$ is left. With $10\asec$ and a \num{20}\,m difference, that is \bnum{\num{0.97}\,mm
      per setup} --- over 100 setups, almost \bnum{10\,cm}.}
    \end{column}
  \end{columns}
\end{frame}

\begin{frame}{\tr{Ida e volta: o nivelamento duplo}{Forward and back: double leveling}}
  \begin{columns}[T]
    \begin{column}{0.50\textwidth}
      \small
      \tr{Cada \textbf{seção} entre duas RRNN é nivelada \bnum{duas vezes}: na ida
      (\emph{nivelamento}) e na volta (\emph{contranivelamento}).}{Each \textbf{section}
      between two benchmarks is leveled \bnum{twice}: forward (\emph{leveling}) and back
      (\emph{counter-leveling}).}
      \medskip

      \renewcommand{\arraystretch}{1.4}
      \begin{tabular}{@{}lr@{}}
        \toprule
        \tr{ida}{forward} & $+\num{12.347}$\,m \\
        \tr{volta}{back} & $-\num{12.343}$\,m \\
        \tr{discrepância}{discrepancy} & \bnum{\num{4}\,mm} \\
        \tr{média}{mean} & $\num{12.345}$\,m \\
        \bottomrule
      \end{tabular}
      \smallskip

      {\footnotesize\tr{Seção de $K=\num{2.8}$\,km: tolerância $3\sqrt{K}=\num{5.0}$\,mm. Aceita.}{Section
      with $K=\num{2.8}$\,km: tolerance $3\sqrt{K}=\num{5.0}$\,mm. Accepted.}}
    \end{column}
    \begin{column}{0.46\textwidth}
      \begin{ideia}[title={\tr{Para que serve}{What it is for}}]
        \small
        \begin{itemize}
          \setlength{\itemsep}{3pt}
          \item \tr{a \bnum{discrepância} é um teste: se for grande, algo deu errado (um erro
                grosseiro, vento, mira mal apoiada) e a seção é \emph{refeita};}{the
                \bnum{discrepancy} is a test: if it is large, something went wrong (a blunder,
                wind, a badly seated rod) and the section is \emph{redone};}
          \item \tr{a \bnum{média} das duas medidas é melhor que cada uma;}{the \bnum{mean} of
                the two measurements is better than either;}
          \item \tr{erros que mudam de sinal na volta (como alguns efeitos da temperatura) se
                cancelam.}{errors that change sign on the way back (like some temperature
                effects) cancel.}
        \end{itemize}
      \end{ideia}
    \end{column}
  \end{columns}
\end{frame}

\begin{frame}{\tr{As ferramentas de alta precisão}{High-precision tools}}
  \begin{columns}[T]
    \begin{column}{0.48\textwidth}
      \begin{regra}[title={\tr{A mira de invar}{The invar rod}}]
        \scriptsize\tr{O \textbf{invar} é uma liga de ferro e níquel que \emph{quase não dilata} com
        o calor. A graduação fica numa fita de invar esticada na mira: o comprimento do
        ``centímetro'' não muda com a temperatura do dia.}{\textbf{Invar} is an iron-nickel
        alloy that \emph{hardly expands} with heat. The graduation sits on an invar strip
        stretched along the rod: the length of the ``centimeter'' does not change with the
        day's temperature.}
      \end{regra}
      \medskip
      \begin{regra}[title={\tr{O nível com micrômetro}{The level with a micrometer}}]
        \scriptsize\tr{Uma placa de vidro plano-paralela, girada por um parafuso, desloca a imagem
        da mira para ``encostar'' no traço mais próximo: lê-se fração de milímetro.}{A
        plane-parallel glass plate, turned by a screw, shifts the rod image to ``touch'' the
        nearest line: a fraction of a millimeter is read.}
      \end{regra}
    \end{column}
    \begin{column}{0.48\textwidth}
      \begin{ideia}[title={\tr{O resto é método}{The rest is method}}]
        \scriptsize
        \begin{itemize}
          \setlength{\itemsep}{2pt}
          \item \tr{nível \bnum{sempre no meio} (visadas iguais);}{level \bnum{always in the
                middle} (equal sights);}
          \item \tr{mira \bnum{apoiada em sapata} (placa metálica), na vertical;}{rod \bnum{on a
                foot plate}, held vertical;}
          \item \tr{guarda-sol sobre o instrumento: o Sol deforma o nível;}{umbrella over the
                instrument: the Sun deforms the level;}
          \item \tr{horas sem vento forte; \emph{ida} e \emph{volta} em horários e condições
                diferentes;}{hours without strong wind; \emph{forward} and \emph{back} at
                different times and conditions;}
          \item \tr{leituras anotadas com \emph{décimos de milímetro}.}{readings written down
                to \emph{tenths of a millimeter}.}
        \end{itemize}
      \end{ideia}
    \end{column}
  \end{columns}
\end{frame}

\begin{frame}{\tr{Experimente: o simulador de nivelamento}{Try it: the leveling simulator}}
  \begin{columns}[T]
    \begin{column}{0.55\textwidth}
      \footnotesize
      \tr{No simulador \texttt{nivelamento/}, cada trecho entre dois pontos recebe um ruído
      \emph{gaussiano de desvio} $\sigma\sqrt{d}$ --- é o $\eta\sqrt{K}$ desta aula.}{In the
      \texttt{nivelamento/} simulator, each section between two points receives \emph{Gaussian
      noise with deviation} $\sigma\sqrt{d}$ --- the $\eta\sqrt{K}$ of this class.}
      \medskip
      \begin{enumerate}
        \setlength{\itemsep}{2pt}
        \item \tr{Gere uma rede (\emph{Gerar Nova Rede}); um ponto é fixo: o datum.}{Generate a
              network (\emph{Generate New Network}); one point is fixed: the datum.}
        \item \tr{Ponha $\sigma$ em \num{1}\,mm$/\sqrt{\text{km}}$ e rode (\emph{Executar
              Ajustamento}).}{Set $\sigma$ to \num{1}\,mm$/\sqrt{\text{km}}$ and run
              (\emph{Run Adjustment}).}
        \item \tr{Dobre o $\sigma$: os resíduos dobram.}{Double $\sigma$: the residuals
              double.}
        \item \tr{Use \emph{Simular Erro Grosseiro} e veja o teste achar o trecho
              errado.}{Use \emph{Simulate Gross Error (Blunder)} and watch the test find the bad
              section.}
      \end{enumerate}
    \end{column}
    \begin{column}{0.41\textwidth}
      \begin{ideia}[title={\tr{O que olhar}{What to look at}}]
        \footnotesize
        \begin{itemize}
          \setlength{\itemsep}{3pt}
          \item \tr{\emph{Resíduos \& Data Snooping}: onde está o erro grosseiro?}{\emph{Residuals
                \& Data Snooping}: where is the blunder?}
          \item \tr{\emph{Teste Global} ($\chi^2$): a rede é compatível com o $\sigma$
                dado?}{\emph{Global Test} ($\chi^2$): is the network compatible with the given
                $\sigma$?}
          \item \tr{\emph{Monte Carlo}: repita o sorteio e veja a nuvem de erros.}{\emph{Monte
                Carlo}: repeat the draw and see the cloud of errors.}
        \end{itemize}
      \end{ideia}
      \medskip
      {\scriptsize\tr{A matemática do ajustamento é da Aula 3.}{The mathematics of the adjustment
      is in Class 3.}}
    \end{column}
  \end{columns}
\end{frame}

\begin{frame}{\tr{Sua vez}{Your turn}}
  \begin{columns}[T]
    \begin{column}{0.50\textwidth}
      \begin{regra}[title={\tr{Um lance}{One setup}}]
        \small
        \begin{enumerate}
          \item[(a)] \tr{$A$ tem altitude \num{100.000}\,m. Ré em $A$: \num{1.612}\,m; vante em
                $B$: \num{2.034}\,m. Qual a altitude de $B$? $B$ é mais alto ou mais baixo que
                $A$?}{$A$ has altitude \num{100.000}\,m. Backsight at $A$: \num{1.612}\,m;
                foresight at $B$: \num{2.034}\,m. What is the altitude of $B$? Is $B$ higher or
                lower than $A$?}
          \item[(b)] \tr{Se o nivelador trocar ré e vante sem querer, que altitude obtém para
                $B$?}{If the leveler swaps backsight and foresight by mistake, what altitude does
                he get for $B$?}
        \end{enumerate}
      \end{regra}
    \end{column}
    \begin{column}{0.46\textwidth}
      \onslide<2->{
      \begin{ideia}[title={\tr{Respostas}{Answers}}]
        \small\raggedright
        (a) $\Delta h=\num{1.612}-\num{2.034}=-\num{0.422}$\,m; $H_B=\bnum{\num{99.578}\,m}$:
        \tr{$B$ é \emph{mais baixo}.}{$B$ is \emph{lower}.}
        \medskip

        (b) $\Delta h=+\num{0.422}$\,m; $H_B=\bnum{\num{100.422}\,m}$ --- \tr{errado por
        $2\times\num{0.422}=\num{0.844}$\,m!}{wrong by $2\times\num{0.422}=\num{0.844}$\,m!}
        \medskip

        \tr{Um erro de sinal é um erro \alert{grosseiro}: o ajustamento (Aula 3) o
        descobre.}{A sign error is a \alert{blunder}: the adjustment (Class 3) finds it.}
      \end{ideia}
      }
    \end{column}
  \end{columns}
\end{frame}

% ------------------------------------------------------------------
\begin{frame}{\tr{Erros pequenos, caminho longo}{Small errors, long road}}
  \begin{columns}[T]
    \begin{column}{0.47\textwidth}
      \begin{ideia}[title={\tr{Erro aleatório}{Random error}}]
        \small\tr{Cada lance erra um pouquinho, \emph{ora para cima, ora para baixo}, sem
        padrão. Os erros se compensam em parte: o total cresce \bnum{devagar}.}{Each setup is
        off by a little, \emph{sometimes up, sometimes down}, with no pattern. The errors
        partly compensate: the total grows \bnum{slowly}.}
        \medskip

        \centering
        \begin{tikzpicture}[x=0.5cm,y=0.5cm]
          \foreach \x/\d in {0/1,1/-1,2/1,3/1,4/-1,5/-1,6/1,7/-1,8/1,9/-1} {
            \draw[seta, corFrame] ({\x},0) -- ({\x},{0.9*\d});
          }
        \end{tikzpicture}
      \end{ideia}
    \end{column}
    \begin{column}{0.47\textwidth}
      \begin{armadilha}[title={\tr{Erro sistemático}{Systematic error}}]
        \small\tr{Cada lance erra \emph{sempre para o mesmo lado} (mira com a graduação
        ``esticada'', refração, colimação). Os erros \emph{se somam}: o total cresce
        \bnum{depressa}.}{Each setup is off \emph{always to the same side} (a rod with a
        ``stretched'' graduation, refraction, collimation). The errors \emph{add up}: the total
        grows \bnum{fast}.}
        \medskip

        \centering
        \begin{tikzpicture}[x=0.5cm,y=0.5cm]
          \foreach \x in {0,1,2,3,4,5,6,7,8,9} {
            \draw[seta, corErro] ({\x},0) -- ({\x},0.9);
          }
        \end{tikzpicture}
      \end{armadilha}
    \end{column}
  \end{columns}
  \bigskip

  \centering
  \tr{Num nivelamento de \num{1000}\,km são milhares de lances. \bnum{Qual dos dois manda}?
  Depende de quantos quilômetros.}{A \num{1000}\,km leveling line is thousands of setups.
  \bnum{Which one rules}? It depends on how many kilometers.}
\end{frame}

\begin{frame}{\tr{Erro aleatório: cresce com a raiz da distância}{Random error: grows with the
square root of the distance}}
  \begin{columns}[T]
    \begin{column}{0.62\textwidth}
      \centering
      \begin{tikzpicture}[x=0.78cm,y=0.95cm,font=\scriptsize]
        \draw[->] (-0.1,0) -- (7.5,0) node[right] {$K$ (km)};
        \draw[->] (0,-2.7) -- (0,2.8) node[above] {\tr{erro acumulado (mm)}{accumulated error (mm)}};
        \foreach \x in {20,40,60,80,100} {
          \draw ({0.07*\x},0.05) -- ({0.07*\x},-0.05) node[below] {\x};
        }
        \foreach \y in {-20,-10,10,20} {
          \draw (0.05,{0.12*\y}) -- (-0.05,{0.12*\y}) node[left] {\y};
        }
        \passeioA{corFrame!60}
        \passeioB{corFrame!60}
        \passeioC{corFrame!60}
        \passeioD{corFrame!60}
        \passeioE{corFrame!60}
        \passeioF{corFrame!60}
        \draw[corErro, very thick, domain=0:100, samples=40, smooth, variable=\k]
              plot ({0.07*\k},{0.12*sqrt(\k)});
        \draw[corErro, very thick, domain=0:100, samples=40, smooth, variable=\k]
              plot ({0.07*\k},{-0.12*sqrt(\k)});
        \node[corErro, above left] at (7.0,1.2) {$\pm\,\eta\sqrt{K}$};
      \end{tikzpicture}\\[1pt]
      {\scriptsize\tr{Seis linhas simuladas ($\eta=\num{1}$\,mm$/\sqrt{\text{km}}$).}{Six simulated
      lines ($\eta=\num{1}$\,mm$/\sqrt{\text{km}}$).}}
    \end{column}
    \begin{column}{0.36\textwidth}
      \small
      \tr{Cada quilômetro soma um erro de \emph{tamanho típico} $\eta$, com sinal ao acaso.
      Depois de $K$ km, o desvio típico é}{Each kilometer adds an error of \emph{typical size}
      $\eta$, with a random sign. After $K$ km, the typical deviation is}
      \[
        \sigma=\eta\sqrt{K}
      \]
      \tr{$\eta$ em mm$/\sqrt{\text{km}}$ é a ``qualidade'' do nivelamento.}{$\eta$ in
      mm$/\sqrt{\text{km}}$ is the ``quality'' of the leveling.}
      \medskip

      \tr{Quadruplicar a distância \bnum{só dobra} o erro típico.}{Quadrupling the distance
      \bnum{only doubles} the typical error.}
    \end{column}
  \end{columns}
\end{frame}

\begin{frame}{\tr{Erro sistemático: cresce com a distância}{Systematic error: grows with the
distance}}
  \begin{columns}[T]
    \begin{column}{0.62\textwidth}
      \centering
      \begin{tikzpicture}[x=0.86cm,y=0.95cm,font=\scriptsize]
        \draw[->] (-0.1,0) -- (7.5,0) node[right] {$K$ (km)};
        \draw[->] (0,-0.1) -- (0,4.2) node[above] {mm};
        \foreach \x in {100,200,300,400} {
          \draw ({0.0175*\x},0.05) -- ({0.0175*\x},-0.05) node[below] {\x};
        }
        \foreach \y in {10,20,30,40} {
          \draw (0.05,{0.09*\y}) -- (-0.05,{0.09*\y}) node[left] {\y};
        }
        \draw[corFrame, very thick, domain=0:400, samples=40, variable=\k]
              plot ({0.0175*\k},{0.09*sqrt(\k)});
        \draw[corErro, very thick] (0,0) -- ({0.0175*400},{0.09*40});
        \fill[mStone] ({0.0175*100},{0.09*10}) circle (0.07);
        \draw[dashed] ({0.0175*100},0) -- ({0.0175*100},{0.09*10}) -- (0,{0.09*10});
        \node[corFrame, below right] at ({0.0175*250},{0.09*sqrt(250)-0.05}) {$\eta\sqrt{K}$, $\eta=\num{1}$};
        \node[corErro, above left] at ({0.0175*330},{0.09*33}) {$\delta K$, $\delta=\num{0.1}$};
        \node[above right] at ({0.0175*100+0.1},{0.09*10+0.05}) {\num{100}\,km};
      \end{tikzpicture}\\[1pt]
      {\scriptsize\tr{Raiz ($\eta=\num{1}$\,mm$/\sqrt{\text{km}}$) contra reta ($\delta=\num{0.1}$\,mm/km).}{Square
      root ($\eta=\num{1}$\,mm$/\sqrt{\text{km}}$) versus line ($\delta=\num{0.1}$\,mm/km).}}
    \end{column}
    \begin{column}{0.36\textwidth}
      \small
      \tr{Se cada quilômetro erra $\delta$ \emph{sempre para o mesmo lado}, o erro em $K$ km
      é}{If each kilometer is off by $\delta$ \emph{always to the same side}, the error over $K$
      km is}
      \[
        e=\delta\,K
      \]
      \tr{Dobrar a distância \bnum{dobra} o erro.}{Doubling the distance \bnum{doubles} the
      error.}
      \medskip

      \tr{As curvas se cruzam em $K^{*}=(\eta/\delta)^{2}=\bnum{\num{100}\,km}$: dali em diante
      \errr{o sistemático vence}.}{The curves cross at $K^{*}=(\eta/\delta)^{2}=\bnum{\num{100}\,km}$:
      from there on \errr{the systematic wins}.}
    \end{column}
  \end{columns}
\end{frame}

\begin{frame}{\tr{Os dois, lado a lado}{Side by side}}
  \begin{columns}[T]
    \begin{column}{0.56\textwidth}
      \renewcommand{\arraystretch}{1.55}
      \small
      \begin{tabular}{@{}rrr@{}}
        \toprule
        $K$ (km) & \tr{aleatório $\eta\sqrt{K}$}{random $\eta\sqrt{K}$} & \tr{sistemático
        $\delta K$}{systematic $\delta K$} \\
        \midrule
        \num{4} & \num{2.0}\,mm & \num{0.4}\,mm \\
        \num{100} & \num{10.0}\,mm & \num{10.0}\,mm \\
        \num{400} & \num{20.0}\,mm & \num{40.0}\,mm \\
        \num{2000} & \bnum{\num{44.7}\,mm} & \errr{\num{200}\,mm} \\
        \bottomrule
      \end{tabular}
      \medskip

      {\footnotesize\tr{$\eta=\num{1}$\,mm$/\sqrt{\text{km}}$; $\delta=\num{0.1}$\,mm/km (didáticos).}{$\eta=\num{1}$\,mm$/\sqrt{\text{km}}$;
      $\delta=\num{0.1}$\,mm/km (didactic).}}
    \end{column}
    \begin{column}{0.40\textwidth}
      \begin{cuidado}[title={\tr{Moral}{Moral}}]
        \small\tr{Num trecho curto, o erro \emph{aleatório} é o que se vê. Numa linha de
        \num{2000}\,km, um erro \emph{sistemático} de só \num{0.1}\,mm por quilômetro faz
        \bnum{\num{20}\,cm}.}{On a short stretch, \emph{random} error is what you see. On a
        \num{2000}\,km line, a \emph{systematic} error of only \num{0.1}\,mm per kilometer makes
        \bnum{\num{20}\,cm}.}
        \medskip

        \tr{Por isso o método --- visadas iguais, invar, calibração, ida e volta --- é
        sobretudo uma \bnum{caça ao sistemático}.}{That is why the method --- equal sights,
        invar, calibration, forward and back --- is above all a \bnum{hunt for the
        systematic}.}
      \end{cuidado}
    \end{column}
  \end{columns}
\end{frame}

\begin{frame}{\tr{Tolerância: quando refazer uma seção}{Tolerance: when to redo a section}}
  \begin{columns}[T]
    \begin{column}{0.49\textwidth}
      \small
      \tr{Uma seção liga duas RRNN. Os desníveis de ida (1) e volta (2) têm
      sinais opostos:}{A section joins two benchmarks. Forward (1) and backward (2)
      height differences have opposite signs:}
      \[
        |d|=|\Delta h_1+\Delta h_2|
        \le 3\,\mathrm{mm}\sqrt K
      \]
      \footnotesize
      \tr{$K$ em km, contado uma vez. Seções de alta precisão fundamental: até 3 km.}{
      $K$ is the one-way length in km. Fundamental high-precision sections: at most 3 km.}
      \par\medskip
      \begin{tabular}{@{}rr@{}}
        \toprule
        $K$ (km) & \tr{limite de discrepância}{discrepancy limit}\\
        \midrule
        \num{0.25} & \num{1.5}\,mm\\
        \num{1} & \num{3.0}\,mm\\
        \num{2} & \num{4.2}\,mm\\
        \num{3} & \num{5.2}\,mm\\
        \bottomrule
      \end{tabular}
    \end{column}
    \begin{column}{0.47\textwidth}
      \begin{ideia}[title={\tr{O erro-padrão após o ajuste}{Standard error after adjustment}}]
        \footnotesize
        \tr{O valor \bnum{$2\,\mathrm{mm}\sqrt K$} é o erro-padrão máximo de uma
        \emph{linha após o ajustamento}. Ele avalia outra grandeza.}{The value
        \bnum{$2\,\mathrm{mm}\sqrt K$} is the maximum standard error of an
        \emph{adjusted line}. It evaluates a different quantity.}
      \end{ideia}
      \par\medskip
      \begin{cuidado}[title={\tr{Linha e circuito}{Lines and circuits}}]
        \footnotesize
        \tr{Discrepância de ida/volta de uma \bnum{linha}: $4\,\mathrm{mm}\sqrt K$.
        Num \bnum{circuito}, a discrepância acumulada dividida pelo perímetro
        deve ser no máximo \num{0.5}\,mm/km.}{Forward/backward discrepancy of a
        \bnum{line}: $4\,\mathrm{mm}\sqrt K$. For a \bnum{circuit}, accumulated
        discrepancy divided by perimeter must not exceed \num{0.5}\,mm/km.}
      \end{cuidado}
    \end{column}
  \end{columns}
  \fonteslide{IBGE (1983), R.PR 22/1983, \tr{Quadro VI, itens 3.1--3.3 e 4, p. 30}{Table VI, items 3.1--3.3 and 4, p. 30}.}
\end{frame}

\begin{frame}{\tr{Sua vez}{Your turn}}
  \begin{columns}[T]
    \begin{column}{0.50\textwidth}
      \begin{regra}[title={\tr{Aceitar ou refazer?}{Accept or redo?}}]
        \small
        \begin{enumerate}
          \item[(a)] \tr{Uma seção de $K=\num{2}$\,km deu ida $+\num{8.912}$\,m e volta
                $-\num{8.904}$\,m. Qual a discrepância? Passa na tolerância
                $3\sqrt{K}$?}{A section of $K=\num{2}$\,km gave forward $+\num{8.912}$\,m and
                back $-\num{8.904}$\,m. What is the discrepancy? Does it pass the
                $3\sqrt{K}$ tolerance?}
          \item[(b)] \tr{Numa linha de \num{2000}\,km, qual o erro típico aleatório
                ($\eta=\num{1}$)? E o sistemático com $\delta=\num{0.1}$\,mm/km? Qual é maior,
                quantas vezes?}{On a \num{2000}\,km line, what is the typical random error
                ($\eta=\num{1}$)? And the systematic one with $\delta=\num{0.1}$\,mm/km? Which
                is larger, by what factor?}
        \end{enumerate}
      \end{regra}
    \end{column}
    \begin{column}{0.46\textwidth}
      \onslide<2->{
      \begin{ideia}[title={\tr{Respostas}{Answers}}]
        \small\raggedright
        (a) $|\num{8.912}-\num{8.904}|=\bnum{\num{8}\,mm}$; $3\sqrt{2}=\num{4.2}$\,mm.
        \tr{$8>4{,}2$: \bnum{refazer} a seção.}{$8>4.2$: \bnum{redo} the section.}
        \medskip

        (b) $\sqrt{2000}=\num{44.7}$\,mm; \ $\num{0.1}\times2000=\num{200}$\,mm; \
        \tr{o sistemático é \bnum{\num{4.5}} vezes maior.}{the systematic one is \bnum{\num{4.5}}
        times larger.}
      \end{ideia}
      }
    \end{column}
  \end{columns}
\end{frame}

\begin{frame}{\tr{Pare e pense}{Stop and think}}
  \begin{columns}[T]
    \begin{column}{0.46\textwidth}
      \begin{armadilha}[title={\tr{Uma ideia tentadora}{A tempting idea}}]
        \small\tr{``O erro aleatório diminui com a média. Então é só \emph{nivelar a mesma linha
        cem vezes} e o erro some.''}{``Random error shrinks with averaging. So just \emph{level
        the same line a hundred times} and the error vanishes.''}
        \medskip

        \tr{O que está certo e o que está errado?}{What is right and what is wrong?}
      \end{armadilha}
    \end{column}
    \begin{column}{0.50\textwidth}
      \onslide<2->{
      \begin{ideia}[title={\tr{Resposta}{Answer}}]
        \footnotesize
        \begin{itemize}
          \setlength{\itemsep}{2pt}
          \item \tr{O \emph{aleatório} cai com $1/\sqrt{n}$: cem vezes dá \bnum{dez vezes}
                menos.}{The \emph{random} part falls as $1/\sqrt{n}$: a hundred times gives
                \bnum{ten times} less.}
          \item \tr{O \emph{sistemático} \bnum{não cai}: se a mira está ``esticada'', as cem
                medidas erram igual.}{The \emph{systematic} part \bnum{does not fall}: if the
                rod is ``stretched'', all hundred measurements are off equally.}
          \item \tr{Por isso se \emph{varia} tudo: ida e volta, horários, miras, instrumentos;
                e se calibra a mira num comparador.}{That is why everything is \emph{varied}:
                forward and back, times, rods, instruments; and the rod is calibrated on a
                comparator.}
        \end{itemize}
      \end{ideia}
      }
    \end{column}
  \end{columns}
\end{frame}

% ------------------------------------------------------------------
\begin{frame}{\tr{A RN: a pedra que guarda a altitude}{The benchmark: the stone that keeps the
altitude}}
  \begin{columns}[T]
    \begin{column}{0.56\textwidth}
      \small
      \begin{itemize}
        \setlength{\itemsep}{4pt}
        \item \tr{Uma \bnum{referência de nível} (RN) é um marco --- em geral uma chapa
              metálica num pilar de concreto ou numa obra --- com \bnum{altitude conhecida}.}{A
              \bnum{benchmark} (RN) is a mark --- usually a metal plaque on a concrete pillar or
              on a structure --- with a \bnum{known altitude}.}
        \item \tr{Elas são postas \emph{ao longo de estradas e ferrovias}, de poucos em poucos
              quilômetros, para o nivelador chegar de carro e de mira.}{They are placed
              \emph{along roads and railways}, every few kilometers, so the leveler can arrive
              with a car and a rod.}
        \item \tr{Cada RN tem um nome (como \textbf{4X} ou \textbf{9329T}) e uma \emph{ficha} no
              \textbf{Banco de Dados Geodésicos} (BDG) do IBGE: descrição, posição, altitude.}{Each RN
              has a name (like \textbf{4X} or \textbf{9329T}) and a \emph{sheet} in IBGE's
              \textbf{Geodetic Database} (BDG): description, position, altitude.}
      \end{itemize}
    \end{column}
    \begin{column}{0.40\textwidth}
      \centering
      \begin{tikzpicture}[x=0.95cm,y=0.95cm,font=\scriptsize]
        \fill[mStoneMid!20] (0,0) rectangle (5.4,1.0);
        \draw[thick, mStoneMid] (0,1.0) -- (5.4,1.0);
        \draw[corMar, dashed, thick] (0,0.2) -- (5.4,0.2);
        \node[corMar, above right] at (0.05,0.2) {\tr{zero (datum)}{zero (datum)}};
        \draw[thick, fill=mStoneLight] (2.3,0) rectangle (3.5,1.9);
        \draw[thick, fill=corTempo!45] (2.45,1.9) rectangle (3.35,2.02);
        \fill[corFrame] (2.9,2.02) circle (0.08);
        \node[corFrame, above] at (2.9,2.15) {RN};
        \draw[corTempo, thick, <->] (4.4,0.2) -- node[right] {$H$} (4.4,2.02);
        \draw[dotted] (3.4,2.02) -- (4.5,2.02);
        \node[left, align=right] at (2.2,1.45) {\tr{pilar de}{concrete}\\\tr{concreto}{pillar}};
        \node[left, align=right] at (2.2,2.35) {\tr{chapa com}{plaque with}\\\tr{o ponto}{the point}};
        \draw[->, thin] (2.25,2.3) -- (2.45,2.05);
      \end{tikzpicture}
    \end{column}
  \end{columns}
\end{frame}

\begin{frame}{\tr{A rede em três escalas}{The network at three scales}}
  \begin{columns}[T]
    \begin{column}{0.58\textwidth}
      \centering
      \begin{tikzpicture}[x=0.86cm,y=0.74cm,font=\scriptsize,
          rn/.style={postaction={decorate, decoration={markings, mark=between positions 0.05 and 1 step 0.45cm with {\fill[corFrame] (0,0) circle (1.4pt);}}}}]
        % circuitos (preenchimento leve)
        \fill[corFrame!8] (0.4,0.4) -- (4.4,0.2) -- (4.1,2.1) -- (0.4,1.9) -- cycle;
        \fill[corFrame!14] (4.4,0.2) -- (7.1,0.9) -- (7.2,2.2) -- (4.1,2.1) -- cycle;
        \fill[corFrame!8] (0.4,1.9) -- (4.1,2.1) -- (3.8,4.0) -- (0.3,3.3) -- cycle;
        \fill[corFrame!14] (4.1,2.1) -- (7.2,2.2) -- (7.3,3.5) -- (3.8,4.0) -- cycle;
        % linhas (ao longo das estradas)
        \draw[rede, rn] (0.4,0.4) -- (4.4,0.2) -- (7.1,0.9) -- (7.2,2.2) -- (7.3,3.5) -- (3.8,4.0) -- (0.3,3.3) -- (0.4,1.9) -- cycle;
        \draw[rede, rn] (0.4,1.9) -- (4.1,2.1) -- (7.2,2.2);
        \draw[rede, rn] (4.4,0.2) -- (4.1,2.1) -- (3.8,4.0);
        % macrocircuito
        \draw[corSist, very thick, dashed] (0.1,0.15) -- (4.5,-0.1) -- (7.4,0.75) -- (7.55,3.65) -- (3.8,4.3) -- (0.0,3.45) -- cycle;
        % datum
        \draw[corMar, very thick] (-0.5,0.0) .. controls (-0.1,1.5) and (-0.7,2.8) .. (-0.2,4.2);
        \node[corMar, rotate=90] at (-0.7,2.1) {\tr{mar}{sea}};
        \fill[corMar] (0.4,0.4) circle (0.09);
        \node[corMar, below right, font=\tiny] at (0.4,0.35) {\tr{datum}{datum}};
        % rótulos
        \node[corFrame] at (2.2,1.1) {\tr{circuito}{loop}};
        \node[corSist, above] at (3.9,4.3) {\tr{macrocircuito}{macro-loop}};
      \end{tikzpicture}
    \end{column}
    \begin{column}{0.38\textwidth}
      \scriptsize
      \begin{description}
        \setlength{\itemsep}{2pt}
        \item[\textcolor{corFrame}{\textbf{\tr{linha}{line}}}] \tr{RRNN seguidas ao longo de uma
              estrada, ligadas por seções niveladas ida e volta.}{benchmarks in sequence along a
              road, connected by sections leveled forward and back.}
        \item[\textcolor{corFrame}{\textbf{\tr{circuito}{loop}}}] \tr{linhas que \emph{fecham uma
              volta}: voltam ao ponto de partida.}{lines that \emph{close a turn}: they return
              to the starting point.}
        \item[\textcolor{corSist}{\textbf{\tr{macrocircuito}{macro-loop}}}] \tr{um circuito
              grande, que contém vários: foi usado em 1991 para dividir a rede em partes do
              ajustamento.}{a big loop containing several: used in 1991 to split the network
              into adjustment pieces.}
      \end{description}
      \smallskip
      {\tiny\tr{Esquema: o desenho não representa nenhuma região real.}{Sketch: the drawing
      does not represent any real region.}}
    \end{column}
  \end{columns}
\end{frame}

\begin{frame}{\tr{Fechamento de circuito}{Loop misclosure}}
  \begin{columns}[T]
    \begin{column}{0.54\textwidth}
      \centering
      \begin{tikzpicture}[x=0.9cm,y=0.9cm,font=\scriptsize]
        \coordinate (A) at (0,0);
        \coordinate (B) at (4.4,0.4);
        \coordinate (C) at (4.2,3.0);
        \coordinate (D) at (-0.2,2.6);
        \draw[seta, corTempo] (A) -- node[below, text=corTempo, align=center] {$+\num{15.4301}$\,m\\$K=\num{10}$\,km} (B);
        \draw[seta, corTempo] (B) -- node[right, text=corTempo, align=left] {$-\num{4.2276}$\,m\\$K=\num{8}$\,km} (C);
        \draw[seta, corTempo] (C) -- node[above, text=corTempo, align=center] {$-\num{13.7819}$\,m\\$K=\num{12}$\,km} (D);
        \draw[seta, corTempo] (D) -- node[left, text=corTempo, align=right] {$+\num{2.5854}$\,m\\$K=\num{9}$\,km} (A);
        \foreach \p/\l in {A/A, B/B, C/C, D/D} {
          \fill[corFrame] (\p) circle (0.1);
        }
        \node[corFrame, below left] at (A) {A};
        \node[corFrame, below right] at (B) {B};
        \node[corFrame, above right] at (C) {C};
        \node[corFrame, above left] at (D) {D};
      \end{tikzpicture}
    \end{column}
    \begin{column}{0.42\textwidth}
      \footnotesize
      \tr{Neste exemplo fictício, os desníveis já corrigidos deveriam somar
      \bnum{zero}.}{In this fictitious example, the corrected height differences
      should sum to \bnum{zero}.}
      \medskip

      \tr{Aqui: $w=\sum\Delta h=\bnum{+\num{6.0}\,mm}$, perímetro $K=\num{39}$\,km.}{Here:
      $w=\sum\Delta h=\bnum{+\num{6.0}\,mm}$, perimeter $K=\num{39}$\,km.}
      \smallskip

      \tr{Com as correções aplicadas: limite $\num{0.5}K=\num{19.5}$\,mm.
      $|w|/K=\num{0.154}$\,mm/km: \bnum{passa no controle de circuito}.}{With the corrections
      applied: limit $\num{0.5}K=\num{19.5}$\,mm. $|w|/K=\num{0.154}$\,mm/km:
      \bnum{passes the circuit check}.}
      \fonteslide{IBGE (1983), VI.3.3, p. 30.}
      \medskip

      \begin{cuidado}[title={\tr{Necessário, não suficiente}{Necessary, not sufficient}}]
        \scriptsize\tr{Fechar bem não garante a ausência de erro: um erro que depende do
        desnível (como a mira ``esticada'') \emph{sobe e desce} e se cancela na volta.}{Closing
        well does not guarantee there is no error: an error that depends on the height
        difference (like the ``stretched'' rod) \emph{goes up and down} and cancels on the
        way back.}
      \end{cuidado}
    \end{column}
  \end{columns}
\end{frame}

\begin{frame}{\tr{Os números da rede e seus universos}{Network counts and their scope}}
  \begin{columns}[T]
    \begin{column}{0.52\textwidth}
      \small
      \renewcommand{\arraystretch}{1.5}
      \begin{tabular}{@{}lr@{}}
        \toprule
        \tr{Conjunto}{Set} & \tr{RRNN}{Benchmarks}\\
        \midrule
        AAGP (1992) & \num{43079}\\
        \tr{Ajuste de 2011}{2011 adjustment} & \num{69590}\\
        \tr{REALT-2018, principal}{REALT-2018, principal} & \num{66882}\\
        \bottomrule
      \end{tabular}
      \par\medskip
      \footnotesize\tr{Bloco principal: \num{70426} seções; \num{62782} resíduos
      em seções formando circuitos.}{Principal block: \num{70426} sections;
      \num{62782} residuals in circuit-forming sections.}
    \end{column}
    \begin{column}{0.44\textwidth}
      \begin{cuidado}[title={\tr{Qual é o conjunto contado?}{Which set is being counted?}}]
        \footnotesize
        \tr{\num{66882} é o \bnum{bloco principal}, vinculado a Imbituba.
        O relatório trata separadamente as linhas rejeitadas desse bloco e a rede
        do Amapá. Esse número não é um total nacional de RRNN.}{
        \num{66882} counts the \bnum{principal block}, tied to Imbituba.
        The report treats lines excluded from that block and the Amapá network
        separately. This is not a national benchmark total.}
      \end{cuidado}
      \par\medskip
      \footnotesize\tr{Os três conjuntos têm abrangências diferentes. Sua comparação
      direta não mede crescimento nem perda de marcos.}{
      The three sets have different scopes. A direct comparison does not measure
      network growth or benchmark loss.}
    \end{column}
  \end{columns}
  \fonteslide{IBGE (2011), \tr{pp. 2 e 12}{pp. 2 and 12}; IBGE (2019), \tr{pp. 35 e 40--44}{pp. 35 and 40--44}.}
\end{frame}

\begin{frame}{\tr{Sua vez}{Your turn}}
  \begin{columns}[T]
    \begin{column}{0.50\textwidth}
      \begin{regra}[title={\tr{Fechamento e voltas}{Misclosure and turns}}]
        \small
        \begin{enumerate}
          \item[(a)] \tr{Um circuito de $K=\num{100}$\,km tem fechamento corrigido
                $w=\num{21}$\,mm. Passa no controle de \num{0.5}\,mm/km?
                Quanto vale $|w|/K$?}{A \num{100}\,km circuit has corrected misclosure
                $w=\num{21}$\,mm. Does it pass the \num{0.5}\,mm/km check?
                What is $|w|/K$?}
          \item[(b)] \tr{Compare dois comprimentos hipotéticos: \num{230000} e
                \num{370000}\,km. Quantas voltas à Terra (\num{40075}\,km) a mais?}{Compare
                two hypothetical lengths: \num{230000} and \num{370000}\,km. How many more
                turns around the Earth (\num{40075}\,km)?}
        \end{enumerate}
      \end{regra}
    \end{column}
    \begin{column}{0.46\textwidth}
      \onslide<2->{
      \begin{ideia}[title={\tr{Respostas}{Answers}}]
        \small\raggedright
        (a) $\num{0.5}\times\num{100}=\num{50}$\,mm:
        \tr{\bnum{passa}}{\bnum{passes}} ($21<50$);
        $|w|/K=\num{21}/\num{100}=\bnum{\num{0.21}}$\,mm/km.
        \medskip

        (b) $\dfrac{\num{370000}-\num{230000}}{\num{40075}}\approx\bnum{\num{3.5}}$
        \tr{voltas.}{turns.}
        \medskip

        \tr{Cumprir o limite não prova ausência de viés sistemático.}{Meeting the limit
        does not prove that systematic bias is absent.}
      \end{ideia}
      }
    \end{column}
  \end{columns}
\end{frame}

% ------------------------------------------------------------------
\begin{frame}{\tr{As superfícies de nível não são paralelas}{Level surfaces are not parallel}}
  \begin{columns}[T]
    \begin{column}{0.56\textwidth}
      \centering
      \begin{tikzpicture}[x=0.78cm,y=1.05cm,font=\scriptsize]
        \draw[corMar, very thick] (0,0.3) -- (8,0.3);
        \draw[corMar, very thick] plot[smooth] coordinates {(0,2.35) (4,2.15) (8,1.95)};
        \node[corMar, right] at (8.05,0.3) {$W_1$};
        \node[corMar, right] at (8.05,1.95) {$W_2$};
        \fill[corFrame] (1,0.3) circle (0.09) node[below, text=corFrame] {$A$};
        \fill[corFrame] (7,2.0) circle (0.09) node[above, text=corFrame] {$B$};
        % caminho 1: sobe à esquerda, anda por W2
        \draw[corTempo, very thick, -{Stealth[length=2mm]}] (1,0.3) -- (1,2.31);
        \draw[corTempo, very thick, -{Stealth[length=2mm]}] (1,2.31) -- (6.9,2.0);
        \node[corTempo, left, align=right, font=\tiny] at (0.85,1.3) {$\Sigma\Delta n$\\$=\num{2.0}$};
        % caminho 2: anda por W1, sobe à direita
        \draw[corErro, very thick, -{Stealth[length=2mm]}] (1.1,0.3) -- (7,0.3);
        \draw[corErro, very thick, -{Stealth[length=2mm]}] (7,0.3) -- (7,1.95);
        \node[corErro, right, align=left, font=\tiny] at (7.1,1.1) {$\Sigma\Delta n$\\$=\num{1.7}$};
      \end{tikzpicture}\\[1pt]
      {\scriptsize\tr{Exagerado. Ao longo de uma superfície de nível o desnível é zero.}{Exaggerated. Along a level surface the height difference is zero.}}
    \end{column}
    \begin{column}{0.40\textwidth}
      \small
      \tr{O nível mede o desnível entre superfícies de nível \emph{vizinhas}. Mas a distância
      entre $W_1$ e $W_2$ \bnum{muda de lugar para lugar}: onde a gravidade é maior, ela é
      menor.}{The level measures the height difference between \emph{neighboring} level
      surfaces. But the distance between $W_1$ and $W_2$ \bnum{changes from place to place}:
      where gravity is larger, it is smaller.}
      \medskip
      \begin{regra}[title={\tr{O que é o mesmo nos dois caminhos}{What is the same on both paths}}]
        \centering\small
        $\Delta W=g\,\Delta n$
      \end{regra}
      \smallskip

      \tr{$\Delta W$ (a diferença de \emph{geopotencial}) é a mesma. $\Sigma\Delta n$ não.}{$\Delta W$ (the
      \emph{geopotential} difference) is the same. $\Sigma\Delta n$ is not.}
    \end{column}
  \end{columns}
\end{frame}

\begin{frame}{\tr{Um circuito perfeito não fecha em zero}{A perfect loop does not close at
zero}}
  \begin{columns}[T]
    \begin{column}{0.50\textwidth}
      \footnotesize
      \tr{Imagine um nivelador \emph{perfeito}, sem erro nenhum, que sobe uma serra por um
      caminho e desce por outro. O desnível que o nível ``vê'' em cada trecho é
      $\Delta n=\Delta W/g$. Se $g$ é um pouco diferente nos dois caminhos, o circuito
      \bnum{não fecha}:}{Imagine a \emph{perfect} leveler, with no error at all, who climbs a
      ridge by one path and comes down by another. The height difference the level ``sees'' in
      each stretch is $\Delta n=\Delta W/g$. If $g$ is slightly different on the two paths, the
      loop \bnum{does not close}:}
      \[
        w\approx H\,\frac{\Delta g}{\gamma}
      \]
      \tr{$H$ é o desnível subido e $\Delta g$ a diferença média de gravidade entre os
      caminhos.}{$H$ is the climbed height and $\Delta g$ the mean gravity difference between
      the paths.}
    \end{column}
    \begin{column}{0.46\textwidth}
      \renewcommand{\arraystretch}{1.3}
      \small
      \begin{tabular}{@{}rrr@{}}
        \toprule
        $H$ (m) & $\Delta g$ (mGal) & $w$ (mm) \\
        \midrule
        \num{500} & \num{20} & \num{10.2} \\
        \num{1000} & \num{20} & \bnum{\num{20.4}} \\
        \num{1000} & \num{50} & \errr{\num{51.1}} \\
        \bottomrule
      \end{tabular}
      \medskip

      \begin{cuidado}[title={\tr{Correções e controle}{Corrections and checks}}]
        \scriptsize\tr{O efeito de \num{20.4}\,mm antecede correções. Num circuito de
        \num{100}\,km, o limite após correções é \num{50}\,mm ($\num{0.5}$\,mm/km).}{
        The \num{20.4}\,mm effect precedes corrections. For \num{100}\,km, the
        corrected-circuit limit is \num{50}\,mm ($\num{0.5}$\,mm/km).}
        \fonteslide{IBGE (1983), pp. 28--30, VI.3.3.}
        \smallskip

        \tr{(1\,mGal $=10^{-5}$\,m/s$^{2}$)}{(1\,mGal $=10^{-5}$\,m/s$^{2}$)}
      \end{cuidado}
    \end{column}
  \end{columns}
\end{frame}

\begin{frame}{\tr{O número geopotencial resolve}{The geopotential number solves it}}
  \begin{columns}[T]
    \begin{column}{0.54\textwidth}
      \footnotesize
      \tr{A saída (já vista nas aulas do geoide): \bnum{multiplicar cada desnível pela gravidade
      medida} no trecho:}{The way out (already seen in the geoid classes): \bnum{multiply each
      height difference by the gravity measured} on the stretch:}
      \[
        \Delta C=g\,\Delta n,\qquad C=\sum g\,\Delta n
      \]
      \begin{itemize}
        \setlength{\itemsep}{3pt}
        \item \tr{$C$ \bnum{não depende do caminho}: todo circuito fecha em zero (a menos do
              erro de medida);}{$C$ \bnum{does not depend on the path}: every loop closes at
              zero (up to measurement error);}
        \item \tr{unidade: m$^{2}$/s$^{2}$ (cerca de \num{9.8} vezes o desnível em
              metros);}{unit: m$^{2}$/s$^{2}$ (about \num{9.8} times the height difference in
              meters);}
        \item \tr{a altitude sai de $C$: $H^{*}=C/\bar\gamma$ (normal).}{the altitude comes
              from $C$: $H^{*}=C/\bar\gamma$ (normal).}
      \end{itemize}
    \end{column}
    \begin{column}{0.42\textwidth}
      \begin{ideia}[title={\tr{O preço}{The price}}]
        \footnotesize\tr{Para somar $g\,\Delta n$ é preciso conhecer $g$ \emph{em cada RN} ---
        gravimetria na rede de nivelamento.}{To add up $g\,\Delta n$ you must know $g$
        \emph{at each benchmark} --- gravimetry on the leveling network.}
      \end{ideia}
      \medskip
      \begin{cuidado}[title={\tr{E o Brasil?}{And Brazil?}}]
        \footnotesize\tr{Até \bnum{2011}, sem gravidade medida nas RRNN, só se aplicava a correção com
        \emph{gravidade teórica}. Só o \bnum{REALT-2018} trouxe a gravidade \emph{real}: é a
        Aula 3.}{Until \bnum{2011}, without gravity measured at the benchmarks, only the
        correction with \emph{theoretical gravity} was applied. Only \bnum{REALT-2018} brought
        \emph{real} gravity: that is Class 3.}
      \end{cuidado}
    \end{column}
  \end{columns}
\end{frame}

\begin{frame}{\tr{Pare e pense}{Stop and think}}
  \begin{columns}[T]
    \begin{column}{0.46\textwidth}
      \begin{armadilha}[title={\tr{O enigma}{The riddle}}]
        \small\tr{``O nivelador mediu \emph{perfeitamente}. Mesmo assim, o circuito de \num{100}\,km
        em terreno de serra fechou em \num{20}\,mm. Quem errou?''}{``The leveler measured
        \emph{perfectly}. Even so, the \num{100}\,km loop in rugged terrain closed at
        \num{20}\,mm. Who made the mistake?''}
      \end{armadilha}
    \end{column}
    \begin{column}{0.50\textwidth}
      \onslide<2->{
      \begin{ideia}[title={\tr{Ninguém}{Nobody}}]
        \footnotesize
        \begin{itemize}
          \setlength{\itemsep}{2pt}
          \item \tr{O erro está na \emph{conta}, não na medida: $\sum\Delta n$ não é uma
                grandeza física que fecha em circuito.}{The error is in the \emph{computation},
                not in the measurement: $\sum\Delta n$ is not a physical quantity that closes
                around a loop.}
          \item \tr{A grandeza que fecha é $\sum g\,\Delta n$ (número geopotencial).}{The
                quantity that closes is $\sum g\,\Delta n$ (geopotential number).}
          \item \tr{Em terreno plano e curto o efeito some, e por isso por décadas se pôde
                ignorá-lo. Num país de \num{8.5} milhões de km$^{2}$, não.}{On flat, short
                terrain the effect vanishes, which is why for decades it could be ignored. In a
                country of \num{8.5} million km$^{2}$, it cannot.}
        \end{itemize}
      \end{ideia}
      }
    \end{column}
  \end{columns}
\end{frame}

\begin{frame}{\tr{As quatro perguntas, respondidas}{The four questions, answered}}
  \small
  \begin{enumerate}
    \setlength{\itemsep}{5pt}
    \item \textbf{\tr{Como se leva o zero do mar até um ponto distante?}{How is the zero
          carried from the sea to a distant point?}} \tr{Por \bnum{nivelamento geométrico}:
          lance a lance, com nível e mira, $\Delta h=\text{ré}-\text{vante}$, somando os
          desníveis ao longo de uma linha de RRNN.}{By \bnum{geometric leveling}: setup by
          setup, with level and rod, $\Delta h=\text{back}-\text{fore}$, adding the height
          differences along a line of benchmarks.}
    \item \textbf{\tr{De onde vem o erro e por que cresce com $\sqrt{K}$?}{Where does the error
          come from and why does it grow with $\sqrt{K}$?}} \tr{Cada lance erra um pouco.
          Se os erros são ao acaso, o total cresce com $\eta\sqrt{K}$; se são sistemáticos,
          com $\delta K$ --- e esse vence nas distâncias longas.}{Each setup is a little off.
          If errors are random, the total grows as $\eta\sqrt{K}$; if systematic, as
          $\delta K$ --- and that one wins over long distances.}
    \item \textbf{\tr{Como é a rede de RRNN?}{What does the benchmark network look like?}}
          \tr{Linhas ao longo das estradas, formando circuitos e macrocircuitos.
          O bloco principal do REALT-2018 tem \bnum{\num{66882}} RRNN; há outros conjuntos
          tratados separadamente.}{Lines along roads, forming loops and macro-loops.
          The REALT-2018 principal block has \bnum{\num{66882}} benchmarks; other sets
          are treated separately.}
    \item \textbf{\tr{Por que um circuito sem erro não fecha?}{Why does an error-free loop not
          close?}} \tr{Porque o nível mede $\Delta n$ e as superfícies de nível não são
          paralelas. O que fecha é o número geopotencial, $\sum g\,\Delta n$.}{Because the level
          measures $\Delta n$ and level surfaces are not parallel. What closes is the
          geopotential number, $\sum g\,\Delta n$.}
  \end{enumerate}
\end{frame}

\begin{frame}{\tr{Para levar para casa}{Take home}}
  \begin{columns}[T]
    \begin{column}{0.54\textwidth}
      \begin{regra}[title={\tr{Quatro ideias}{Four ideas}}]
        \small
        \begin{enumerate}
          \setlength{\itemsep}{3pt}
          \item \tr{O zero chega por \bnum{nivelamento}: $\Delta h=$ ré $-$ vante, lance a
                lance.}{The zero arrives by \bnum{leveling}: $\Delta h=$ back $-$ fore, setup by
                setup.}
          \item \tr{Erro aleatório: $\eta\sqrt{K}$. Sistemático: $\delta K$. Tolerância de seção:
                $3\sqrt{K}$ mm (IBGE, 1983, p. 30).}{Random error: $\eta\sqrt{K}$. Systematic:
                $\delta K$. Section tolerance: $3\sqrt{K}$ mm (IBGE, 1983, p. 30).}
          \item \tr{A rede é de linhas e circuitos; fechar bem é \bnum{necessário, não
                suficiente}.}{The network is made of lines and loops; closing well is
                \bnum{necessary, not sufficient}.}
          \item \tr{$\sum\Delta n$ não fecha; $\sum g\,\Delta n$ fecha.}{$\sum\Delta n$ does not
                close; $\sum g\,\Delta n$ does.}
        \end{enumerate}
      \end{regra}
    \end{column}
    \begin{column}{0.42\textwidth}
      \begin{ideia}[title={\tr{Ponte para a Aula 3}{Bridge to Class 3}}]
        \small\tr{Cada seção foi nivelada, mas a rede tem \emph{milhares} de circuitos que não
        fecham exatamente. Quem decide a altitude de cada RN? \bnum{O ajustamento.}}{Each
        section was leveled, but the network has \emph{thousands} of loops that do not close
        exactly. Who decides the altitude of each benchmark? \bnum{The adjustment.}}
        \medskip

        \tr{E por que \emph{a mesma RN} já teve várias altitudes ao longo da história?}{And why
        has \emph{the same benchmark} had several altitudes throughout history?}
      \end{ideia}
    \end{column}
  \end{columns}
\end{frame}

% ==================================================================
\section{\tr{Aula 3: Os ajustamentos e o REALT}{Class 3: The adjustments and REALT}}

\begin{frame}{\tr{As perguntas de hoje}{Today's questions}}
  \begin{columns}[T]
    \begin{column}{0.58\textwidth}
      \begin{enumerate}
        \setlength{\itemsep}{8pt}
        \item \tr{Por que a \emph{mesma} RN já teve várias altitudes?}{Why has the \emph{same}
              benchmark had several altitudes?}
        \item \tr{O que é ``ajustar'' uma rede de nivelamento --- e como se faz?}{What does
              ``adjusting'' a leveling network mean --- and how is it done?}
        \item \tr{Como foram os ajustamentos do Brasil, de 1945 a 2011?}{How were Brazil's
              adjustments, from 1945 to 2011?}
        \item \tr{O que o REALT-2018 mudou por dentro?}{What did REALT-2018 change on the
              inside?}
      \end{enumerate}
    \end{column}
    \begin{column}{0.38\textwidth}
      \begin{ideia}[title={\tr{Combinado}{Deal}}]
        \tr{No fim da aula, você vai saber responder às quatro --- com as suas
        palavras.}{By the end of the class you will be able to answer all four --- in your own
        words.}
        \medskip

        \tr{Elas voltam no último slide.}{They come back on the last slide.}
      \end{ideia}
    \end{column}
  \end{columns}
\end{frame}

\begin{frame}{\tr{O que a aula anterior deixou}{What the previous class left us}}
  \begin{columns}[T]
    \begin{column}{0.54\textwidth}
      \small
      \tr{Na Aula 2, quatro ideias:}{In Class 2, four ideas:}
      \begin{enumerate}
        \setlength{\itemsep}{3pt}
        \item \tr{o zero chega por \bnum{nivelamento}: $\Delta h=$ ré $-$ vante;}{the zero
              arrives by \bnum{leveling}: $\Delta h=$ back $-$ fore;}
        \item \tr{cada seção tem erro: $\eta\sqrt{K}$ (aleatório) e $\delta K$
              (sistemático);}{each section has error: $\eta\sqrt{K}$ (random) and $\delta K$
              (systematic);}
        \item \tr{a rede tem \bnum{circuitos} que não fecham exatamente;}{the network has
              \bnum{loops} that do not close exactly;}
        \item \tr{$\sum\Delta n$ não fecha; $\sum g\,\Delta n$ fecha.}{$\sum\Delta n$ does not
              close; $\sum g\,\Delta n$ does.}
      \end{enumerate}
    \end{column}
    \begin{column}{0.42\textwidth}
      \begin{ideia}[title={\tr{A pergunta de hoje}{Today's question}}]
        \small\tr{Cada circuito fecha com um \emph{pequeno desacordo}. Num país, são milhares de
        circuitos que se cruzam. Que altitude damos a cada RN? \bnum{Uma só}, que concorde com
        todas as medidas ao mesmo tempo.}{Each loop closes with a \emph{small disagreement}. In
        a country there are thousands of loops crossing each other. Which altitude do we give
        each benchmark? \bnum{One only}, which agrees with all measurements at once.}
      \end{ideia}
    \end{column}
  \end{columns}
\end{frame}

\begin{frame}{\tr{A mesma RN, várias altitudes}{The same benchmark, several altitudes}}
  \centering
  \begin{tikzpicture}[x=1cm,y=1cm,font=\scriptsize]
    \pic[scale=2.2] at (0.2,0.7) {marco};
    \node[corFrame, below] at (0.2,0.65) {RN};
    \draw[->, thick] (1.0,0.9) -- (12.6,0.9);
    \foreach \a/\x/\y/\i in {1952/1.8/1.50/1, 1962/3.9/1.30/2, 1975/5.8/1.62/3, 1992/7.7/1.40/4, 2011/9.6/1.55/5, 2018/11.5/1.25/6} {
      \draw (\x,0.8) -- (\x,1.0) node[below=2pt, pos=0] {\a};
      \node[caixa, draw=corTempo, fill=corTempo!10, minimum width=0.9cm] at (\x,\y+0.45) {$H_{\i}$};
    }
    \draw[corTempo, thick] (1.8,1.95) -- (3.9,1.75) -- (5.8,2.07) -- (7.7,1.85) -- (9.6,2.0) -- (11.5,1.7);
    \node[corTempo, above] at (6.6,2.5) {\tr{a altitude de \emph{uma} RN, em seis ajustamentos (alturas ilustrativas)}{the altitude of \emph{one} benchmark, in six adjustments (illustrative heights)}};
  \end{tikzpicture}
  \bigskip

  \begin{columns}[T]
    \begin{column}{0.30\textwidth}
      \begin{ideia}[title={\tr{Dados novos}{New data}}]
        \footnotesize\tr{Mais linhas, remedições, erros achados e corrigidos.}{More lines,
        re-leveling, errors found and corrected.}
      \end{ideia}
    \end{column}
    \begin{column}{0.30\textwidth}
      \begin{regra}[title={\tr{Método novo}{New method}}]
        \footnotesize\tr{Blocos $\to$ ajuste simultâneo; sem gravidade $\to$ com gravidade.}{Blocks
        $\to$ simultaneous adjustment; no gravity $\to$ with gravity.}
      \end{regra}
    \end{column}
    \begin{column}{0.30\textwidth}
      \begin{cuidado}[title={\tr{Zero novo}{New zero}}]
        \footnotesize\tr{Torres $\to$ Imbituba; Santana para o Amapá.}{Torres $\to$ Imbituba;
        Santana for Amapá.}
      \end{cuidado}
    \end{column}
  \end{columns}
\end{frame}

\begin{frame}{\tr{Por que ajustar?}{Why adjust?}}
  \begin{columns}[T]
    \begin{column}{0.46\textwidth}
      \centering
      \begin{tikzpicture}[x=0.82cm,y=0.82cm,font=\scriptsize]
        \coordinate (A) at (0,0);
        \coordinate (B) at (4.0,0.2);
        \coordinate (C) at (4.2,3.2);
        \coordinate (D) at (0.1,3.0);
        \draw[corFrame, thick] (A) -- (B) -- (C) -- (D) -- cycle;
        \draw[corFrame, thick] (A) -- (C);
        \draw[corFrame, thick] (B) -- (D);
        \foreach \p in {A,B,C,D} { \fill[corFrame] (\p) circle (0.1); }
        \node[below left] at (A) {$A$};
        \node[below right] at (B) {$B$};
        \node[above right] at (C) {$C$};
        \node[above left] at (D) {$D$};
        \node[corErro, font=\small] at (2.05,1.7) {$\neq0$};
        \node[corErro, below] at (2.0,-0.55) {\tr{cada circuito fecha com um desacordo}{each loop closes with a disagreement}};
      \end{tikzpicture}
    \end{column}
    \begin{column}{0.50\textwidth}
      \small
      \tr{Seis seções medidas, \bnum{três} altitudes desconhecidas ($B$, $C$, $D$; $A$ é o
      datum). Há \emph{mais medidas do que o necessário}: três a mais.}{Six sections measured,
      \bnum{three} unknown altitudes ($B$, $C$, $D$; $A$ is the datum). There are \emph{more
      measurements than needed}: three extra.}
      \medskip

      \tr{Essas medidas \emph{extras} (redundância) não concordam exatamente entre si:
      cada circuito fecha com um pequeno desacordo.}{Those \emph{extra} measurements
      (redundancy) do not agree exactly with each other: each loop closes with a small
      disagreement.}
      \medskip
      \begin{regra}[title={\tr{Ajustar é...}{Adjusting is...}}]
        \small\tr{achar \bnum{uma só} altitude para cada RN, repartindo o desacordo de forma
        justa --- e dizer \bnum{quão incerta} cada altitude ficou.}{finding \bnum{one single}
        altitude for each benchmark, spreading the disagreement fairly --- and saying \bnum{how
        uncertain} each altitude is.}
      \end{regra}
    \end{column}
  \end{columns}
\end{frame}

\begin{frame}{\tr{O modelo: cada medida é uma equação}{The model: each measurement is an equation}}
  \begin{columns}[T]
    \begin{column}{0.52\textwidth}
      \small
      \tr{Para a seção de $i$ a $j$, o desnível medido $\Delta h_{ij}$ deveria ser $H_j-H_i$.
      Como há erro, soma-se uma \bnum{correção} $v_{ij}$:}{For the section from $i$ to $j$,
      the measured height difference $\Delta h_{ij}$ should be $H_j-H_i$. Since there is error,
      a \bnum{correction} $v_{ij}$ is added:}
      \[
        \boxed{\;\Delta h_{ij}+v_{ij}=H_j-H_i\;}
      \]
      \tr{Cada seção recebe um \bnum{peso} $p_{ij}=\dfrac{1}{\eta^{2}K_{ij}}$: seção longa $=$
      medida menos confiável $=$ peso menor.}{Each section gets a \bnum{weight}
      $p_{ij}=\dfrac{1}{\eta^{2}K_{ij}}$: long section $=$ less reliable measurement $=$
      smaller weight.}
    \end{column}
    \begin{column}{0.44\textwidth}
      \begin{regra}[title={\tr{Mínimos quadrados}{Least squares}}]
        \small\tr{Escolhem-se as altitudes que \bnum{minimizam}}{The altitudes are chosen that
        \bnum{minimize}}
        \[
          \sum p_{ij}\,v_{ij}^{2}
        \]
        \tr{As correções grandes custam caro (ao quadrado) e as medidas pouco confiáveis custam
        barato (peso pequeno).}{Large corrections are expensive (squared) and unreliable
        measurements are cheap (small weight).}
      \end{regra}
      \medskip
      {\footnotesize\tr{É o mesmo método do simulador \texttt{nivelamento/}.}{It is the same
      method as in the \texttt{nivelamento/} simulator.}}
    \end{column}
  \end{columns}
\end{frame}

\begin{frame}{\tr{O datum entra como um ponto fixo}{The datum enters as a fixed point}}
  \begin{columns}[T]
    \begin{column}{0.52\textwidth}
      \centering
      \begin{tikzpicture}[x=0.9cm,y=0.9cm,font=\scriptsize]
        % rede livre
        \begin{scope}
          \foreach \x/\y in {0.4/0.6, 1.8/0.3, 2.4/1.5, 1.0/1.7} { \fill[corFrame] (\x,\y) circle (0.09); }
          \draw[rede] (0.4,0.6) -- (1.8,0.3) -- (2.4,1.5) -- (1.0,1.7) -- cycle (0.4,0.6) -- (2.4,1.5);
          \draw[<->, corTempo, thick] (3.1,0.3) -- node[right, text=corTempo] {$+c$} (3.1,1.8);
          \node[below] at (1.5,0) {\tr{rede \textbf{livre}}{\textbf{free} network}};
        \end{scope}
        % rede injuncionada
        \begin{scope}[shift={(4.6,0)}]
          \foreach \x/\y in {1.8/0.3, 2.4/1.5, 1.0/1.7} { \fill[corFrame] (\x,\y) circle (0.09); }
          \draw[rede] (0.4,0.6) -- (1.8,0.3) -- (2.4,1.5) -- (1.0,1.7) -- cycle (0.4,0.6) -- (2.4,1.5);
          \pic[scale=1.8] at (0.4,0.45) {marco};
          \draw[corMar, dashed] (-0.2,0.45) -- (1.0,0.45);
          \node[corMar, left, font=\tiny] at (-0.2,0.45) {\tr{datum}{datum}};
          \node[below] at (1.5,0) {\tr{rede \textbf{injuncionada}}{\textbf{constrained} network}};
        \end{scope}
      \end{tikzpicture}
    \end{column}
    \begin{column}{0.44\textwidth}
      \small
      \tr{Todas as medidas são \emph{diferenças}. Somar o mesmo número $c$ a todas as
      altitudes \bnum{não muda nenhuma medida}: a rede ``flutua''.}{All measurements are
      \emph{differences}. Adding the same number $c$ to all altitudes \bnum{changes no
      measurement}: the network ``floats''.}
      \medskip

      \tr{Para fixar, prende-se \bnum{um ponto} ao datum: a RN de Imbituba (4X) ou a de Santana
      (9329T) recebe altitude dada. Aí a solução é única.}{To fix it, \bnum{one point} is pinned
      to the datum: the Imbituba benchmark (4X) or the Santana one (9329T) is given an altitude.
      Then the solution is unique.}
      \medskip

      {\footnotesize\tr{É o mesmo ``defeito de posto'' do datum 3D da Parte 1.}{It is the same
      ``rank defect'' as the 3D datum in Part 1.}}
    \end{column}
  \end{columns}
\end{frame}

\begin{frame}{\tr{Um exemplo: a rede de 4 nós}{An example: the 4-node network}}
  \begin{columns}[T]
    \begin{column}{0.52\textwidth}
      \centering
      \begin{tikzpicture}[x=0.85cm,y=0.85cm,font=\scriptsize]
        \coordinate (A) at (0,0);
        \coordinate (B) at (4.6,0.2);
        \coordinate (C) at (4.8,3.6);
        \coordinate (D) at (0.1,3.4);
        \draw[seta, corFrame] (A) -- node[below, text=corTempo, align=center] {$+\num{2.350}$} (B);
        \draw[seta, corFrame] (B) -- node[right, text=corTempo] {$-\num{2.483}$} (C);
        \draw[seta, corFrame] (C) -- node[above, text=corTempo] {$+\num{1.254}$} (D);
        \draw[seta, corFrame] (D) -- node[left, text=corTempo] {$-\num{1.121}$} (A);
        \draw[seta, corFrame] (A) -- node[pos=0.28, above, sloped, text=corTempo] {$-\num{0.133}$} (C);
        \draw[seta, corFrame] (B) -- node[pos=0.28, above, sloped, text=corTempo] {$-\num{1.223}$} (D);
        \fill[corMar] (A) circle (0.12);
        \foreach \p in {B,C,D} { \fill[corFrame] (\p) circle (0.1); }
        \node[below left] at (A) {$A$};
        \node[below right] at (B) {$B$};
        \node[above right] at (C) {$C$};
        \node[above left] at (D) {$D$};
      \end{tikzpicture}
    \end{column}
    \begin{column}{0.44\textwidth}
      \footnotesize
      \renewcommand{\arraystretch}{1.25}
      \begin{tabular}{@{}lrr@{}}
        \toprule
        \tr{seção}{section} & $K$ (km) & $\sigma$ (mm) \\
        \midrule
        $A\to B$ & \num{4} & \num{2.00} \\
        $B\to C$ & \num{6} & \num{2.45} \\
        $C\to D$ & \num{5} & \num{2.24} \\
        $D\to A$ & \num{3} & \num{1.73} \\
        $A\to C$ & \num{5} & \num{2.24} \\
        $B\to D$ & \num{7} & \num{2.65} \\
        \bottomrule
      \end{tabular}
      \medskip

      \tr{$A$ é o datum: $H_A=\num{10.000}$\,m. Desníveis em metros; $\eta=\num{1}$\,mm$/\sqrt{\text{km}}$ ($\sigma=\eta\sqrt{K}$).}{$A$ is
      the datum: $H_A=\num{10.000}$\,m. Height differences in meters; $\eta=\num{1}$\,mm$/\sqrt{\text{km}}$
      ($\sigma=\eta\sqrt{K}$).}
      \smallskip

      \tr{O circuito $A\!-\!B\!-\!D\!-\!A$ fecha em \bnum{$+\num{6}$\,mm}; os outros dois fecham
      em zero.}{The loop $A\!-\!B\!-\!D\!-\!A$ closes at \bnum{$+\num{6}$\,mm}; the other two
      close at zero.}
    \end{column}
  \end{columns}
\end{frame}

\begin{frame}{\tr{As contas do exemplo}{The computations of the example}}
  \begin{columns}[T]
    \begin{column}{0.58\textwidth}
      \footnotesize
      \tr{Monta-se $N\,\hat x=b$ para as três incógnitas ($B$, $C$, $D$), com $p=1/K$:}{Build
      $N\,\hat x=b$ for the three unknowns ($B$, $C$, $D$), with $p=1/K$:}
      \[
        N=\begin{pmatrix}
          \num{0.5595} & -\num{0.1667} & -\num{0.1429}\\
          -\num{0.1667} & \num{0.5667} & -\num{0.2000}\\
          -\num{0.1429} & -\num{0.2000} & \num{0.6762}
        \end{pmatrix}\,\mathrm{km}^{-1}
      \]
      \tr{Cada seção $ij$ põe $+p$ na diagonal de $i$ e de $j$ e $-p$ fora dela: $N$ é o
      ``mapa'' da rede.}{Each section $ij$ puts $+p$ on the diagonal of $i$ and of $j$ and
      $-p$ off it: $N$ is the ``map'' of the network.}
      \medskip

      \tr{Resolvendo:}{Solving:}
      \[
        \hat H_B=\bnum{\num{12.3487}}\,\mathrm{m},\quad
        \hat H_C=\bnum{\num{9.8670}}\,\mathrm{m},\quad
        \hat H_D=\bnum{\num{11.1220}}\,\mathrm{m}
      \]
    \end{column}
    \begin{column}{0.38\textwidth}
      \begin{ideia}[title={\tr{Nada de mágica}{No magic}}]
        \footnotesize\tr{É só um sistema de $3$ equações com $3$ incógnitas. Na rede real do
        Brasil são \bnum{dezenas de milhares}: o método é o mesmo, o computador é que
        muda.}{It is only a system of $3$ equations with $3$ unknowns. In the real Brazilian
        network there are \bnum{tens of thousands}: the method is the same, the computer is
        what changes.}
      \end{ideia}
      \medskip
      {\scriptsize\tr{Conta completa e conferência: \texttt{sgb\_altimetria\_numeros.py} e
      apêndice A2.}{Full computation and check: \texttt{sgb\_altimetria\_numeros.py} and
      appendix A2.}}
    \end{column}
  \end{columns}
\end{frame}

\begin{frame}{\tr{O que o ajuste fez}{What the adjustment did}}
  \begin{columns}[T]
    \begin{column}{0.52\textwidth}
      \centering
      \begin{tikzpicture}[x=1cm,y=1cm,font=\scriptsize]
        \draw[thick] (3.6,0.3) -- (3.6,-4.6);
        \foreach \i/\n/\v in {0/{$A\to B$}/-1.29, 1/{$B\to C$}/1.26, 2/{$C\to D$}/1.02,
                             3/{$D\to A$}/-0.99, 4/{$A\to C$}/-0.03, 5/{$B\to D$}/-3.72} {
          \pgfmathsetmacro{\w}{0.6*\v}
          \ifdim\w pt<0pt
            \fill[corErro!75] ({3.6+\w},{-0.75*\i}) rectangle (3.6,{-0.75*\i-0.5});
            \node[left] at ({3.6+\w-0.05},{-0.75*\i-0.25}) {$\num{\v}$};
          \else
            \fill[corFrame!75] (3.6,{-0.75*\i}) rectangle ({3.6+\w},{-0.75*\i-0.5});
            \node[right] at ({3.6+\w+0.05},{-0.75*\i-0.25}) {$+\num{\v}$};
          \fi
          \node[left] at (0.5,{-0.75*\i-0.25}) {\n};
        }
        \node[below] at (3.6,-4.65) {\tr{correção $v$ (mm)}{correction $v$ (mm)}};
      \end{tikzpicture}
    \end{column}
    \begin{column}{0.44\textwidth}
      \footnotesize
      \tr{O desacordo de \num{6}\,mm em $A\!-\!B\!-\!D\!-\!A$ foi \bnum{repartido}: a maior
      correção caiu em \bnum{$B\to D$}, a seção mais longa ($K=\num{7}$\,km, menor peso).}{The
      \num{6}\,mm disagreement in $A\!-\!B\!-\!D\!-\!A$ was \bnum{spread}: the largest
      correction fell on \bnum{$B\to D$}, the longest section ($K=\num{7}$\,km, smallest
      weight).}
      \smallskip

      \tr{O \bnum{fator de variância} $\hat\sigma_0=\num{1.03}$ é quase $1$: os dados concordam
      com $\eta=\num{1}$. As altitudes saem com $\sigma\approx\num{1.4}$ a $\num{1.6}$\,mm.}{The
      \bnum{variance factor} $\hat\sigma_0=\num{1.03}$ is almost $1$: the data agree with
      $\eta=\num{1}$. The altitudes come out with $\sigma\approx\num{1.4}$ to
      $\num{1.6}$\,mm.}
      \smallskip

      \tr{Conferência: o exemplo foi \emph{inventado} a partir de altitudes conhecidas
      ($\num{12.348}$; $\num{9.868}$; $\num{11.121}$\,m). O ajuste ficou a cerca de \bnum{1\,mm}
      de cada uma.}{Check: the example was \emph{made up} from known altitudes
      ($\num{12.348}$; $\num{9.868}$; $\num{11.121}$\,m). The adjustment came within about
      \bnum{1\,mm} of each.}
    \end{column}
  \end{columns}
\end{frame}

\begin{frame}{\tr{Sua vez}{Your turn}}
  \begin{columns}[T]
    \begin{column}{0.50\textwidth}
      \begin{regra}[title={\tr{Ajustar um circuito}{Adjusting one loop}}]
        \small\tr{Um circuito de três seções, com $K=\num{2}$, $\num{3}$ e $\num{5}$\,km, fechou
        em $w=+\num{5}$\,mm. Os pesos são $p=1/K$. Quanto se corrige em cada seção, de modo
        que o circuito feche em zero?}{A loop of three sections, with $K=\num{2}$, $\num{3}$ and
        $\num{5}$\,km, closed at $w=+\num{5}$\,mm. The weights are $p=1/K$. How much is each
        section corrected so that the loop closes at zero?}
        \medskip

        \tr{Dica: a correção é proporcional a $K$.}{Hint: the correction is proportional to
        $K$.}
      \end{regra}
    \end{column}
    \begin{column}{0.46\textwidth}
      \onslide<2->{
      \begin{ideia}[title={\tr{Resposta}{Answer}}]
        \small\raggedright
        $v_i=-w\,\dfrac{K_i}{\sum K}$ \tr{com}{with} $\sum K=\num{10}$:
        \medskip

        $v_1=-\num{1.0}$, \ $v_2=-\num{1.5}$, \ $v_3=-\num{2.5}$\,mm.
        \medskip

        \tr{Soma: $-\num{5}$\,mm $=-w$: o circuito fecha. Quem tem mais quilômetros
        \bnum{absorve mais}.}{Sum: $-\num{5}$\,mm $=-w$: the loop closes. Whoever has more
        kilometers \bnum{absorbs more}.}
      \end{ideia}
      }
    \end{column}
  \end{columns}
\end{frame}

\begin{frame}{\tr{Pare e pense}{Stop and think}}
  \begin{columns}[T]
    \begin{column}{0.46\textwidth}
      \begin{armadilha}[title={\tr{Uma ideia tentadora}{A tempting idea}}]
        \small\tr{``Depois do ajustamento, os circuitos fecham em zero e as altitudes têm
        incerteza de poucos milímetros. Então os \emph{erros acabaram}.''}{``After the
        adjustment, the loops close at zero and the altitudes have an uncertainty of a few
        millimeters. So the \emph{errors are gone}.''}
        \medskip

        \tr{O que está errado?}{What is wrong?}
      \end{armadilha}
    \end{column}
    \begin{column}{0.50\textwidth}
      \onslide<2->{
      \begin{ideia}[title={\tr{Resposta}{Answer}}]
        \footnotesize
        \begin{itemize}
          \setlength{\itemsep}{2pt}
          \item \tr{O ajuste \bnum{reparte} o desacordo; ele não o faz desaparecer.}{The
                adjustment \bnum{spreads} the disagreement; it does not make it disappear.}
          \item \tr{Ele reduz o erro \emph{aleatório}; o \emph{sistemático} que fecha bem em
                circuito continua lá.}{It reduces the \emph{random} error; the \emph{systematic}
                one that closes well in a loop is still there.}
          \item \tr{Um erro \emph{grosseiro} pode ser espalhado pelas vizinhas e
                escondido --- por isso se faz o teste dos resíduos (\emph{data snooping}), como no
                simulador.}{A \emph{blunder} can be spread over the neighbors and hidden --- which
                is why the residual test (\emph{data snooping}) is run, as in the simulator.}
        \end{itemize}
      \end{ideia}
      }
    \end{column}
  \end{columns}
\end{frame}

% ------------------------------------------------------------------
\begin{frame}{\tr{A história dos ajustamentos do Brasil}{The history of Brazil's adjustments}}
  \centering
  \begin{tikzpicture}[x=0.172cm,y=0.9cm,font=\scriptsize]
    \draw[thick, ->] (-1,0) -- (79,0);
    \foreach \a in {1950,1960,1970,1980,1990,2000,2010,2020} {
      \draw ({\a-1945},0.1) -- ({\a-1945},-0.1) node[below] {\a};
    }
    % oito ajustamentos parciais
    \foreach \a in {1948,1952} { \fill[corErro] ({\a-1945},0) circle (0.1); }
    \foreach \a in {1959,1962,1963,1966,1970,1975} { \fill[corFrame] ({\a-1945},0) circle (0.1); }
    \draw[decorate, decoration={brace, amplitude=4pt}] ({1948-1945},0.4) -- ({1975-1945},0.4);
    \node[above, align=center] at ({1961.5-1945},0.7) {\tr{\textbf{8 ajustamentos}}{\textbf{8 adjustments}}\\\tr{(1945--75)}{(1945--75)}};
    \node[corErro, below, font=\tiny] at ({1950-1945},-0.65) {\tr{datum de Torres}{Torres datum}};
    \node[corFrame, below, font=\tiny] at ({1967-1945},-0.65) {\tr{datum de Imbituba}{Imbituba datum}};
    % macrocircuitos
    \fill[corSist!80] ({1991-1945},0.15) rectangle ({1992-1945},0.45);
    \node[corSist, above, align=center] at ({1991.5-1945},0.5) {\tr{1991--92: macro-}{1991--92: macro-}\\\tr{circuitos (AAGP)}{loops (AAGP)}};
    % DALTI e ajuste simultâneo
    \fill[corTempo!70] ({2005-1945},0.15) rectangle ({2011-1945},0.45);
    \fill[corFrame] ({2011-1945},0) circle (0.13);
    \node[corFrame, above, align=center] at ({2007-1945},0.5) {\tr{2005--11: DALTI,}{2005--11: DALTI,}\\\tr{ajuste simultâneo}{simultaneous adj.}};
    % REALT
    \fill[corMar] ({2018-1945},0) circle (0.13);
    \node[corMar, below, align=center, font=\tiny] at ({2018-1945},-0.65) {REALT-2018};
  \end{tikzpicture}
  \bigskip

  \begin{columns}[T]
    \begin{column}{0.31\textwidth}
      \begin{cuidado}[title={\tr{Até 1975}{Up to 1975}}]
        \footnotesize\tr{Oito ajustamentos; a partir do quarto, por blocos.}{Eight adjustments;
        from the fourth onward, by blocks.}
      \end{cuidado}
    \end{column}
    \begin{column}{0.31\textwidth}
      \begin{regra}[title={\tr{1991--2011}{1991--2011}}]
        \footnotesize\tr{De macrocircuitos a um ajuste \emph{simultâneo}.}{From macro-loops to a
        \emph{simultaneous} adjustment.}
      \end{regra}
    \end{column}
    \begin{column}{0.31\textwidth}
      \begin{ideia}[title={\tr{2018}{2018}}]
        \footnotesize\tr{O mesmo ajuste, agora com \emph{gravidade real}.}{The same kind of
        adjustment, now with \emph{real gravity}.}
      \end{ideia}
    \end{column}
  \end{columns}
\end{frame}

\begin{frame}{\tr{1945--1975: oito ajustamentos}{1945--1975: eight adjustments}}
  \begin{columns}[T]
    \begin{column}{0.54\textwidth}
      \small
      \begin{itemize}
        \setlength{\itemsep}{4pt}
        \item \tr{Entre 1945 e 1975, o IBGE fez \bnum{oito ajustamentos}: 1948, 1952, 1959,
              1962, 1963, 1966, 1970 e 1975.}{Between 1945 and 1975, IBGE made \bnum{eight
              adjustments}: 1948, 1952, 1959, 1962, 1963, 1966, 1970 and 1975.}
        \item \tr{Só os \bnum{três primeiros} usaram \emph{todas} as observações disponíveis.}{Only
              the \bnum{first three} used \emph{all} the available observations.}
        \item \tr{A partir do quarto, com observações demais, a rede foi \bnum{partida em
              blocos}: algumas RRNN ficavam com altitude \emph{fixa}, calculada nos blocos
              anteriores.}{From the fourth on, with too many observations, the network was
              \bnum{split into blocks}: some benchmarks kept a \emph{fixed} altitude, computed in
              the previous blocks.}
        \item \tr{Os de 1948 e 1952 foram no datum de \textcolor{corErro}{\textbf{Torres}}; os
              seis de 1959 a 1975, no de \textcolor{corFrame}{\textbf{Imbituba}}.}{Those of 1948
              and 1952 were on the \textcolor{corErro}{\textbf{Torres}} datum; the six from 1959
              to 1975, on \textcolor{corFrame}{\textbf{Imbituba}}.}
      \end{itemize}
    \end{column}
    \begin{column}{0.42\textwidth}
      \begin{ideia}[title={\tr{Por que blocos?}{Why blocks?}}]
        \small\tr{O cálculo era feito com os meios da época: a rede crescia mais depressa do que
        se podia resolver, de uma vez, o sistema de equações. Partir em blocos era a saída
        prática.}{The computation used the means of the day: the network grew faster than the
        system of equations could be solved at once. Splitting into blocks was the practical
        way out.}
      \end{ideia}
      \medskip
      {\scriptsize\tr{Fonte: IBGE (2011), pp. 10--12; Siqueira et al. (2018), slides 3--4.}{
      Source: IBGE (2011), pp. 10--12; Siqueira et al. (2018), slides 3--4.}}
    \end{column}
  \end{columns}
\end{frame}

\begin{frame}{\tr{Blocos em cascata}{Blocks in cascade}}
  \begin{columns}[T]
    \begin{column}{0.58\textwidth}
      \centering
      \begin{tikzpicture}[x=0.76cm,y=0.9cm,font=\scriptsize]
        \foreach \x/\c/\n in {0/10/1, 3.3/18/2, 6.6/26/3} {
          \fill[corFrame!\c] (\x,0) rectangle (\x+3.4,2.2);
          \draw[corFrame, thick] (\x,0) rectangle (\x+3.4,2.2);
          \node[corFrame] at (\x+1.7,1.1) {\tr{bloco}{block} \n};
        }
        \foreach \x in {3.3,3.4,6.6,6.7} {
          \fill[corErro] (\x+0.05,0.5) circle (0.07);
          \fill[corErro] (\x+0.05,1.1) circle (0.07);
          \fill[corErro] (\x+0.05,1.7) circle (0.07);
        }
        \draw[seta, corTempo] (3.0,2.5) -- node[above, text=corTempo, font=\tiny] {\tr{altitudes fixadas}{fixed altitudes}} (4.5,2.5);
        \draw[seta, corTempo] (6.3,2.5) -- node[above, text=corTempo, font=\tiny] {\tr{altitudes fixadas}{fixed altitudes}} (7.8,2.5);
        \node[corErro, below] at (3.5,-0.1) {\tr{RRNN de fronteira}{boundary benchmarks}};
      \end{tikzpicture}
    \end{column}
    \begin{column}{0.38\textwidth}
      \small
      \begin{enumerate}
        \setlength{\itemsep}{3pt}
        \item \tr{Ajusta-se o bloco 1.}{Block 1 is adjusted.}
        \item \tr{As RRNN da fronteira viram \emph{pontos fixos} no bloco 2.}{The boundary
              benchmarks become \emph{fixed points} in block 2.}
        \item \tr{E assim por diante.}{And so on.}
      \end{enumerate}
      \medskip
      \tr{O bloco 1 \bnum{nunca} se beneficia das medidas dos blocos seguintes, e o desacordo
      que sobra ``empurra'' para o fim da cascata.}{Block 1 \bnum{never} benefits from the
      measurements of the following blocks, and the leftover disagreement gets ``pushed'' to the
      end of the cascade.}
    \end{column}
  \end{columns}
\end{frame}

\begin{frame}{\tr{Um teste pequeno: cascata $\times$ ajuste simultâneo}{A small test: cascade $\times$
simultaneous adjustment}}
  \begin{columns}[T]
    \begin{column}{0.54\textwidth}
      \small
      \tr{Com a rede de 4 nós, fazemos \emph{dois} ajustamentos:}{With the 4-node network, we do
      \emph{two} adjustments:}
      \begin{itemize}
        \setlength{\itemsep}{2pt}
        \item \tr{\textbf{simultâneo}: as 6 seções de uma vez;}{\textbf{simultaneous}: all 6
              sections at once;}
        \item \tr{\textbf{cascata}: bloco 1 ($A$, $B$, $C$: seções $AB$, $BC$, $AC$); depois $B$
              e $C$ \emph{fixos}, e o bloco 2 (seções $CD$, $DA$, $BD$) acha $D$.}{\textbf{cascade}:
              block 1 ($A$, $B$, $C$: sections $AB$, $BC$, $AC$); then $B$ and $C$ \emph{fixed},
              and block 2 (sections $CD$, $DA$, $BD$) finds $D$.}
      \end{itemize}
      \medskip
      \renewcommand{\arraystretch}{1.3}
      \begin{tabular}{@{}lrrr@{}}
        \toprule
        & $H_B$ (m) & $H_C$ (m) & $H_D$ (m) \\
        \midrule
        \tr{simultâneo}{simultaneous} & \num{12.3487} & \num{9.8670} & \num{11.1220} \\
        \tr{cascata}{cascade} & \num{12.3500} & \num{9.8670} & \num{11.1223} \\
        \tr{diferença (mm)}{difference (mm)} & \bnum{$+\num{1.3}$} & $+\num{0.0}$ & $+\num{0.3}$ \\
        \bottomrule
      \end{tabular}
    \end{column}
    \begin{column}{0.42\textwidth}
      \begin{cuidado}[title={\tr{Duas diferenças}{Two differences}}]
        \footnotesize
        \begin{itemize}
          \setlength{\itemsep}{2pt}
          \item \tr{As \emph{altitudes} diferem: $B$ mudou \num{1.3}\,mm, sem nenhuma medida nova
                --- só pela ordem de cálculo.}{The \emph{altitudes} differ: $B$ changed
                \num{1.3}\,mm, with no new measurement --- only because of the order of
                computation.}
          \item \tr{A \emph{incerteza} de $D$ sai \bnum{otimista}: $\sigma_D=\num{1.22}$\,mm na
                cascata (que trata $B$ e $C$ como sem erro) contra $\num{1.38}$\,mm no
                simultâneo --- \num{12}\,\% menor.}{The \emph{uncertainty} of $D$ comes out
                \bnum{optimistic}: $\sigma_D=\num{1.22}$\,mm in the cascade (which treats $B$ and
                $C$ as error-free) versus $\num{1.38}$\,mm in the simultaneous one --- \num{12}\,\%
                smaller.}
        \end{itemize}
      \end{cuidado}
    \end{column}
  \end{columns}
\end{frame}

\begin{frame}{\tr{1991--92: macrocircuitos}{1991--92: macro-loops}}
  \begin{columns}[T]
    \begin{column}{0.54\textwidth}
      \small
      \begin{itemize}
        \setlength{\itemsep}{4pt}
        \item \tr{O trabalho iniciado em \bnum{1988} divide a rede em \emph{sub-redes}
              de \bnum{macrocircuitos}, para reduzir distorções.}{The work begun in
              \bnum{1988} divides the network into \bnum{macro-loop} sub-networks
              to reduce distortions.}
        \item \tr{O resultado é o \bnum{AAGP}, o Ajustamento Altimétrico Global Preliminar: ainda
              \emph{não simultâneo}.}{The result is the \bnum{AAGP}, the Preliminary Global
              Altimetric Adjustment: still \emph{not simultaneous}.}
        \item \tr{O BDG usa \bnum{01/12/1992} como data de referência de cálculo do AAGP:
              \num{43079} RRNN. Depois vieram ajustes isolados e ramais.}{The BDG uses
              \bnum{1 December 1992} as the AAGP calculation reference date:
              \num{43079} benchmarks. Isolated adjustments and spurs followed.}
              \fonteslide{IBGE (2011), pp. 11--12.}
      \end{itemize}
    \end{column}
    \begin{column}{0.42\textwidth}
      \begin{ideia}[title={\tr{O macrocircuito, outra vez}{The macro-loop, again}}]
        \small\tr{Um macrocircuito é um circuito \emph{grande}, que contém vários. Ajustar cada
        sub-rede antes de juntá-las tira os piores empurrões da cascata, mas ainda deixa uma
        \bnum{costura} entre as partes.}{A macro-loop is a \emph{large} loop containing several.
        Adjusting each sub-network before joining them removes the worst pushes of the cascade,
        but still leaves a \bnum{seam} between the parts.}
      \end{ideia}
    \end{column}
  \end{columns}
\end{frame}

\begin{frame}{\tr{2005--2011: o ajuste simultâneo}{2005--2011: the simultaneous adjustment}}
  \begin{columns}[T]
    \begin{column}{0.54\textwidth}
      \small
      \begin{itemize}
        \setlength{\itemsep}{4pt}
        \item \tr{A equipe de altimetria do IBGE (\textbf{DALTI}, na Coordenação de Geodésia)
              começa em \bnum{2005} um ajustamento \bnum{de tudo ao mesmo tempo}.}{IBGE's
              altimetry team (\textbf{DALTI}, in the Geodesy Coordination) starts in
              \bnum{2005} an adjustment of \bnum{everything at the same time}.}
        \item \tr{Resultado, divulgado em \bnum{junho de 2011}: altitudes para \bnum{\num{69590}
              RRNN}, a partir de mais de \bnum{\num{74000} observações} feitas desde
              1945.}{Result, released in \bnum{June 2011}: altitudes for \bnum{\num{69590}
              benchmarks}, from more than \bnum{\num{74000} observations} made since
              1945.}
        \item \tr{Programa: o \textbf{GHOST} (canadense), feito para ajustar redes enormes.}{Software: \textbf{GHOST}
              (Canadian), made to adjust huge networks.}
        \item \tr{Ainda \emph{sem gravidade medida}: só a correção com gravidade
              teórica.}{Still \emph{without measured gravity}: only the correction with
              theoretical gravity.}
      \end{itemize}
    \end{column}
    \begin{column}{0.42\textwidth}
      \begin{regra}[title={\tr{66 anos depois}{66 years later}}]
        \small\tr{De 1945 a 2011 são \bnum{\num{66} anos}. O novo processamento retoma o
        ajuste simultâneo após a fase particionada; os três primeiros ajustes históricos
        também foram globais.}{From 1945 to 2011 is \bnum{\num{66} years}.
        The new processing returns to simultaneous adjustment after the partitioned phase;
        the first three historical adjustments were global too.}
      \end{regra}
      \medskip
      {\scriptsize\tr{Fonte: IBGE (2011), apresentação, p. 2, e introdução, p. 7.}{
      Source: IBGE (2011), presentation, p. 2, and introduction, p. 7.}}
    \end{column}
  \end{columns}
\end{frame}

\begin{frame}{\tr{Blocos que conversam}{Blocks that talk}}
  \begin{columns}[T]
    \begin{column}{0.58\textwidth}
      \centering
      \begin{tikzpicture}[x=0.82cm,y=0.95cm,font=\scriptsize]
        % cascata
        \node[corTempo, above] at (1.6,2.7) {\textbf{\tr{cascata}{cascade}}};
        \node[caixa, draw=corFrame, fill=corFrame!10, minimum width=0.9cm] (c1) at (0.4,1.6) {1};
        \node[caixa, draw=corFrame, fill=corFrame!10, minimum width=0.9cm] (c2) at (1.7,1.6) {2};
        \node[caixa, draw=corFrame, fill=corFrame!10, minimum width=0.9cm] (c3) at (3.0,1.6) {3};
        \draw[seta, corTempo] (c1) -- (c2);
        \draw[seta, corTempo] (c2) -- (c3);
        \node[corTempo, below=0.3cm, align=center] at (1.7,1.3) {\tr{uma mão só}{one way only}};
        % Helmert
        \node[corFrame, above] at (7.0,2.7) {\textbf{\tr{Helmert (GHOST)}{Helmert (GHOST)}}};
        \node[caixa, draw=corFrame, fill=corFrame!10, minimum width=0.9cm] (h1) at (5.4,1.9) {1};
        \node[caixa, draw=corFrame, fill=corFrame!10, minimum width=0.9cm] (h2) at (7.0,1.9) {2};
        \node[caixa, draw=corFrame, fill=corFrame!10, minimum width=0.9cm] (h3) at (8.6,1.9) {3};
        \node[caixa, draw=corErro, fill=corErro!8, minimum width=3.6cm] (hc) at (7.0,0.6) {\tr{sistema comum (fronteiras)}{common system (boundaries)}};
        \draw[seta2, corFrame] (h1) -- (hc);
        \draw[seta2, corFrame] (h2) -- (hc);
        \draw[seta2, corFrame] (h3) -- (hc);
        \node[corFrame, below=0.3cm, align=center] at (7.0,0.2) {\tr{ida e volta}{there and back}};
      \end{tikzpicture}
    \end{column}
    \begin{column}{0.38\textwidth}
      \small
      \tr{No \bnum{Helmert blocking}, os blocos são resolvidos \emph{por dentro} e só as
      incógnitas da \emph{fronteira} formam um sistema comum, resolvido de uma vez; depois cada
      bloco recebe a resposta de volta.}{In \bnum{Helmert blocking}, the blocks are solved
      \emph{inside} and only the \emph{boundary} unknowns form a common system, solved at once;
      then each block gets the answer back.}
      \medskip
      \begin{ideia}[title={\tr{O ponto}{The point}}]
        \footnotesize\tr{A solução é \bnum{exatamente} a do ajuste simultâneo: o bloco é só um
        truque de \emph{cálculo}. Na cascata, é uma \emph{aproximação}.}{The solution is
        \bnum{exactly} that of the simultaneous adjustment: the block is only a trick of
        \emph{computation}. In the cascade, it is an \emph{approximation}.}
      \end{ideia}
    \end{column}
  \end{columns}
\end{frame}

\begin{frame}{\tr{Sua vez}{Your turn}}
  \begin{columns}[T]
    \begin{column}{0.50\textwidth}
      \begin{regra}[title={\tr{Quantos circuitos tem o Brasil?}{How many loops does Brazil
        have?}}]
        \small\tr{O ajustamento de 2011 tinha \num{69590} RRNN e mais de \num{74000}
        observações (seções). Há \bnum{dois} blocos de rede (Imbituba e Santana), cada um com um
        ponto fixo. Quantas incógnitas? Quantas observações \emph{a mais}, no mínimo (a
        redundância)? Que fração é isso?}{The 2011 adjustment had \num{69590} benchmarks and more
        than \num{74000} observations (sections). There are \bnum{two} network blocks (Imbituba
        and Santana), each with one fixed point. How many unknowns? How many observations
        \emph{extra}, at least (the redundancy)? What fraction is that?}
      \end{regra}
    \end{column}
    \begin{column}{0.46\textwidth}
      \onslide<2->{
      \begin{ideia}[title={\tr{Respostas}{Answers}}]
        \small\raggedright
        $\num{69590}-2=\bnum{\num{69588}}$ \tr{incógnitas.}{unknowns.}
        \medskip

        $r\ge\num{74000}-\num{69588}=\bnum{\num{4412}}$.
        \medskip

        $\num{4412}/\num{74000}\approx\bnum{\num{6}\,\%}$: \tr{só \emph{uma em cada 17}
        seções ``fecha'' um circuito.}{only \emph{one in every 17} sections ``closes'' a
        loop.}
        \medskip

        \tr{A rede é quase uma \emph{árvore}: cada circuito vale muito.}{The network is almost a
        \emph{tree}: each loop is worth a lot.}
      \end{ideia}
      }
    \end{column}
  \end{columns}
\end{frame}

% ------------------------------------------------------------------
\begin{frame}{\tr{REALT-2018: o que mudou}{REALT-2018: what changed}}
  \begin{columns}[T]
    \begin{column}{0.31\textwidth}
      \begin{ideia}[title={\tr{1. Gravidade real}{1. Real gravity}}]
        \footnotesize\tr{A gravidade observada, complementada por \bnum{interpolação}, entra
        no cálculo em todas as RRNN. Em 2011 usava-se a correção com gravidade
        \emph{teórica}.}{Observed gravity, supplemented by \bnum{interpolation},
        enters the calculation at all benchmarks. In 2011 the correction used
        \emph{theoretical} gravity.}
      \end{ideia}
    \end{column}
    \begin{column}{0.31\textwidth}
      \begin{regra}[title={\tr{2. Outra incógnita}{2. A different unknown}}]
        \footnotesize\tr{O ajuste calcula \bnum{números geopotenciais} $C$ (m$^{2}$/s$^{2}$), não
        altitudes. A altitude \emph{normal} sai depois: $H^{*}=C/\bar\gamma$.}{The adjustment
        computes \bnum{geopotential numbers} $C$ (m$^{2}$/s$^{2}$), not altitudes. The
        \emph{normal} height comes afterward: $H^{*}=C/\bar\gamma$.}
      \end{regra}
    \end{column}
    \begin{column}{0.31\textwidth}
      \begin{cuidado}[title={\tr{3. Dados revistos}{3. Revised data}}]
        \footnotesize\tr{Revisão \bnum{completa} das bases, novas observações de nivelamento e de
        gravimetria, e inconsistências corrigidas.}{\bnum{Complete} review of the databases, new
        leveling and gravimetry observations, and inconsistencies corrected.}
      \end{cuidado}
    \end{column}
  \end{columns}
  \bigskip

  \centering\small
  \tr{\bnum{O que não mudou}: os dois zeros. O REALT-2018 manteve os data de Imbituba e de
  Santana; mudou o que se \emph{propaga} a partir deles.}{\bnum{What did not change}: the two
  zeros. REALT-2018 kept the Imbituba and Santana datums; what changed is what is
  \emph{propagated} from them.}
  \smallskip

  {\footnotesize\tr{Divulgado em 30/07/2018; 2ª edição do relatório em 30/09/2019.}{Released on
  30/07/2018; 2nd edition of the report on 30/09/2019.}}
\end{frame}

\begin{frame}{\tr{O caminho dos dados no REALT-2018}{The data path in REALT-2018}}
  \centering
  \begin{tikzpicture}[x=0.9cm,y=0.95cm,font=\scriptsize,
      cx/.style={caixa, text width=2.55cm, minimum height=1.3cm, font=\scriptsize}]
    \node[cx, draw=corFrame, fill=corFrame!10] (a) at (0,1.9) {\textbf{\tr{desníveis $\Delta n$}{height differences $\Delta n$}}\\\tr{(nivelamento, por seção)}{(leveling, per section)}};
    \node[cx, draw=corMar, fill=corMar!10] (b) at (0,0) {\textbf{\tr{gravidade $g$}{gravity $g$}}\\\tr{$>\num{100000}$ pontos $\to$ grade}{$>\num{100000}$ points $\to$ grid}};
    \node[cx, draw=corSist, fill=corSist!10] (c) at (3.9,0.95) {\textbf{$\Delta C=\bar g\,\Delta n$}\\\tr{por seção}{per section}};
    \node[cx, draw=corTempo, fill=corTempo!10] (d) at (7.8,0.95) {\textbf{\tr{ajustamento (MMQ)}{adjustment (LSQ)}}\\\tr{dois blocos: 4X e 9329T}{two blocks: 4X and 9329T}};
    \node[cx, draw=corErro, fill=corErro!8] (e) at (11.7,0.95) {\textbf{$C\pm\sigma_C\to H^{*}$}\\\tr{publicação (BDG)}{publication (BDG)}};
    \draw[seta] (a.east) -- (c.north west);
    \draw[seta] (b.east) -- (c.south west);
    \draw[seta] (c) -- (d);
    \draw[seta] (d) -- (e);
  \end{tikzpicture}
  \bigskip

  \small
  \tr{\bnum{Até 2011}: $\Delta n$ $\to$ correção com gravidade teórica $\to$ ajuste de altitudes.
  \bnum{Em 2018}: $\Delta n$ e $g$ observada/interpolada $\to$ $\Delta C$ $\to$ ajuste de $C$ $\to$
  $H^{*}$.}{\bnum{Up to 2011}: $\Delta n$ $\to$ correction with theoretical gravity $\to$
  adjustment of altitudes. \bnum{In 2018}: $\Delta n$ and observed/interpolated $g$ $\to$ $\Delta C$
  $\to$ adjustment of $C$ $\to$ $H^{*}$.}
\end{frame}

\begin{frame}{\tr{Onde a gravidade entra}{Where gravity comes in}}
  \begin{columns}[T]
    \begin{column}{0.54\textwidth}
      \small
      \tr{Para cada seção $ij$, a diferença de geopotencial é o desnível \emph{vezes a gravidade
      média} entre as duas pontas:}{For each section $ij$, the geopotential difference is the
      height difference \emph{times the mean gravity} between its two ends:}
      \[
        \Delta C_{ij}=\bar g_{ij}\,\Delta n_{ij},\qquad
        \bar g_{ij}\approx\frac{g_i+g_j}{2}
      \]
      \tr{Usaram-se dados de gravidade \bnum{observados e interpolados}.
      A interpolação reuniu \bnum{mais de \num{100000} pontos} do IBGE e do BNDG,
      com o programa Geogrid\_7c.}{Both \bnum{observed and interpolated} gravity values
      were used. Interpolation combined \bnum{more than \num{100000} points}
      from IBGE and BNDG using Geogrid\_7c.}
      \fonteslide{IBGE (2019), pp. 29--32; Siqueira et al. (2018), slide 20.}
      \medskip

      \tr{A mesma \emph{gravimetria} que alimenta o geoide agora também alimenta a rede de
      nivelamento.}{The same \emph{gravimetry} that feeds the geoid now also feeds the leveling
      network.}
    \end{column}
    \begin{column}{0.42\textwidth}
      \begin{cuidado}[title={\tr{Quanto a gravidade varia?}{How much does gravity vary?}}]
        \small\tr{Na Terra toda, $g$ vai de \num{9.764} a \num{9.834}\,m/s$^{2}$, cerca de
        \num{0.7}\,\% (valores citados pelo IBGE). Parece pouco, mas $\num{0.7}\,\%$ de uma
        altitude de \num{1000}\,m são \bnum{\num{7}\,m}: é por isso que a gravidade \emph{não}
        pode ser tratada como constante.}{Over the whole Earth, $g$ ranges from \num{9.764} to
        \num{9.834}\,m/s$^{2}$, about \num{0.7}\,\% (values quoted by IBGE). It looks small, but
        $\num{0.7}\,\%$ of an altitude of \num{1000}\,m is \bnum{\num{7}\,m}: that is why
        gravity can \emph{not} be treated as constant.}
      \end{cuidado}
      \medskip
      {\scriptsize\tr{O que importa para o nivelamento é a variação \emph{entre caminhos}
      (Aula 2); a variação com a latitude já é tratada pela gravidade normal.}{What matters for
      leveling is the variation \emph{between paths} (Class 2); the variation with latitude is
      already handled by normal gravity.}}
    \end{column}
  \end{columns}
\end{frame}

\begin{frame}{\tr{De $C$ para a altitude normal}{From $C$ to the normal height}}
  \begin{columns}[T]
    \begin{column}{0.54\textwidth}
      \small
      \tr{O resultado do ajuste é o número geopotencial $C$ de cada RN, com sua incerteza
      $\sigma_C$. A altitude sai por \bnum{uma divisão}:}{The result of the adjustment is the
      geopotential number $C$ of each benchmark, with its uncertainty $\sigma_C$. The altitude
      comes out by \bnum{one division}:}
      \[
        \boxed{\;H^{*}=\frac{C}{\bar\gamma}\;}
      \]
      \tr{$\bar\gamma$ é a gravidade \emph{normal} média ao longo da vertical da RN (fórmula do
      elipsoide, sem medir nada).}{$\bar\gamma$ is the mean \emph{normal} gravity along the
      benchmark's vertical (ellipsoid formula, nothing measured).}
    \end{column}
    \begin{column}{0.42\textwidth}
      \begin{ideia}[title={\tr{Exemplo (RN fictícia)}{Example (fictitious benchmark)}}]
        \small\raggedright
        $C=\num{68.392}$\,m$^{2}$/s$^{2}$
        \medskip

        $\bar\gamma=\num{9.78}$\,m/s$^{2}$
        \medskip

        $H^{*}=\dfrac{\num{68.392}}{\num{9.78}}=\bnum{\num{6.993}\,m}$
        \medskip

        \tr{Regra de bolso: \bnum{1\,m $\approx$ \num{9.8}\,m$^{2}$/s$^{2}$}.}{Rule of thumb:
        \bnum{1\,m $\approx$ \num{9.8}\,m$^{2}$/s$^{2}$}.}
      \end{ideia}
      \smallskip
      {\scriptsize\tr{Valores de exemplo, inventados; não são de uma RN real.}{Example values, made
      up; not from a real benchmark.}}
    \end{column}
  \end{columns}
\end{frame}

\begin{frame}{\tr{Dois blocos, dois datums, um método}{Two blocks, two datums, one method}}
  \begin{columns}[T]
    \begin{column}{0.52\textwidth}
      \centering
      \begin{tikzpicture}[x=0.8cm,y=0.8cm,font=\scriptsize]
        \fill[corFrame!12] (0,0) -- (4.6,0.2) -- (5.2,2.6) -- (3.0,3.9) -- (0.2,3.3) -- cycle;
        \draw[corFrame, thick] (0,0) -- (4.6,0.2) -- (5.2,2.6) -- (3.0,3.9) -- (0.2,3.3) -- cycle;
        \node[corFrame] at (2.6,2.0) {\tr{bloco principal}{main block}};
        \pic[scale=1.6] at (0.5,0.35) {marco};
        \node[corFrame, right, font=\tiny] at (0.85,0.55) {RN 4X \tr{(Imbituba)}{(Imbituba)}};
        \fill[corSist!12] (6.4,5.0) -- (8.0,4.8) -- (8.3,6.0) -- (6.8,6.4) -- cycle;
        \draw[corSist, thick] (6.4,5.0) -- (8.0,4.8) -- (8.3,6.0) -- (6.8,6.4) -- cycle;
        \node[corSist] at (7.3,5.7) {Amapá};
        \pic[scale=1.4] at (7.1,4.85) {marco};
        \node[corSist, below right, font=\tiny] at (7.0,4.75) {RN 9329T};
        \draw[corErro, dashed, thick, <->] (4.6,3.2) -- (6.6,4.9);
        \node[corErro, left, font=\tiny, align=right] at (5.4,3.9) {$\approx$\num{330}\,km\\\tr{sem nivelamento}{no leveling}};
      \end{tikzpicture}
    \end{column}
    \begin{column}{0.44\textwidth}
      \small
      \tr{O REALT-2018 faz \bnum{dois ajustamentos}:}{REALT-2018 makes \bnum{two
      adjustments}:}
      \begin{itemize}
        \setlength{\itemsep}{3pt}
        \item \tr{o \textbf{bloco principal}, preso à RN \textbf{4X} (Imbituba);}{the \textbf{main
              block}, pinned to benchmark \textbf{4X} (Imbituba);}
        \item \tr{a rede do \textbf{Amapá}, presa à RN \textbf{9329T} (Santana).}{the
              \textbf{Amapá} network, pinned to benchmark \textbf{9329T} (Santana).}
      \end{itemize}
      \medskip
      \tr{Não existe nivelamento entre eles. Cada bloco tem o \emph{seu} zero --- e as altitudes
      de um não são diretamente comparáveis às do outro.}{There is no leveling between them.
      Each block has \emph{its} zero --- and the altitudes of one are not directly comparable
      to those of the other.}
    \end{column}
  \end{columns}
\end{frame}

\begin{frame}{\tr{REALT-2018: resíduos do bloco principal}{REALT-2018: principal-block residuals}}
  \begin{columns}[T]
    \begin{column}{0.54\textwidth}
      \centering
      \begin{tikzpicture}[x=0.84cm,y=1cm,font=\scriptsize]
        \foreach \i/\k/\p in {0/1/92.7, 1/2/98.0, 2/3/99.1, 3/4/99.5} {
          \fill[corFrame!70] (0,{-0.8*\i}) rectangle ({0.05*\p},{-0.8*\i-0.5});
          \node[left] at (0,{-0.8*\i-0.25}) {$|v|\le\k\,\mathrm{mm}\sqrt K$};
          \node[right] at ({0.05*\p},{-0.8*\i-0.25}) {\num{\p}\,\%};
        }
        \draw[thick] (0,0.1) -- (0,-3.3);
        \node[below, text width=5.2cm, align=center] at (2.6,-3.3)
          {\tr{fração das seções dentro de cada intervalo simétrico}{share of sections within each symmetric interval}};
      \end{tikzpicture}
    \end{column}
    \begin{column}{0.42\textwidth}
      \footnotesize
      \tr{O IBGE apresenta as correções dos desníveis em função do comprimento:
      $v/\sqrt K$, com $K$ em km. Os percentuais são do \bnum{bloco principal do
      REALT-2018}.}{IBGE presents height-difference corrections as a function of length:
      $v/\sqrt K$, with $K$ in km. These percentages describe the \bnum{principal block
      of REALT-2018}.}
      \par\medskip
      \begin{cuidado}[title={\tr{Resultados do ajuste}{Adjustment results}}]
        \footnotesize
        \tr{Essas frequências avaliam a consistência das observações no modelo ajustado.
        A aceitação em campo usa as tolerâncias de seção, linha e circuito.
        A incerteza das altitudes exige também suas covariâncias.}{These frequencies
        assess observation consistency in the adjusted model. Field acceptance uses
        section, line and circuit tolerances. Height uncertainty also requires the
        covariance information.}
      \end{cuidado}
    \end{column}
  \end{columns}
  \fonteslide{IBGE (2019), REALT-2018, \tr{2.ª ed., p. 36, Figura 16; Siqueira et al. (2018), slide 31}{2nd ed., p. 36, Figure 16; Siqueira et al. (2018), slide 31}.}
\end{frame}

\begin{frame}{\tr{O que mudou nas RRNN, na prática}{What changed at the benchmarks, in practice}}
  \begin{columns}[T]
    \begin{column}{0.48\textwidth}
      \begin{regra}[title={\tr{Isolar a inserção da gravidade}{Isolating gravity integration}}]
        \small\tr{Com os mesmos dados e geometria, inserir a gravidade real produziu
        diferenças entre $-15$ e $+5$\,mm em \bnum{94\,\%} das RRNN.}{With the same
        data and geometry, integrating real gravity produced differences between
        $-15$ and $+5$\,mm at \bnum{94\,\%} of benchmarks.}
      \end{regra}
    \end{column}
    \begin{column}{0.48\textwidth}
      \begin{cuidado}[title={\tr{O reajustamento inteiro}{The whole readjustment}}]
        \small\tr{Cerca de \num{50000} RRNN (76\,\%) mudaram de $+20$ a $+30$\,cm, em boa parte
        do N, NE, Centro-Oeste e SE.}{About \num{50000} benchmarks (76\,\%) changed by $+20$ to
        $+30$\,cm, in much of the North, Northeast, Center-West and Southeast.}
      \end{cuidado}
    \end{column}
  \end{columns}
  \bigskip

  \centering\small
  \tr{No teste controlado do IBGE, o efeito isolado da gravidade ficou em milímetros
  para 94\,\% das RRNN. A depuração e a geometria da rede explicam grande parte das
  diferenças do reajustamento completo.}{In IBGE's controlled test, the isolated gravity
  effect was millimetric at 94\,\% of benchmarks. Data cleaning and network geometry
  explain much of the difference in the complete readjustment.}
  \smallskip

  {\footnotesize\tr{Fonte: IBGE (2019), p. 37, Figura 18; p. 40, Figura 20.
  São duas comparações diferentes.}{Source: IBGE (2019), p. 37, Figure 18;
  p. 40, Figure 20. These are two different comparisons.}}
\end{frame}

\begin{frame}{\tr{Sua vez}{Your turn}}
  \begin{columns}[T]
    \begin{column}{0.50\textwidth}
      \begin{regra}[title={\tr{De $C$ para $H^{*}$}{From $C$ to $H^{*}$}}]
        \small
        \begin{enumerate}
          \item[(a)] \tr{Uma RN tem $C=\num{68.392}$\,m$^{2}$/s$^{2}$ e $\bar\gamma=\num{9.78}$\,m/s$^{2}$.
                Qual a altitude normal?}{A benchmark has $C=\num{68.392}$\,m$^{2}$/s$^{2}$ and
                $\bar\gamma=\num{9.78}$\,m/s$^{2}$. What is the normal height?}
          \item[(b)] \tr{Se por engano se usasse $\bar\gamma=\num{9.80}$, quanto mudaria?}{If
                $\bar\gamma=\num{9.80}$ were used by mistake, how much would it change?}
          \item[(c)] \tr{E numa RN de $H^{*}=\num{1500}$\,m, com o mesmo erro relativo?}{And for a
                benchmark with $H^{*}=\num{1500}$\,m, with the same relative error?}
        \end{enumerate}
      \end{regra}
    \end{column}
    \begin{column}{0.46\textwidth}
      \onslide<2->{
      \begin{ideia}[title={\tr{Respostas}{Answers}}]
        \small\raggedright
        (a) $\num{68.392}/\num{9.78}=\bnum{\num{6.993}\,m}$.
        \medskip

        (b) $\num{68.392}/\num{9.80}=\num{6.979}$\,m: \tr{muda}{changes by}
        \bnum{\num{14}\,mm}.
        \medskip

        (c) \tr{O erro é \emph{relativo}: $\num{0.2}\,\%\times\num{1500}\approx\bnum{\num{3}\,m}$.
        Por isso $\bar\gamma$ é calculada com cuidado, ponto a ponto.}{The error is
        \emph{relative}: $\num{0.2}\,\%\times\num{1500}\approx\bnum{\num{3}\,m}$. That is why
        $\bar\gamma$ is computed carefully, point by point.}
      \end{ideia}
      }
    \end{column}
  \end{columns}
\end{frame}

\begin{frame}{\tr{Pare e pense}{Stop and think}}
  \begin{columns}[T]
    \begin{column}{0.46\textwidth}
      \begin{armadilha}[title={\tr{Um projeto, duas versões}{One project, two versions}}]
        \small\tr{Um projeto de drenagem usa a RN $P$ com altitude do \emph{ajustamento de
        2011} e a RN $Q$ com altitude do \emph{REALT-2018}. As duas são oficiais. O que pode dar
        errado?}{A drainage project uses benchmark $P$ with the altitude from the \emph{2011
        adjustment} and benchmark $Q$ with the altitude from \emph{REALT-2018}. Both are
        official. What can go wrong?}
      \end{armadilha}
    \end{column}
    \begin{column}{0.50\textwidth}
      \onslide<2->{
      \begin{ideia}[title={\tr{Resposta}{Answer}}]
        \footnotesize
        \begin{itemize}
          \setlength{\itemsep}{2pt}
          \item \tr{Se a região subiu $\approx$ $+25$\,cm de uma versão para a outra, misturar as
                duas cria um \bnum{degrau falso} de decímetros entre $P$ e $Q$ --- e a água corre
                pelo degrau.}{If the region rose about $+25$\,cm from one version to the other,
                mixing the two creates a \bnum{false step} of decimeters between $P$ and $Q$ ---
                and the water flows down the step.}
          \item \tr{A \emph{diferença} entre RRNN da \emph{mesma} versão é o que serve.}{The
                \emph{difference} between benchmarks of the \emph{same} version is what is
                useful.}
          \item \tr{Regra: anotar sempre \bnum{qual ajustamento} (ano e datum) deu a
                altitude.}{Rule: always note \bnum{which adjustment} (year and datum) gave the
                altitude.}
        \end{itemize}
      \end{ideia}
      }
    \end{column}
  \end{columns}
\end{frame}

\begin{frame}{\tr{As quatro perguntas, respondidas}{The four questions, answered}}
  \small
  \begin{enumerate}
    \setlength{\itemsep}{3pt}
    \item \textbf{\tr{Por que a mesma RN já teve várias altitudes?}{Why has the same benchmark
          had several altitudes?}} \tr{Porque cada ajustamento usa \bnum{dados novos}, um
          \bnum{método novo} e às vezes um \bnum{zero novo} --- e a altitude é o \emph{resultado}
          do ajustamento, não uma propriedade da pedra.}{Because each adjustment uses \bnum{new
          data}, a \bnum{new method} and sometimes a \bnum{new zero} --- and the altitude is the
          \emph{result} of the adjustment, not a property of the stone.}
    \item \textbf{\tr{O que é ajustar uma rede?}{What is adjusting a network?}}
          \tr{Achar uma altitude para cada RN, por \bnum{mínimos quadrados}: $\Delta h+v=H_j-H_i$,
          peso $1/K$, um ponto fixo no datum; repartir o desacordo e dar a incerteza.}{Finding one
          altitude for each benchmark by \bnum{least squares}: $\Delta h+v=H_j-H_i$, weight
          $1/K$, one fixed point at the datum; spreading the disagreement and giving the
          uncertainty.}
    \item \textbf{\tr{Como foram os ajustamentos de 1945 a 2011?}{How were the adjustments from
          1945 to 2011?}} \tr{Oito ajustes, os três primeiros globais (1948--75), macrocircuitos (1991--92)
          e a retomada do \bnum{ajuste simultâneo} (2005--11): \num{69590} RRNN, mais de \num{74000}
          observações.}{Eight adjustments, the first three global (1948--75), macro-loops (1991--92)
          and the return to \bnum{simultaneous adjustment} (2005--11): \num{69590} benchmarks, more
          than \num{74000} observations.}
    \item \textbf{\tr{O que o REALT-2018 mudou?}{What did REALT-2018 change?}}
          \tr{A incógnita virou o \bnum{número geopotencial}, com \bnum{gravidade real} em todas
          as RRNN, dados revistos e dois blocos (Imbituba e Santana); $H^{*}=C/\bar\gamma$ é
          derivado.}{The unknown became the \bnum{geopotential number}, with \bnum{real gravity}
          at all benchmarks, revised data and two blocks (Imbituba and Santana); $H^{*}=C/\bar\gamma$
          is derived.}
  \end{enumerate}
\end{frame}

\begin{frame}{\tr{Para levar para casa}{Take home}}
  \begin{columns}[T]
    \begin{column}{0.54\textwidth}
      \begin{regra}[title={\tr{Quatro ideias}{Four ideas}}]
        \small
        \begin{enumerate}
          \setlength{\itemsep}{3pt}
          \item \tr{Altitude \bnum{é resultado de um ajustamento}, não propriedade da
                pedra.}{Altitude \bnum{is the result of an adjustment}, not a property of the
                stone.}
          \item \tr{Ajustar $=$ mínimos quadrados com pesos $1/K$ e um datum fixo.}{Adjusting
                $=$ least squares with weights $1/K$ and a fixed datum.}
          \item \tr{Cascata de blocos $\neq$ ajuste simultâneo (Helmert blocking é
                \emph{exato}).}{Cascade of blocks $\neq$ simultaneous adjustment (Helmert
                blocking is \emph{exact}).}
          \item \tr{REALT-2018: $C$ com gravidade real; $H^{*}=C/\bar\gamma$; dois
                zeros.}{REALT-2018: $C$ with real gravity; $H^{*}=C/\bar\gamma$; two
                zeros.}
        \end{enumerate}
      \end{regra}
    \end{column}
    \begin{column}{0.42\textwidth}
      \begin{ideia}[title={\tr{Ponte para a Aula 4}{Bridge to Class 4}}]
        \small\tr{Os dois zeros ficaram. Mas \emph{como se relacionam}? Os marégrafos de hoje
        medem o mar com GNSS ao lado. E o futuro é um zero \bnum{global}: o IHRS.}{The two zeros
        remained. But \emph{how do they relate}? Today's tide gauges measure the sea with GNSS
        next to them. And the future is a \bnum{global} zero: the IHRS.}
      \end{ideia}
    \end{column}
  \end{columns}
\end{frame}

% ==================================================================
\section{\tr{Aula 4: O zero hoje e amanhã}{Class 4: The zero, now and next}}

\begin{frame}{\tr{As perguntas de hoje}{Today's questions}}
  \begin{columns}[T]
    \begin{column}{0.58\textwidth}
      \begin{enumerate}
        \setlength{\itemsep}{8pt}
        \item \tr{Como o IBGE vigia o mar de hoje em relação ao zero de ontem?}{How does IBGE
              watch today's sea against yesterday's zero?}
        \item \tr{Os dois zeros do Brasil concordam? Quanto diferem?}{Do Brazil's two zeros
              agree? How much do they differ?}
        \item \tr{O que o SIRGAS e o IHRS propõem para o futuro?}{What do SIRGAS and the IHRS
              propose for the future?}
        \item \tr{Na prática: nivelar, usar GNSS e modelo --- e o que fazer se a RN
              sumir?}{In practice: level, use GNSS and a model --- and what to do if the
              benchmark disappears?}
      \end{enumerate}
    \end{column}
    \begin{column}{0.38\textwidth}
      \begin{ideia}[title={\tr{Combinado}{Deal}}]
        \tr{No fim da aula, você vai saber responder às quatro --- com as suas
        palavras.}{By the end of the class you will be able to answer all four --- in your own
        words.}
        \medskip

        \tr{Elas voltam no último slide.}{They come back on the last slide.}
      \end{ideia}
    \end{column}
  \end{columns}
\end{frame}

\begin{frame}{\tr{O que a aula anterior deixou}{What the previous class left us}}
  \begin{columns}[T]
    \begin{column}{0.54\textwidth}
      \small
      \tr{Na Aula 3, quatro ideias:}{In Class 3, four ideas:}
      \begin{enumerate}
        \setlength{\itemsep}{3pt}
        \item \tr{a altitude é o \bnum{resultado de um ajustamento};}{the altitude is the
              \bnum{result of an adjustment};}
        \item \tr{ajustar é mínimos quadrados, com pesos $1/K$ e um ponto fixo no datum;}{adjusting
              is least squares, with weights $1/K$ and one point fixed at the datum;}
        \item \tr{a história foi: blocos $\to$ macrocircuitos $\to$ ajuste simultâneo
              (2011);}{the history was: blocks $\to$ macro-loops $\to$ simultaneous adjustment
              (2011);}
        \item \tr{o REALT-2018 ajusta $C$ com gravidade real; $H^{*}=C/\bar\gamma$.}{REALT-2018
              adjusts $C$ with real gravity; $H^{*}=C/\bar\gamma$.}
      \end{enumerate}
    \end{column}
    \begin{column}{0.42\textwidth}
      \begin{ideia}[title={\tr{O que ficou aberto}{What stayed open}}]
        \small\tr{Os dois \bnum{zeros} foram mantidos --- Imbituba e Santana --- e o mar de
        1949--57 já não é o de hoje. Como saber \emph{quanto} eles estão defasados? E como se
        relacionam entre si?}{The two \bnum{zeros} were kept --- Imbituba and Santana --- and
        the sea of 1949--57 is no longer today's. How do we know \emph{how much} they are off?
        And how do they relate to each other?}
      \end{ideia}
    \end{column}
  \end{columns}
\end{frame}

\begin{frame}{\tr{O fio da aula}{The thread of the class}}
  \centering
  \begin{tikzpicture}[x=1cm,y=1cm,font=\small]
    \node[caixa, draw=corMar, fill=corMar!8, text width=2.8cm, minimum height=2.1cm] (a) at (0,0)
      {\textbf{\tr{1.\ Vigiar o mar}{1.\ Watching the sea}}\\[2pt]
       \scriptsize\tr{RMPG, GNSS, nivelamento ao marégrafo}{RMPG, GNSS, leveling to the gauge}};
    \node[caixa, draw=corErro, fill=corErro!8, text width=2.8cm, minimum height=2.1cm] (b) at (3.5,0)
      {\textbf{\tr{2.\ Dois zeros}{2.\ Two zeros}}\\[2pt]
       \scriptsize\tr{quanto diferem, e como se estima}{how much they differ, and how it is estimated}};
    \node[caixa, draw=corSist, fill=corSist!8, text width=2.8cm, minimum height=2.1cm] (c) at (7.0,0)
      {\textbf{\tr{3.\ Um zero global}{3.\ A global zero}}\\[2pt]
       \scriptsize\tr{SIRGAS, IHRS, IHRF}{SIRGAS, IHRS, IHRF}};
    \node[caixa, draw=corFrame, fill=corFrame!8, text width=2.8cm, minimum height=2.1cm] (d) at (10.5,0)
      {\textbf{\tr{4.\ Na prática}{4.\ In practice}}\\[2pt]
       \scriptsize\tr{nivelar ou GNSS; RRNN; época; as três redes}{level or GNSS; benchmarks;
       epoch; the three networks}};
    \draw[seta] (a) -- (b);
    \draw[seta] (b) -- (c);
    \draw[seta] (c) -- (d);
  \end{tikzpicture}
\end{frame}

% ------------------------------------------------------------------
\begin{frame}{\tr{Por que vigiar o mar?}{Why watch the sea?}}
  \begin{columns}[T]
    \begin{column}{0.56\textwidth}
      \small
      \begin{itemize}
        \setlength{\itemsep}{4pt}
        \item \tr{Os níveis médios do mar que definem as RRNN foram calculados com dados de
              \bnum{épocas diferentes}, muitas vezes de décadas atrás.}{The mean sea levels that
              define the benchmarks were computed with data from \bnum{different epochs}, often
              decades ago.}
        \item \tr{Por isso o IBGE criou a \bnum{RMPG}, a Rede Maregráfica Permanente para
              Geodésia: ideia de \bnum{1996} (UFRJ e UFPR), a cargo da Coordenação de Geodésia
              desde \bnum{2001}.}{That is why IBGE created the \bnum{RMPG}, the Permanent Tide
              Gauge Network for Geodesy: an idea from \bnum{1996} (UFRJ and UFPR), run by the
              Geodesy Coordination since \bnum{2001}.}
        \item \tr{Objetivo: relacionar o datum vertical (Imbituba, Santana) com o \emph{nível do
              mar de hoje} e com outros níveis de referência da costa.}{Goal: to relate the
              vertical datum (Imbituba, Santana) to \emph{today's sea level} and to other
              coastal reference levels.}
      \end{itemize}
    \end{column}
    \begin{column}{0.40\textwidth}
      \begin{ideia}[title={\tr{Uma série de mais de um ciclo}{A series of more than one cycle}}]
        \small\tr{A série consolidada da RMPG vai de \bnum{julho de 2001 a dezembro de 2024}:
        cerca de \bnum{\num{23.5}\,anos}, mais que um ciclo nodal de \num{18.6}\,anos.}{The
        consolidated RMPG series runs from \bnum{July 2001 to December 2024}: about
        \bnum{\num{23.5}\,years}, more than one \num{18.6}-year nodal cycle.}
      \end{ideia}
      \medskip
      {\scriptsize\tr{Dados públicos, atualizados todo dia; os dados de tempo (meteorologia) são
      publicados desde 1/4/2022.}{Public data, updated daily; weather (meteorological) data have
      been published since 1/4/2022.}}
    \end{column}
  \end{columns}
\end{frame}

\begin{frame}{\tr{As seis estações}{The six stations}}
  \begin{columns}[T]
    \begin{column}{0.34\textwidth}
      \centering
      \begin{tikzpicture}[x=0.72cm,y=0.72cm,font=\scriptsize]
        \costaBR{corMar, thick}
        \coordinate (PImbituba) at (\ptImbituba);
        \coordinate (PArraialdoCabo) at (\ptArraialdoCabo);
        \coordinate (PSalvador) at (\ptSalvador);
        \coordinate (PFortaleza) at (\ptFortaleza);
        \coordinate (PBelem) at (\ptBelem);
        \coordinate (PSantana) at (\ptSantana);
        \foreach \q in {PImbituba, PArraialdoCabo, PSalvador, PFortaleza, PBelem} { \fill[corFrame] (\q) circle (0.1); }
        \fill[corSist] (PSantana) circle (0.1);
        \node[corFrame, right=2pt, font=\tiny] at (PImbituba) {Imbituba};
        \node[corFrame, right=2pt, font=\tiny] at (PArraialdoCabo) {Arraial do Cabo};
        \node[corFrame, right=2pt, font=\tiny] at (PSalvador) {Salvador};
        \node[corFrame, right=2pt, font=\tiny] at (PFortaleza) {Fortaleza};
        \node[corFrame, below right=1pt, font=\tiny] at (PBelem) {Belém};
        \node[corSist, left=2pt, font=\tiny] at (PSantana) {Santana};
      \end{tikzpicture}
    \end{column}
    \begin{column}{0.62\textwidth}
      \small
      \tr{Rede do IBGE na Nota Técnica 01/2024:}{IBGE network in Technical Note 01/2024:}
      \medskip

      \renewcommand{\arraystretch}{1.2}
      \begin{tabular}{@{}lL{4.8cm}@{}}
        \toprule
        Imbituba (SC) & \tr{o \bnum{datum} do Sul e Centro do país}{the \bnum{datum} of the South and Center of the country} \\
        Arraial do Cabo (RJ) & \tr{costa do Sudeste}{Southeast coast} \\
        Salvador (BA) & \tr{costa do Nordeste}{Northeast coast} \\
        Fortaleza (CE) & \tr{costa do Nordeste}{Northeast coast} \\
        Belém (PA) & \tr{foz do Amazonas}{Amazon mouth} \\
        Santana (AP) & \tr{o \bnum{datum} do Amapá}{the \bnum{datum} of Amapá} \\
        \bottomrule
      \end{tabular}
      \bigskip

      {\footnotesize\tr{Parceiros: Marinha (DHN), UFPR, Univ.\ do Havaí, Petrobras, FURG e
      administrações portuárias. Uma fonte mais antiga fala em cinco estações.}{Partners: Navy
      (DHN), UFPR, Univ.\ of Hawaii, Petrobras, FURG and port authorities. An older source
      mentions five stations.}}
    \end{column}
  \end{columns}
\end{frame}

\begin{frame}{\tr{O que há numa estação}{What is in a station}}
  \begin{columns}[T]
    \begin{column}{0.56\textwidth}
      \centering
      \begin{tikzpicture}[x=0.9cm,y=0.9cm,font=\scriptsize]
        \fill[mStoneMid!20] (0,0) rectangle (3.4,0.9);
        \draw[thick, mStoneMid] (0,0.9) -- (3.4,0.9);
        \fill[corMar!14] (3.4,0) rectangle (7.0,0.9);
        \draw[corMar, very thick, decorate, decoration={snake, amplitude=0.7pt, segment length=8pt}] (3.4,0.9) -- (7.0,0.9);
        % GNSS
        \draw[thick] (1.0,0.9) -- (1.0,2.2);
        \pic[scale=1.7] at (1.0,2.2) {antena};
        \node[corFrame, right] at (1.3,2.7) {GNSS};
        % sensor de nível
        \draw[thick] (4.4,0.0) -- (4.4,2.4) (4.9,0.0) -- (4.9,2.4);
        \draw[thick, fill=white] (4.48,0.82) rectangle (4.82,1.0);
        \node[left, font=\tiny] at (4.35,2.3) {\tr{sensor de nível}{level sensor}};
        % meteorologia
        \draw[thick] (6.3,0.9) -- (6.3,2.0);
        \draw[thick, fill=white] (6.1,2.0) circle (0.12);
        \draw[thick, fill=white] (6.5,2.0) circle (0.12);
        \node[above, font=\tiny] at (6.3,2.15) {\tr{tempo}{weather}};
        % RN e nivelamento
        \pic[scale=1.5] at (2.3,0.9) {marco};
        \node[corFrame, below] at (2.3,0.85) {RN};
        \draw[corTempo, thick, <->, dashed] (2.5,1.2) -- (4.35,1.2);
        \node[corTempo, above, font=\tiny] at (3.45,1.2) {\tr{nivelamento}{leveling}};
      \end{tikzpicture}
    \end{column}
    \begin{column}{0.40\textwidth}
      \scriptsize
      \begin{description}
        \setlength{\itemsep}{2pt}
        \item[\textcolor{corMar}{\textbf{\tr{nível}{level}}}] \tr{o sensor mede o mar \emph{em
              relação ao chão}.}{the sensor measures the sea \emph{relative to the ground}.}
        \item[\textcolor{corTempo}{\textbf{\tr{nivelamento}{leveling}}}] \tr{liga o zero do
              sensor a uma RN do SGB: \emph{assim a leitura vira altitude}.}{ties the sensor's zero
              to an SGB benchmark: \emph{so the reading becomes altitude}.}
        \item[\textcolor{corFrame}{\textbf{GNSS}}] \tr{o movimento vertical do \emph{chão}.}{the
              vertical motion of the \emph{ground}.}
        \item[\textbf{\tr{tempo}{weather}}] \tr{vento e pressão ajudam a entender saltos no
              nível.}{wind and pressure help understand jumps in the level.}
      \end{description}
      \medskip
      {\scriptsize\tr{O manual do IBGE (Manual Técnico n.º 11, 2010) trata também de
      gravimetria nas estações.}{IBGE's manual (Technical Manual No.~11, 2010) also covers
      gravimetry at the stations.}}
    \end{column}
  \end{columns}
\end{frame}

\begin{frame}{\tr{Da leitura do sensor à altitude do mar}{From the sensor reading to the sea's altitude}}
  \begin{columns}[T]
    \begin{column}{0.54\textwidth}
      \small
      \tr{Um exemplo \emph{inventado}, com a RN da estação:}{An \emph{invented} example, with the
      station's benchmark:}
      \medskip

      \renewcommand{\arraystretch}{1.15}
      \footnotesize
      \begin{tabular}{@{}lr@{}}
        \toprule
        \tr{altitude da RN (SGB)}{altitude of the benchmark (SGB)} & $\num{3.250}$\,m \\
        \tr{a RN está acima do zero do sensor}{the benchmark is above the sensor's zero} & $\num{4.100}$\,m \\
        \tr{$\Rightarrow$ altitude do zero do sensor}{$\Rightarrow$ altitude of the sensor's zero} & $-\num{0.850}$\,m \\
        \tr{leitura média do nível do mar}{mean sea level reading} & $+\num{0.995}$\,m \\
        \midrule
        \textbf{\tr{NMM local, em altitude do SGB}{local MSL, in SGB altitude}} & \bnum{$+\num{0.145}$\,m} \\
        \bottomrule
      \end{tabular}
    \end{column}
    \begin{column}{0.42\textwidth}
      \begin{ideia}[title={\tr{O que significa}{What it means}}]
        \scriptsize\tr{Neste exemplo, o mar \emph{daquela janela} está \num{14.5}\,cm \emph{acima}
        do zero do datum. É exatamente o que a RMPG calcula, estação por estação, e que os
        documentos de ``níveis de referência'' de cada porto trazem.}{In this example, the sea
        \emph{of that window} is \num{14.5}\,cm \emph{above} the datum's zero. That is exactly
        what the RMPG computes, station by station, and what each port's ``reference levels''
        documents contain.}
      \end{ideia}
      \medskip
      \begin{cuidado}[title={\tr{Cada porto é um caso}{Each port is a case}}]
        \scriptsize\tr{O valor muda de estação para estação: topografia da superfície do mar,
        época e movimento do chão.}{The value changes from station to station: sea surface
        topography, epoch and ground motion.}
      \end{cuidado}
    \end{column}
  \end{columns}
\end{frame}

\begin{frame}{\tr{O GNSS no marégrafo: o chão manda}{GNSS at the tide gauge: the ground rules}}
  \begin{columns}[T]
    \begin{column}{0.54\textwidth}
      \small
      \tr{A antena GNSS dá a \bnum{velocidade vertical do chão} (mm/ano), em relação ao centro
      da Terra. Pequena por ano, \emph{grande em décadas}:}{The GNSS antenna gives the
      \bnum{vertical velocity of the ground} (mm/year), relative to the Earth's center. Small
      per year, \emph{large over decades}:}
      \renewcommand{\arraystretch}{1.4}
      \medskip

      \begin{tabular}{@{}rrr@{}}
        \toprule
        \tr{velocidade}{velocity} & \tr{em 10 anos}{in 10 years} & \tr{em 30 anos}{in 30 years} \\
        \midrule
        $+\num{3}$\,mm/\tr{ano}{yr} & $+\num{30}$\,mm & $+\num{90}$\,mm \\
        $-\num{3}$\,mm/\tr{ano}{yr} & $-\num{30}$\,mm & $-\num{90}$\,mm \\
        \bottomrule
      \end{tabular}
    \end{column}
    \begin{column}{0.42\textwidth}
      \begin{ideia}[title={\tr{Separar mar e chão}{Separating sea and ground}}]
        \footnotesize\tr{O marégrafo vê \emph{mar menos chão}. O GNSS vê só o chão. A diferença dá o
        \bnum{mar de verdade}.}{The tide gauge sees \emph{sea minus ground}. GNSS sees only the
        ground. The difference gives the \bnum{real sea}.}
      \end{ideia}
      \medskip
      \begin{cuidado}[title={\tr{Imbituba}{Imbituba}}]
        \scriptsize\tr{Um estudo da UFPR (2017) achou, com campanhas GNSS, um afundamento do
        marégrafo de cerca de \num{3}\,mm/ano --- em 72,5 anos seriam $\approx$ \num{22}\,cm.
        Valor de um estudo; campanhas curtas.}{A UFPR study (2017) found, with GNSS campaigns, a
        sinking of the tide gauge of about \num{3}\,mm/year --- over 72.5 years that would be
        $\approx$ \num{22}\,cm. Value from one study; short campaigns.}
      \end{cuidado}
    \end{column}
  \end{columns}
\end{frame}

\begin{frame}{\tr{Sua vez}{Your turn}}
  \begin{columns}[T]
    \begin{column}{0.50\textwidth}
      \begin{regra}[title={\tr{O mar numa estação}{The sea at a station}}]
        \small\tr{Numa estação, a RN tem altitude \num{2.870}\,m e está \num{3.540}\,m acima
        do zero do sensor. A leitura média do mar, numa janela de 19 anos, foi
        $+\num{0.812}$\,m sobre o zero do sensor. Qual a altitude do NMM local no sistema do
        SGB? Está acima ou abaixo do datum?}{At a station, the benchmark has altitude
        \num{2.870}\,m and is \num{3.540}\,m above the sensor's zero. The mean sea reading, over
        a 19-year window, was $+\num{0.812}$\,m above the sensor's zero. What is the altitude of
        the local MSL in the SGB system? Is it above or below the datum?}
      \end{regra}
    \end{column}
    \begin{column}{0.46\textwidth}
      \onslide<2->{
      \begin{ideia}[title={\tr{Resposta}{Answer}}]
        \small\raggedright
        \tr{Zero do sensor:}{Sensor's zero:} $\num{2.870}-\num{3.540}=-\num{0.670}$\,m.
        \medskip

        \tr{NMM:}{MSL:} $-\num{0.670}+\num{0.812}=\bnum{+\num{0.142}\,m}$.
        \medskip

        \tr{Está \bnum{\num{14.2}\,cm acima} do zero do datum --- coerente com um mar que
        subiu desde a época do datum.}{It is \bnum{\num{14.2}\,cm above} the datum's zero ---
        consistent with a sea that has risen since the datum's epoch.}
      \end{ideia}
      }
    \end{column}
  \end{columns}
\end{frame}

% ------------------------------------------------------------------
\begin{frame}{\tr{Os dois zeros concordam?}{Do the two zeros agree?}}
  \begin{columns}[T]
    \begin{column}{0.52\textwidth}
      \centering
      \begin{tikzpicture}[x=1cm,y=1.15cm,font=\scriptsize]
        \draw[corFrame, very thick] (0,0) -- (3.0,0);
        \node[corFrame, below] at (1.5,-0.05) {\tr{zero de Imbituba}{Imbituba zero}};
        \draw[corSist, very thick] (4.2,1.30) -- (7.2,1.30);
        \node[corSist, above] at (5.7,1.35) {\tr{zero de Santana}{Santana zero}};
        \fill[corErro!15] (3.1,0.45) rectangle (4.1,1.40);
        \draw[corErro, thick, <->] (3.6,0) -- node[right, text=corErro, font=\small] {$\approx\num{1.3}$\,m} (3.6,1.30);
        \draw[dotted] (3.0,0) -- (3.6,0);
        \draw[dotted] (3.6,1.30) -- (4.2,1.30);
        \draw[mStoneMid, dashed] (0,1.30) -- (3.0,1.30);
        \node[mStoneMid, font=\tiny, above] at (1.5,1.3) {\tr{(se fossem iguais, aqui)}{(if equal, here)}};
        \node[corErro, below, font=\tiny, align=center] at (3.6,-0.6) {\tr{faixa dos estudos: 0,45 a 1,40\,m}{range of studies: 0.45 to 1.40\,m}};
      \end{tikzpicture}
    \end{column}
    \begin{column}{0.44\textwidth}
      \small
      \tr{Não há nivelamento entre as duas redes. Então a diferença entre os zeros só pode ser
      \emph{estimada}, com GNSS, gravimetria e modelos.}{There is no leveling between the two
      networks. So the difference between the zeros can only be \emph{estimated}, with GNSS,
      gravimetry and models.}
      \medskip

      \tr{Um estudo com um modelo global do satélite GOCE achou a rede de Santana
      \bnum{\num{1.30}\,m} acima da de Imbituba; outros achados vão de \num{0.45} a
      \num{1.40}\,m, conforme a estratégia.}{A study with a global model from the GOCE satellite
      found the Santana network \bnum{\num{1.30}\,m} above the Imbituba one; other findings range
      from \num{0.45} to \num{1.40}\,m, depending on the strategy.}
      \medskip

      {\scriptsize\tr{Valores de estudos (2013--2019), não oficiais; do resumo de busca (apêndice
      A8).}{Values from studies (2013--2019), not official; from the search summary (appendix
      A8).}}
    \end{column}
  \end{columns}
\end{frame}

\begin{frame}{\tr{Como estimar a diferença sem nivelar}{How to estimate the difference without leveling}}
  \begin{columns}[T]
    \begin{column}{0.54\textwidth}
      \footnotesize
      \tr{Em cada RN com GNSS, compara-se a altura do satélite com a altitude da rede e com um
      \emph{modelo} de quase-geoide:}{At each benchmark with GNSS, the satellite height is
      compared with the network altitude and with a quasigeoid \emph{model}:}
      \[
        r_i = h_i - H^{*}_i - \zeta_i^{\text{\tr{modelo}{model}}}
      \]
      \tr{Se o zero da rede coincidisse com o do modelo, os $r_i$ ficariam em torno de
      \emph{zero}. A \bnum{média} dos $r_i$ em cada rede é o \emph{deslocamento do seu datum}
      $\delta$:}{If the network's zero coincided with the model's, the $r_i$ would scatter
      around \emph{zero}. The \bnum{mean} of the $r_i$ in each network is the \emph{offset of its
      datum} $\delta$:}
      \[
        \delta_S-\delta_I=\overline{r}_{S}-\overline{r}_{I}
      \]
    \end{column}
    \begin{column}{0.42\textwidth}
      \begin{ideia}[title={\tr{A ideia}{The idea}}]
        \footnotesize\tr{É a mesma conta do \emph{GNSS/nivelamento} da aula do quase-geoide, usada
        para um fim novo: medir a \bnum{defasagem entre zeros}, em vez de testar o
        modelo.}{It is the same \emph{GNSS/leveling} computation of the quasigeoid class, used
        for a new purpose: measuring the \bnum{lag between zeros}, instead of testing the
        model.}
      \end{ideia}
      \medskip
      \begin{cuidado}[title={\tr{O limite}{The limit}}]
        \scriptsize\tr{O resultado vale tanto quanto o \emph{modelo}: onde a gravimetria é
        esparsa (a Amazônia), o erro do modelo pode ser de decímetros.}{The result is only as good
        as the \emph{model}: where gravimetry is sparse (the Amazon), the model error can be
        decimeters.}
      \end{cuidado}
    \end{column}
  \end{columns}
\end{frame}

\begin{frame}{\tr{Por que os estudos divergem?}{Why do the studies differ?}}
  \begin{columns}[T]
    \begin{column}{0.50\textwidth}
      \small
      \begin{itemize}
        \setlength{\itemsep}{4pt}
        \item \tr{Cada estudo usa um \bnum{modelo global} e uma \bnum{estratégia} diferentes
              (GNSS e modelo; nivelamento com gravimetria; ligação física).}{Each study uses a
              different \bnum{global model} and \bnum{strategy} (GNSS and model; leveling with
              gravimetry; physical connection).}
        \item \tr{Na região entre as duas redes (a foz do Amazonas e Marajó), a gravimetria é
              \bnum{rara} e o nivelamento não passa.}{In the region between the two networks (the
              Amazon mouth and Marajó), gravimetry is \bnum{scarce} and leveling cannot
              pass.}
        \item \tr{Os dois datums são NMM de \bnum{épocas e janelas} diferentes (9 anos $\times$ 2
              anos).}{The two datums are MSL of different \bnum{epochs and windows} (9 years
              $\times$ 2 years).}
      \end{itemize}
    \end{column}
    \begin{column}{0.46\textwidth}
      \begin{ideia}[title={\tr{Uma régua do problema}{A yardstick of the problem}}]
        \small\tr{Para os sistemas verticais sul-americanos como um todo, estimativas globais
        dão exatidão de \bnum{centímetros} onde há bons dados, \bnum{decímetros} onde são
        esparsos e até \bnum{1\,m} em casos extremos.}{For the South American vertical systems as
        a whole, global estimates give an accuracy of \bnum{centimeters} where data are good,
        \bnum{decimeters} where they are sparse and up to \bnum{1\,m} in extreme cases.}
      \end{ideia}
      \medskip
      {\scriptsize\tr{Valores de apresentação do DGFI-TUM (resumo de busca).}{Values from a
      DGFI-TUM presentation (search summary).}}
    \end{column}
  \end{columns}
\end{frame}

\begin{frame}{\tr{Sua vez}{Your turn}}
  \begin{columns}[T]
    \begin{column}{0.50\textwidth}
      \begin{regra}[title={\tr{A defasagem entre as redes}{The lag between networks}}]
        \small\tr{Em três RRNN do bloco de Imbituba, $r=h-H^{*}-\zeta$ deu $\num{0.10}$,
        $\num{0.14}$ e $\num{0.12}$\,m. Em três RRNN do Amapá: $\num{1.38}$, $\num{1.44}$ e
        $\num{1.42}$\,m (números inventados). Qual a defasagem $\delta_S-\delta_I$?}{At three
        benchmarks of the Imbituba block, $r=h-H^{*}-\zeta$ gave $\num{0.10}$, $\num{0.14}$ and
        $\num{0.12}$\,m. At three benchmarks in Amapá: $\num{1.38}$, $\num{1.44}$ and
        $\num{1.42}$\,m (made-up numbers). What is the lag $\delta_S-\delta_I$?}
      \end{regra}
    \end{column}
    \begin{column}{0.46\textwidth}
      \onslide<2->{
      \begin{ideia}[title={\tr{Resposta}{Answer}}]
        \small\raggedright
        $\overline r_I=\num{0.120}$\,m; \quad $\overline r_S=\num{1.413}$\,m.
        \medskip

        $\delta_S-\delta_I=\bnum{\num{1.293}\,m}\approx\num{1.3}$\,m.
        \medskip

        \tr{Dentro de cada rede os $r$ concordam entre si (alguns cm); entre as redes há um
        \emph{degrau} de $\approx$ \num{1.3}\,m.}{Within each network the $r$ agree with each
        other (a few cm); between networks there is a \emph{step} of $\approx$
        \num{1.3}\,m.}
      \end{ideia}
      }
    \end{column}
  \end{columns}
\end{frame}

\begin{frame}{\tr{O que fazer enquanto isso?}{What to do meanwhile?}}
  \begin{columns}[T]
    \begin{column}{0.54\textwidth}
      \small
      \begin{enumerate}
        \setlength{\itemsep}{4pt}
        \item \tr{\textbf{Não misture} altitudes das duas redes sem dizer: a diferença pode ser de
              \emph{mais de um metro}.}{\textbf{Do not mix} altitudes of the two networks without
              saying so: the difference may be \emph{over a meter}.}
        \item \tr{O \textbf{hgeoHNOR2020} entrega $H^{*}$ no datum \emph{local}: foi feito para
              concordar com as referências oficiais, cada qual no seu zero.}{\textbf{hgeoHNOR2020}
              delivers $H^{*}$ in the \emph{local} datum: it was made to agree with the official
              benchmarks, each in its own zero.}
        \item \tr{Um projeto que cruza o Amazonas deve \bnum{declarar} o datum de cada trecho.}{A
              project crossing the Amazon must \bnum{state} the datum of each stretch.}
      \end{enumerate}
    \end{column}
    \begin{column}{0.42\textwidth}
      \begin{cuidado}[title={\tr{O erro clássico}{The classic mistake}}]
        \small\tr{Somar altitudes de Macapá (datum Santana) com cotas do Pará (datum Imbituba) e
        achar que ``a obra fecha''. Fecha no papel; no chão, a água corre por um degrau
        escondido.}{Adding altitudes of Macapá (Santana datum) to elevations in Pará (Imbituba
        datum) and thinking ``the job closes''. It closes on paper; on the ground, water flows
        down a hidden step.}
      \end{cuidado}
    \end{column}
  \end{columns}
\end{frame}

% ------------------------------------------------------------------
\begin{frame}{\tr{Por que não simplesmente trocar o zero?}{Why not simply change the zero?}}
  \begin{columns}[T]
    \begin{column}{0.50\textwidth}
      \centering
      \begin{tikzpicture}[x=1cm,y=1.25cm,font=\scriptsize]
        \draw[corSist, very thick, dashed] (0,1.3) -- (7.0,1.3);
        \node[corSist, above right] at (0,1.3) {\tr{zero global (IHRS)}{global zero (IHRS)}};
        \foreach \x/\y/\n/\pos in {0.3/0.35/A/below, 1.9/2.0/B/above, 3.5/0.8/C/below, 5.1/2.45/D/above} {
          \draw[corFrame, very thick] (\x,\y) -- (\x+1.1,\y);
          \node[corFrame, \pos, font=\tiny] at (\x+0.55,\y) {\tr{país}{country} \n};
        }
        \draw[<->, corErro] (0.85,0.35) -- (0.85,1.3);
        \draw[<->, corErro] (5.65,1.3) -- (5.65,2.45);
      \end{tikzpicture}\\[1pt]
      {\scriptsize\tr{Cada país tem o \emph{seu} zero local; entre vizinhos a diferença pode chegar a 1--2\,m.}{Each country has \emph{its} local zero; between neighbors the difference can reach 1--2\,m.}}
    \end{column}
    \begin{column}{0.46\textwidth}
      \small
      \begin{itemize}
        \setlength{\itemsep}{3pt}
        \item \tr{Trocar o zero muda \bnum{todas} as altitudes: mapas, projetos, leis.}{Changing
              the zero changes \bnum{all} altitudes: maps, projects, laws.}
        \item \tr{Faltava um zero \emph{melhor}: só em \bnum{2015} a IAG definiu o \textbf{IHRS},
              com um \emph{valor de geopotencial} $W_0=\num{62636853.4}$\,m$^{2}$/s$^{2}$,
              não um marégrafo.}{A \emph{better} zero was missing: only in \bnum{2015} did the IAG
              define the \textbf{IHRS}, with a \emph{geopotential value}
              $W_0=\num{62636853.4}$\,m$^{2}$/s$^{2}$, not a tide gauge.}
        \item \tr{A rota é \bnum{ligar} cada datum ao IHRS (um deslocamento por datum), e só
              depois decidir o que mudar.}{The route is to \bnum{link} each datum to the IHRS (one
              offset per datum), and only then decide what to change.}
      \end{itemize}
    \end{column}
  \end{columns}
\end{frame}

\begin{frame}{\tr{O SIRGAS e o GT III}{SIRGAS and WG III}}
  \centering
  \begin{tikzpicture}[x=0.92cm,y=1cm,font=\scriptsize,
      ev/.style={caixa, text width=2.9cm, minimum height=1.9cm, font=\scriptsize}]
    \draw[thick, ->] (-1.2,0) -- (13.3,0);
    % Keep \c available for Portuguese cedillas inside the timeline labels.
    % Preserve \c para as cedilhas nos textos da linha do tempo.
    \foreach \x/\a/\eventColor/\t in {0.9/1997/corTempo/{\tr{o GT III nasce na assembleia da IAG, no Rio}{WG III is born at the IAG assembly, in Rio}},
                            4.5/2000/corFrame/{\tr{campanha SIRGAS2000, com GPS em marégrafos}{SIRGAS2000 campaign, with GPS at tide gauges}},
                            8.1/2001/corSist/{\tr{o GT III se consolida, em Cartagena}{WG III is consolidated, in Cartagena}},
                            11.7/2015/corMar/{\tr{IHRS (IAG): definição do sistema global}{IHRS (IAG): global system definition}}} {
      \fill[\eventColor] (\x,0) circle (0.13);
      \node[above=2pt, text=\eventColor, font=\small] at (\x,0) {\textbf{\a}};
      \node[ev, draw=\eventColor, fill=\eventColor!8, below=10pt] at (\x,0) {\t};
    }
  \end{tikzpicture}
  \bigskip

  \small
  \tr{O \bnum{GT III} (Datum Vertical) do SIRGAS nasceu em \bnum{1997} para \emph{unificar} as
  altitudes das Américas: cada país tinha o seu marégrafo, época e janela.}{SIRGAS's
  \bnum{Working Group III} (Vertical Datum) was born in \bnum{1997} to \emph{unify} the heights
  of the Americas: each country had its own tide gauge, epoch and window.}
\end{frame}

\begin{frame}{\tr{O que o IHRF pede de cada país}{What the IHRF asks of each country}}
  \begin{columns}[T]
    \begin{column}{0.54\textwidth}
      \small
      \tr{A realização vertical do SIRGAS será uma \bnum{densificação regional do IHRF}, com:}{The
      vertical realization of SIRGAS will be a \bnum{regional densification of the IHRF}, with:}
      \begin{itemize}
        \setlength{\itemsep}{3pt}
        \item \tr{as estações de referência GNSS do SIRGAS;}{the SIRGAS GNSS reference
              stations;}
        \item \tr{os \bnum{principais marégrafos} de cada país;}{the \bnum{main tide gauges} of
              each country;}
        \item \tr{estações de \bnum{fronteira}, que ligam as redes de nivelamento de primeira
              ordem dos vizinhos.}{\bnum{border} stations, linking the first-order leveling
              networks of neighbors.}
      \end{itemize}
      \smallskip
      \tr{A componente física usa \bnum{números geopotenciais} referidos ao $W_0$ comum,
      com uma época definida. O nivelamento e a gravimetria ajudam a ligar os datums locais.}{
      The physical component uses \bnum{geopotential numbers} referred to a common $W_0$,
      at a specified epoch. Leveling and gravimetry help connect the local datums.}
      \fonteslide{SIRGAS, \href{https://sirgas.ipgh.org/en/organization/working-groups/working-group-iii/}{Working Group III}, \tr{definição e realização vertical}{vertical definition and realization}.}
    \end{column}
    \begin{column}{0.42\textwidth}
      \begin{ideia}[title={\tr{E o Brasil?}{And Brazil?}}]
        \footnotesize\tr{O \bnum{REALT-2018} foi ajustado em \emph{números geopotenciais} justamente
        para preparar a rede ao IHRS. Imbituba e Santana são marégrafos principais. A vinculação exige expressar os dois datums em relação ao $W_0$ global.}{\bnum{REALT-2018} was adjusted in
        \emph{geopotential numbers} precisely to prepare the network for the IHRS. Imbituba and
        Santana are main tide gauges. The connection requires expressing both datums relative to the
        global $W_0$.}
      \end{ideia}
      \medskip
      \begin{regra}[title={\tr{Em metros e em $\mathrm{m^2/s^2}$}{In meters and in $\mathrm{m^2/s^2}$}}]
        \footnotesize\tr{$\num{1.30}\,\text{m}\times\num{9.78}\approx\bnum{\num{12.7}}$\,m$^{2}$/s$^{2}$.}{$\num{1.30}\,\text{m}\times\num{9.78}\approx\bnum{\num{12.7}}$\,m$^{2}$/s$^{2}$.}
      \end{regra}
    \end{column}
  \end{columns}
\end{frame}

\begin{frame}{\tr{Sua vez}{Your turn}}
  \begin{columns}[T]
    \begin{column}{0.50\textwidth}
      \begin{regra}[title={\tr{Metros e geopotencial}{Meters and geopotential}}]
        \small
        \begin{enumerate}
          \item[(a)] \tr{Quantos m$^{2}$/s$^{2}$ valem \num{0.15}\,m (a altura do NMM acima do
                datum, no exemplo anterior)? Use $\num{9.78}$.}{How many m$^{2}$/s$^{2}$ is
                \num{0.15}\,m worth (the MSL height above the datum, in the earlier example)? Use
                $\num{9.78}$.}
          \item[(b)] \tr{Um datum tem deslocamento de \num{12.7}\,m$^{2}$/s$^{2}$ em relação ao
                $W_0$. Quantos metros?}{A datum has an offset of \num{12.7}\,m$^{2}$/s$^{2}$ from
                $W_0$. How many meters?}
        \end{enumerate}
      \end{regra}
    \end{column}
    \begin{column}{0.46\textwidth}
      \onslide<2->{
      \begin{ideia}[title={\tr{Respostas}{Answers}}]
        \small\raggedright
        (a) $\num{0.15}\times\num{9.78}=\bnum{\num{1.47}}$\,m$^{2}$/s$^{2}$.
        \medskip

        (b) $\num{12.7}/\num{9.78}=\bnum{\num{1.30}\,m}$.
        \medskip

        \tr{Um degrau de \emph{metro e meio} é só \num{13} a \num{15} unidades de geopotencial:
        por isso as incertezas do IHRS vêm em m$^{2}$/s$^{2}$.}{A step of \emph{a meter and a
        half} is only \num{13} to \num{15} geopotential units: that is why IHRS uncertainties come
        in m$^{2}$/s$^{2}$.}
      \end{ideia}
      }
    \end{column}
  \end{columns}
\end{frame}

% ------------------------------------------------------------------
\begin{frame}{\tr{Nivelar, ou GNSS e modelo?}{Level, or GNSS and a model?}}
  \footnotesize
  \renewcommand{\arraystretch}{1.25}
  \begin{tabular}{@{}L{3.6cm}L{4.4cm}L{5.0cm}@{}}
    \toprule
    \textbf{\tr{Tarefa}{Task}} & \textbf{\tr{O que importa}{What matters}} & \textbf{\tr{Ferramenta}{Tool}} \\
    \midrule
    \tr{drenagem, canal, esgoto}{drainage, canal, sewer} & \tr{\bnum{diferenças pequenas} entre pontos próximos}{\bnum{small differences} between nearby points} & \tr{nivelamento a partir de RRNN próximas}{leveling from nearby benchmarks} \\
    \tr{mapa; projeto longo; ponto longe de RRNN}{map; long project; point far from benchmarks} & \tr{altitude \bnum{absoluta} no SGB}{\bnum{absolute} altitude in the SGB} & \tr{GNSS + hgeoHNOR2020 (aula do quase-geoide)}{GNSS + hgeoHNOR2020 (quasigeoid class)} \\
    \tr{obra costeira, porto, cheia}{coastal work, port, flood} & \tr{relação com o \bnum{mar local}}{relation to the \bnum{local sea}} & \tr{RMPG e documentos de níveis de referência do IBGE}{RMPG and IBGE's reference-levels documents} \\
    \bottomrule
  \end{tabular}
  \smallskip

  \begin{ideia}[title={\tr{Regra de bolso}{Rule of thumb}}]
    \footnotesize\tr{O nivelamento é bom para \emph{diferenças} entre pontos vizinhos; GNSS com modelo é bom
    para \emph{chegar} a um valor absoluto onde não há RN. Em obras sérias, usa-se os
    dois: GNSS para chegar à área, nivelamento para a obra.}{Leveling is good for
    \emph{differences} between neighboring points; GNSS with a model is good for \emph{reaching}
    an absolute value where there is no benchmark. In serious works, both are used: GNSS to get
    to the area, leveling for the work.}
  \end{ideia}
\end{frame}

\begin{frame}{\tr{A vida das RRNN}{The life of benchmarks}}
  \begin{columns}[T]
    \begin{column}{0.54\textwidth}
      \small
      \tr{Uma RN é uma pedra no meio da estrada. Obras, ampliações e vandalismo as
      \bnum{destroem}; o chão pode \bnum{afundar}. Antes de usar uma:}{A benchmark is a stone by
      the road. Roadworks, expansions and vandalism \bnum{destroy} them; the ground may
      \bnum{sink}. Before using one:}
      \begin{enumerate}
        \setlength{\itemsep}{3pt}
        \item \tr{leia a \textbf{ficha} no BDG (descrição, época, altitude);}{read its \textbf{sheet}
              in the BDG (description, epoch, altitude);}
        \item \tr{confirme \emph{no campo} que a RN existe e bate com a descrição;}{confirm
              \emph{in the field} that the benchmark exists and matches the description;}
        \item \tr{ligue-se a \bnum{duas} RRNN e compare o desnível medido com o publicado.}{tie to
              \bnum{two} benchmarks and compare the measured height difference with the published
              one.}
      \end{enumerate}
    \end{column}
    \begin{column}{0.42\textwidth}
      \begin{ideia}[title={\tr{Um teste de campo (números inventados)}{A field test (made-up numbers)}}]
        \footnotesize\tr{Publicado: $H_1=\num{12.345}$ e $H_2=\num{20.100}$\,m;
        $\Delta H=\num{7.755}$\,m. Nova diferença \emph{já corrigida}:
        $\num{7.759}$\,m; discrepância $\num{4}$\,mm.
        Com incerteza combinada ilustrativa $\sigma_r=\num{5}$\,mm,
        $|r|/\sigma_r=\num{0.8}$.
        $\sigma_r$ combina a nova medição com incertezas e covariância das RRNN.
        O limite de ida/volta de uma seção avalia outra grandeza.}{Published:
        $H_1=\num{12.345}$ and $H_2=\num{20.100}$\,m;
        $\Delta H=\num{7.755}$\,m. New \emph{corrected} difference:
        $\num{7.759}$\,m; discrepancy $\num{4}$\,mm.
        With illustrative combined uncertainty $\sigma_r=\num{5}$\,mm,
        $|r|/\sigma_r=\num{0.8}$.
        $\sigma_r$ combines the new measurement with benchmark uncertainties and covariance.
        The section forward/backward limit evaluates a different quantity.}
      \end{ideia}
    \end{column}
  \end{columns}
\end{frame}

\begin{frame}{\tr{Três datas em toda altitude}{Three dates in every altitude}}
  \centering
  \begin{tikzpicture}[x=1cm,y=1cm,font=\small]
    \node[caixa, draw=corTempo, fill=corTempo!8, text width=3.6cm, minimum height=2.2cm] (a) at (0,0)
      {\textbf{\tr{A época do datum}{The datum's epoch}}\\[3pt]
       \scriptsize\tr{o NMM de 1949--57 (Imbituba) ou 1957--58 (Santana)}{the MSL of 1949--57 (Imbituba) or 1957--58 (Santana)}};
    \node[caixa, draw=corFrame, fill=corFrame!8, text width=3.6cm, minimum height=2.2cm] (b) at (4.7,0)
      {\textbf{\tr{A versão da rede}{The network version}}\\[3pt]
       \scriptsize\tr{o ajustamento: 2011, REALT-2018, 2ª edição 2019}{the adjustment: 2011, REALT-2018, 2nd edition 2019}};
    \node[caixa, draw=corErro, fill=corErro!8, text width=3.6cm, minimum height=2.2cm] (c) at (9.4,0)
      {\textbf{\tr{O dia da medida}{The day of the measurement}}\\[3pt]
       \scriptsize\tr{o chão se mexeu desde então: mm por ano}{the ground has moved since: mm per year}};
  \end{tikzpicture}
  \bigskip

  \small
  \tr{Uma altitude honesta diz as três: ``\bnum{$H^{*}=\num{6.993}$\,m, datum de Imbituba,
  REALT-2018, medida em 2026}''. Sem isso, ela só é boa até onde as três se
  parecem.}{An honest altitude states all three: ``\bnum{$H^{*}=\num{6.993}$\,m, Imbituba datum,
  REALT-2018, measured in 2026}''. Without that, it is only as good as the three look
  alike.}
\end{frame}

\begin{frame}{\tr{As três redes do SGB}{The three networks of the SGB}}
  \centering
  \begin{tikzpicture}[x=0.88cm,y=1cm,font=\small,
      cx/.style={caixa, text width=3.6cm, minimum height=1.7cm, font=\small}]
    \node[cx, draw=corSist, fill=corSist!8] (p) at (0,2.7) {\textbf{\tr{planimétrica}{planimetric}}\\[2pt]\scriptsize\tr{SIRGAS2000, RBMC}{SIRGAS2000, RBMC}\\\scriptsize$(\varphi,\lambda,h)$};
    \node[cx, draw=corFrame, fill=corFrame!8] (a) at (6.2,2.7) {\textbf{\tr{altimétrica}{vertical}}\\[2pt]\scriptsize\tr{RRNN, REALT-2018}{benchmarks, REALT-2018}\\\scriptsize$H^{*}$};
    \node[cx, draw=corMar, fill=corMar!8] (g) at (3.1,0) {\textbf{\tr{gravimétrica}{gravimetric}}\\[2pt]\scriptsize\tr{Rede Gravimétrica, grades de $g$}{Gravity network, $g$ grids}\\\scriptsize$g$};
    \node[cx, draw=corTempo, fill=corTempo!8, minimum height=1.2cm] (m) at (10.6,2.7) {\textbf{RMPG}\\\scriptsize\tr{o zero e o mar}{the zero and the sea}};
    \draw[seta2] (p) -- node[above, font=\scriptsize] {$H^{*}\approx h-\zeta$} (a);
    \draw[seta] (g.north west) -- node[left, font=\scriptsize] {$\zeta$} (p.south);
    \draw[seta] (g.north east) -- node[right, font=\scriptsize] {$C=\sum g\,\Delta n$} (a.south);
    \draw[seta, corTempo] (m) -- (a);
  \end{tikzpicture}
  \bigskip

  \small
  \tr{As três redes se ajudam: a \bnum{gravidade} dá o geoide e o número geopotencial; o
  \bnum{GNSS} dá $h$; a \bnum{rede de nivelamento} dá $H^{*}$; e a \bnum{RMPG} amarra o zero ao
  mar.}{The three networks help each other: \bnum{gravity} gives the geoid and the geopotential
  number; \bnum{GNSS} gives $h$; the \bnum{leveling network} gives $H^{*}$; and the \bnum{RMPG}
  ties the zero to the sea.}
\end{frame}

\begin{frame}{\tr{Pare e pense}{Stop and think}}
  \begin{columns}[T]
    \begin{column}{0.46\textwidth}
      \begin{armadilha}[title={\tr{A pergunta do final}{The closing question}}]
        \small\tr{``O IBGE tem um datum em Imbituba, outro em Santana, uma rede nova de
        \num{2018} e um IHRS no horizonte. Para uma \emph{casa} em Florianópolis, nada disso
        importa, certo?''}{``IBGE has a datum at Imbituba, another at Santana, a new network of
        \num{2018} and an IHRS on the horizon. For a \emph{house} in Florianópolis, none of this
        matters, right?''}
        \medskip

        \tr{O que você responderia?}{What would you answer?}
      \end{armadilha}
    \end{column}
    \begin{column}{0.50\textwidth}
      \onslide<2->{
      \begin{ideia}[title={\tr{Resposta}{Answer}}]
        \footnotesize
        \begin{itemize}
          \setlength{\itemsep}{2pt}
          \item \tr{Para a casa \emph{sozinha}, quase nada: a obra usa a diferença com uma RN
                próxima.}{For the house \emph{alone}, almost nothing: the work uses the difference
                with a nearby benchmark.}
          \item \tr{Mas \bnum{a rede de esgoto, a rua e o rio} usam altitudes do mesmo sistema: se
                cada projeto escolhesse um zero, nada se encaixaria.}{But \bnum{the sewer, the
                street and the river} use altitudes of the same system: if each project chose a
                zero, nothing would fit.}
          \item \tr{O trabalho dos datums, ajustamentos e do REALT é para que
                \bnum{todo mundo tenha o mesmo zero} --- e para que se saiba \emph{qual}
                é.}{The work of datums, adjustments and REALT is so that \bnum{everyone has the
                same zero} --- and so that it is known \emph{which} one it is.}
        \end{itemize}
      \end{ideia}
      }
    \end{column}
  \end{columns}
\end{frame}

\begin{frame}{\tr{As quatro perguntas, respondidas}{The four questions, answered}}
  \small
  \begin{enumerate}
    \setlength{\itemsep}{3pt}
    \item \textbf{\tr{Como o IBGE vigia o mar de hoje?}{How does IBGE watch today's sea?}}
          \tr{Com a \bnum{RMPG}: seis marégrafos com sensor de nível, tempo, \bnum{GNSS} e
          nivelamento a uma RN. Isso dá o NMM local \emph{em altitude do SGB}, estação por estação,
          e o movimento do chão.}{With the \bnum{RMPG}: six tide gauges with a level sensor,
          weather, \bnum{GNSS} and leveling to a benchmark. This gives the local MSL \emph{in SGB
          altitude}, station by station, and the ground motion.}
    \item \textbf{\tr{Os dois zeros concordam?}{Do the two zeros agree?}}
          \tr{Não se sabe ao certo: \bnum{estudos} dão de \num{0.45} a \num{1.40}\,m (cerca de
          \num{1.3}\,m, Santana acima). Sem nivelamento entre as redes, só dá para
          \emph{estimar}.}{It is not known for sure: \bnum{studies} give \num{0.45} to
          \num{1.40}\,m (about \num{1.3}\,m, Santana above). With no leveling between the
          networks, it can only be \emph{estimated}.}
    \item \textbf{\tr{O que SIRGAS e IHRS propõem?}{What do SIRGAS and IHRS propose?}}
          \tr{Um zero \bnum{global}, um valor de geopotencial ($W_0$), e um quadro de estações
          (\bnum{IHRF}) com números geopotenciais e marégrafos principais; o GT III trabalha nisso
          desde 1997.}{A \bnum{global} zero, a geopotential value ($W_0$), and a frame of stations
          (\bnum{IHRF}) with geopotential numbers and main tide gauges; WG III has worked on it
          since 1997.}
    \item \textbf{\tr{E na prática?}{And in practice?}}
          \tr{Nivelar para \emph{diferenças}, GNSS com modelo para \emph{chegar}; conferir a RN no
          campo; e declarar sempre \bnum{datum, versão e data}.}{Level for \emph{differences}, GNSS
          with a model to \emph{get there}; check the benchmark in the field; and always state
          \bnum{datum, version and date}.}
  \end{enumerate}
\end{frame}

\begin{frame}{\tr{Resumo da Parte 2}{Summary of Part 2}}
  \begin{columns}[T]
    \begin{column}{0.47\textwidth}
      \begin{regra}[title={\tr{Aula 1 · O zero}{Class 1 · The zero}}]
        \scriptsize\tr{Altitude é a leitura de uma régua. O zero é um \bnum{NMM, numa
        estação, numa janela}: Torres, depois Imbituba (1949--57) e Santana (1957--58).}{Altitude
        is the reading of a ruler. The zero is \bnum{one MSL, at one station, in one window}:
        Torres, then Imbituba (1949--57) and Santana (1957--58).}
      \end{regra}
      \vspace{2pt}
      \begin{regra}[title={\tr{Aula 3 · Os ajustamentos}{Class 3 · The adjustments}}]
        \scriptsize\tr{A altitude é o \bnum{resultado de um ajustamento}: blocos $\to$
        macrocircuitos $\to$ simultâneo (2011) $\to$ REALT-2018 com $C$ e gravidade real.}{The
        altitude is the \bnum{result of an adjustment}: blocks $\to$ macro-loops $\to$
        simultaneous (2011) $\to$ REALT-2018 with $C$ and real gravity.}
      \end{regra}
    \end{column}
    \begin{column}{0.47\textwidth}
      \begin{regra}[title={\tr{Aula 2 · Lance a lance}{Class 2 · Leg by leg}}]
        \scriptsize\tr{O zero chega por \bnum{nivelamento}: $\Delta h=$ ré $-$ vante, erro
        $\eta\sqrt{K}$ (aleatório) ou $\delta K$ (sistemático), rede de RRNN, $\sum g\,\Delta n$
        fecha.}{The zero arrives by \bnum{leveling}: $\Delta h=$ back $-$ fore, error
        $\eta\sqrt{K}$ (random) or $\delta K$ (systematic), network of benchmarks, $\sum
        g\,\Delta n$ closes.}
      \end{regra}
      \vspace{2pt}
      \begin{regra}[title={\tr{Aula 4 · O zero hoje e amanhã}{Class 4 · The zero, now and next}}]
        \scriptsize\tr{A RMPG vigia o mar; os \bnum{dois zeros} diferem em cerca de
        \num{1.3}\,m; o futuro é o \bnum{IHRS}; e toda altitude tem \bnum{três datas}.}{The RMPG
        watches the sea; the \bnum{two zeros} differ by about \num{1.3}\,m; the future is the
        \bnum{IHRS}; and every altitude has \bnum{three dates}.}
      \end{regra}
    \end{column}
  \end{columns}
  \smallskip
  \centering\footnotesize
  \tr{Parte 1: \bnum{onde} (SIRGAS2000). Parte 2: \bnum{a que altura} (o zero, o nivelamento, o
  REALT).}{Part 1: \bnum{where} (SIRGAS2000). Part 2: \bnum{how high} (the zero, leveling,
  REALT).}
\end{frame}

\begin{frame}{\tr{Para levar para casa}{Take home}}
  \begin{columns}[T]
    \begin{column}{0.54\textwidth}
      \begin{regra}[title={\tr{Quatro ideias}{Four ideas}}]
        \small
        \begin{enumerate}
          \setlength{\itemsep}{3pt}
          \item \tr{A RMPG liga o \bnum{mar de hoje} ao zero de ontem.}{The RMPG ties \bnum{today's
                sea} to yesterday's zero.}
          \item \tr{Os dois zeros diferem em cerca de \bnum{\num{1.3}\,m} (estimativa, não
                oficial).}{The two zeros differ by about \bnum{\num{1.3}\,m} (an estimate, not
                official).}
          \item \tr{O futuro é um zero \bnum{global}: o IHRS, com $W_0=\num{62636853.4}$\,m$^{2}$/s$^{2}$.}{The
                future is a \bnum{global} zero: the IHRS, with $W_0=\num{62636853.4}$\,m$^{2}$/s$^{2}$.}
          \item \tr{Toda altitude diz \bnum{datum, versão e data}.}{Every altitude states
                \bnum{datum, version and date}.}
        \end{enumerate}
      \end{regra}
    \end{column}
    \begin{column}{0.42\textwidth}
      \begin{ideia}[title={\tr{Onde ler mais}{Where to read more}}]
        \small\tr{As contas de cada slide estão em \texttt{sgb\_altimetria\_numeros.py}; as
        fórmulas e a história, no apêndice. O geoide e o quase-geoide têm aulas próprias:
        \texttt{geoide\_e\_alturas}, \texttt{quase\_geoide} e \texttt{gravimetria}.}{The
        computations of each slide are in \texttt{sgb\_altimetria\_numeros.py}; the formulas and
        the history, in the appendix. The geoid and the quasigeoid have their own classes:
        \texttt{geoide\_e\_alturas}, \texttt{quase\_geoide} and \texttt{gravimetria}.}
      \end{ideia}
    \end{column}
  \end{columns}
\end{frame}

% ==================================================================
\appendix

\begin{frame}[standout]
  \tr{Apêndice Técnico}{Technical Appendix}\\[6pt]
  {\normalsize \tr{Fórmulas, contas, linha do tempo, glossário e fontes}{Formulas, computations,
  timeline, glossary and sources}}
\end{frame}

\begin{frame}{\tr{A1 · Número geopotencial e tipos de altitude}{A1 · Geopotential number and kinds of height}}
  \scriptsize
  \tr{Só as fórmulas; a dedução está nas aulas do geoide (\texttt{geoide\_e\_alturas},
  \texttt{quase\_geoide}).}{Formulas only; the derivation is in the geoid classes
  (\texttt{geoide\_e\_alturas}, \texttt{quase\_geoide}).}
  \begin{columns}[onlytextwidth]
    \begin{column}{0.48\textwidth}
      \begin{tcolorbox}[colback=mTeal!8, colframe=mTeal, title={\tr{Número geopotencial}{Geopotential number}}]
        \scriptsize
        \begin{equation*}
          C_P=W_0-W_P=\int_{0}^{P}g\,\mathrm{d}n\;\approx\;\sum_{i}\bar g_i\,\Delta n_i
        \end{equation*}
        \tr{Unidade: m$^{2}$/s$^{2}$ (1 m de altitude $\approx\num{9.8}$ m$^{2}$/s$^{2}$). Não depende do caminho; fecha em
        qualquer circuito.}{Unit: m$^{2}$/s$^{2}$ (1 m of height $\approx\num{9.8}$ m$^{2}$/s$^{2}$). Path-independent;
        closes around any loop.}
      \end{tcolorbox}
    \end{column}
    \begin{column}{0.48\textwidth}
      \begin{tcolorbox}[colback=mIndigo!8, colframe=mIndigo, title={\tr{Três altitudes a partir de $C$}{Three heights from $C$}}]
        \scriptsize
        \begin{align*}
          H^{*}&=\frac{C}{\bar\gamma}\quad\text{(\tr{normal}{normal})}\\
          H&=\frac{C}{\bar g}\quad\text{(\tr{ortométrica}{orthometric})}\\
          H^{\mathrm{dyn}}&=\frac{C}{\gamma_{45}}\quad\text{(\tr{dinâmica}{dynamic})}
        \end{align*}
      \end{tcolorbox}
    \end{column}
  \end{columns}
  \medskip
  \tr{$\bar\gamma$: gravidade \emph{normal} média ao longo da vertical normal; $\bar g$: gravidade \emph{real} média ao longo da linha de prumo; $\gamma_{45}$: gravidade normal a $45\gr$ de latitude (valor fixo).}{$\bar\gamma$: mean \emph{normal} gravity along the normal plumb line; $\bar g$: mean \emph{actual} gravity along the plumb line; $\gamma_{45}$: normal gravity at $45\gr$ latitude (fixed value).}
  \smallskip

  \tr{No Brasil, desde 30/07/2018, as altitudes oficiais são $H^{*}$ (REALT-2018).}{In Brazil, since 30/07/2018, the official heights are $H^{*}$ (REALT-2018).}
\end{frame}

\begin{frame}{\tr{A2 · Mínimos quadrados do nivelamento}{A2 · Least squares for leveling}}
  \scriptsize
  \begin{columns}[onlytextwidth]
    \begin{column}{0.48\textwidth}
      \begin{tcolorbox}[colback=mStoneLight, colframe=mStone, title={\tr{Modelo e solução}{Model and solution}}]
        \scriptsize
        \tr{Para cada seção $ij$: $\Delta h_{ij}+v_{ij}=H_j-H_i$, ou $\ell+v=A\,x$, com $A$ de linhas $(\dots,-1,\dots,+1,\dots)$ e $P=\mathrm{diag}(1/K_{ij})$ (o $\eta^{2}$ sai como fator).}{For each section $ij$: $\Delta h_{ij}+v_{ij}=H_j-H_i$, or $\ell+v=A\,x$, with $A$ having rows $(\dots,-1,\dots,+1,\dots)$ and $P=\mathrm{diag}(1/K_{ij})$ ($\eta^{2}$ factors out).}
        \begin{equation*}
          N=A^{\mathsf T}PA,\quad b=A^{\mathsf T}P\ell,\quad \hat x=N^{-1}b
        \end{equation*}
        \tr{$N$ é o \emph{laplaciano ponderado} do grafo da rede (sem as linhas e colunas dos pontos fixos).}{$N$ is the \emph{weighted Laplacian} of the network graph (without the rows and columns of fixed points).}
      \end{tcolorbox}
    \end{column}
    \begin{column}{0.48\textwidth}
      \begin{tcolorbox}[colback=mTeal!8, colframe=mTeal, title={\tr{Precisão}{Precision}}]
        \scriptsize
        \begin{equation*}
          \hat\sigma_0^{2}=\frac{v^{\mathsf T}Pv}{m-n},\qquad
          \Sigma_{\hat x}=\hat\sigma_0^{2}\,\eta^{2}N^{-1}
        \end{equation*}
        \tr{$m$ observações, $n$ incógnitas, $r=m-n$ graus de liberdade. Teste global: $r\,\hat\sigma_0^{2}\sim\chi^{2}_{r}$ se o modelo estocástico vale (``Teste Global'' do simulador).}{$m$ observations, $n$ unknowns, $r=m-n$ degrees of freedom. Global test: $r\,\hat\sigma_0^{2}\sim\chi^{2}_{r}$ if the stochastic model holds (the simulator's ``Global Test'').}
        \smallskip

        \tr{Resíduos padronizados (\emph{data snooping}): $w_i=v_i/\sigma_{v_i}$, com $Q_{vv}=P^{-1}-AN^{-1}A^{\mathsf T}$.}{Standardized residuals (\emph{data snooping}): $w_i=v_i/\sigma_{v_i}$, with $Q_{vv}=P^{-1}-AN^{-1}A^{\mathsf T}$.}
      \end{tcolorbox}
    \end{column}
  \end{columns}
  \medskip
  \tr{\textbf{Rede livre}: sem ponto fixo, $N$ é singular, com núcleo $(1,\dots,1)^{\mathsf T}$ (somar $c$ a todas as altitudes): defeito de posto \num{1} por componente conexa. \textbf{Injunção}: fixar $H_A$ elimina o defeito.}{\textbf{Free network}: with no fixed point, $N$ is singular, with null space $(1,\dots,1)^{\mathsf T}$ (adding $c$ to all altitudes): rank defect \num{1} per connected component. \textbf{Constraint}: fixing $H_A$ removes the defect.}
\end{frame}

\begin{frame}{\tr{A2 · O exemplo de 4 nós, completo}{A2 · The 4-node example, in full}}
  \scriptsize
  \begin{columns}[onlytextwidth]
    \begin{column}{0.46\textwidth}
      \renewcommand{\arraystretch}{1.25}
      \begin{tabular}{@{}lrrr@{}}
        \toprule
        \tr{seção}{section} & $\Delta h$ (m) & $K$ (km) & $v$ (mm) \\
        \midrule
        $A\to B$ & $+\num{2.350}$ & \num{4} & $-\num{1.29}$ \\
        $B\to C$ & $-\num{2.483}$ & \num{6} & $+\num{1.26}$ \\
        $C\to D$ & $+\num{1.254}$ & \num{5} & $+\num{1.02}$ \\
        $D\to A$ & $-\num{1.121}$ & \num{3} & $-\num{0.99}$ \\
        $A\to C$ & $-\num{0.133}$ & \num{5} & $-\num{0.03}$ \\
        $B\to D$ & $-\num{1.223}$ & \num{7} & $-\num{3.72}$ \\
        \bottomrule
      \end{tabular}
      \medskip

      \tr{$H_A=\num{10.000}$\,m fixo; $p=1/K$; $m=6$, $n=3$, $r=3$.}{$H_A=\num{10.000}$\,m fixed; $p=1/K$; $m=6$, $n=3$, $r=3$.}
    \end{column}
    \begin{column}{0.50\textwidth}
      \[
        N=\begin{pmatrix}
          \num{0.5595} & -\num{0.1667} & -\num{0.1429}\\
          -\num{0.1667} & \num{0.5667} & -\num{0.2000}\\
          -\num{0.1429} & -\num{0.2000} & \num{0.6762}
        \end{pmatrix}\mathrm{km}^{-1}
      \]
      \[
        b=\begin{pmatrix}\num{3.6761} & \num{1.3088} & \num{3.7831}\end{pmatrix}^{\mathsf T}\mathrm{m/km}
      \]
      \renewcommand{\arraystretch}{1.2}
      \begin{tabular}{@{}lrr@{}}
        \toprule
        & $\hat H$ (m) & $\sigma$ (mm) \\
        \midrule
        $B$ & \num{12.3487} & \num{1.55} \\
        $C$ & \num{9.8670} & \num{1.58} \\
        $D$ & \num{11.1220} & \num{1.42} \\
        \bottomrule
      \end{tabular}
      \smallskip

      \tr{$\hat\sigma_0=\num{1.03}$; $\det N=\num{0.1521}$.}{$\hat\sigma_0=\num{1.03}$; $\det N=\num{0.1521}$.}
    \end{column}
  \end{columns}
\end{frame}

\begin{frame}{\tr{A3 · Propagação de erro no nivelamento}{A3 · Error propagation in leveling}}
  \scriptsize
  \begin{columns}[onlytextwidth]
    \begin{column}{0.48\textwidth}
      \begin{tcolorbox}[colback=mTeal!8, colframe=mTeal, title={\tr{Aleatório e sistemático}{Random and systematic}}]
        \tiny
        \begin{align*}
          \sigma_{\Delta H}^{2}&=\sum_i\sigma_i^{2}=\eta^{2}\sum_iK_i=\eta^{2}K\;\Rightarrow\;\sigma=\eta\sqrt K\\
          e&=\delta K\\
          \operatorname{RMSE}^{2}&=\eta^{2}K+(\delta K)^{2},\qquad K^{*}=(\eta/\delta)^{2}
        \end{align*}
        \tr{Com $\eta=\num{1}$\,mm$/\sqrt{\text{km}}$ e $\delta=\num{0.1}$\,mm/km: $K^{*}=\num{100}$\,km; em \num{2000}\,km, $\num{44.7}$\,mm contra $\num{200}$\,mm.}{With $\eta=\num{1}$\,mm$/\sqrt{\text{km}}$ and $\delta=\num{0.1}$\,mm/km: $K^{*}=\num{100}$\,km; at \num{2000}\,km, $\num{44.7}$\,mm versus $\num{200}$\,mm.}
      \end{tcolorbox}
    \end{column}
    \begin{column}{0.48\textwidth}
      \begin{tcolorbox}[colback=mIndigo!8, colframe=mIndigo, title={\tr{Efeitos por visada}{Per-sight effects}}]
        \tiny
        \begin{align*}
          c&=(1-k)\frac{s^{2}}{2R},\quad k\approx\num{0.13}\\
          e_{\text{col}}&=\varepsilon\,s,\qquad \Delta e=\varepsilon\,(s_{\text{ré}}-s_{\text{vante}})\\
          w&\approx H\,\frac{\Delta g}{\gamma}\quad\text{(\tr{circuito}{loop})}
        \end{align*}
        \tr{$s=\num{50}$ m: $c=\num{0.20}$ mm ($\num{0.17}$ com refração); $\varepsilon=10\asec$ e $\Delta s=\num{20}$ m: \num{0.97} mm por lance; $H=\num{1000}$ m, $\Delta g=\num{20}$ mGal: $w=\num{20.4}$ mm.}{$s=\num{50}$ m: $c=\num{0.20}$ mm ($\num{0.17}$ with refraction); $\varepsilon=10\asec$ and $\Delta s=\num{20}$ m: \num{0.97} mm per setup; $H=\num{1000}$ m, $\Delta g=\num{20}$ mGal: $w=\num{20.4}$ mm.}
      \end{tcolorbox}
    \end{column}
  \end{columns}
  \medskip
  \tr{\textbf{Ida e volta}: $d=\Delta h_{\text{ida}}+\Delta h_{\text{volta}}$ (com a volta de sinal trocado); $\mathrm{Var}(d)=2\sigma_1^{2}$ se cada percurso tem $\sigma_1=\eta_1\sqrt K$. Tolerância de seção: $|d|\le 3\sqrt{K}$ mm; erro-padrão máximo da linha ajustada: $2\sqrt K$ mm (IBGE, 1983, Quadro VI, p. 30).}{\textbf{Forward and back}: $d=\Delta h_{\text{forward}}+\Delta h_{\text{back}}$ (with the return of opposite sign); $\mathrm{Var}(d)=2\sigma_1^{2}$ if each run has $\sigma_1=\eta_1\sqrt K$. Section tolerance: $|d|\le 3\sqrt{K}$ mm; maximum standard error of an adjusted line: $2\sqrt K$ mm (IBGE, 1983, Table VI, p. 30).}
\end{frame}

\begin{frame}{\tr{A4 · Linha do tempo da altimetria brasileira (1)}{A4 · Timeline of Brazilian leveling (1)}}
  \tiny
  \renewcommand{\arraystretch}{1.25}
  \begin{tabular}{@{}lL{11.8cm}@{}}
    \toprule
    \textbf{\tr{Ano}{Year}} & \textbf{\tr{Marco}{Milestone}} \\
    \midrule
    \textbf{1919--20} & \tr{Marégrafo de \textbf{Torres} (RS), da Comissão da Carta Geral: NMM de fev/1919 a fev/1920, datum provisório.}{\textbf{Torres} (RS) tide gauge, of the Comissão da Carta Geral: MSL from Feb 1919 to Feb 1920, provisional datum.} \\
    \textbf{1945} & \tr{Outubro: o CNG (IBGE) começa o nivelamento de precisão (Seção de Nivelamento; origem em Urussanga, SC).}{October: the CNG (IBGE) starts precision leveling (Leveling Section; origin at Urussanga, SC).} \\
    \textbf{1946} & \tr{Dezembro: a rede é ligada ao marégrafo de Torres.}{December: the network is tied to the Torres gauge.} \\
    \textbf{1948} & \tr{Começa a série do marégrafo de Imbituba (PSMSL: 1948--68). Primeiro ajustamento (datum de Torres).}{Imbituba gauge series begins (PSMSL: 1948--68). First adjustment (Torres datum).} \\
    \textbf{1949--57} & \tr{Janela de observação do NMM de \textbf{Imbituba} (9 anos): o datum vertical.}{MSL observation window of \textbf{Imbituba} (9 years): the vertical datum.} \\
    \textbf{1952} & \tr{Segundo ajustamento (datum de Torres).}{Second adjustment (Torres datum).} \\
    \textbf{1957--58} & \tr{Janela de observação do NMM de \textbf{Santana} (2 anos).}{MSL observation window of \textbf{Santana} (2 years).} \\
    \textbf{1958 / 1959} & \tr{Substituição em 1958 (IBGE, 2011, p. 8); adoção e terceiro ajustamento em 1959 (IBGE, 2019, p. 13). Desde 1959, ajustes no datum de Imbituba.}{Replacement in 1958 (IBGE, 2011, p. 8); adoption and third adjustment in 1959 (IBGE, 2019, p. 13). Adjustments use Imbituba from 1959 onward.} \\
    \textbf{1980--81} & \tr{Adoção de Santana em 1980; RRNN implantadas em 1980--81, com NMM de 1957--58 (IBGE, 2019, pp. 13 e 44).}{Santana adopted in 1980; benchmarks established in 1980--81, using 1957--58 MSL (IBGE, 2019, pp. 13 and 44).} \\
    \textbf{1991--92} & \tr{Macrocircuitos e o AAGP (ajustamento global preliminar); último ajustamento particionado: \num{43079} RRNN, com referência de cálculo no BDG em 01/12/1992 (IBGE, 2011, p. 12).}{Macro-loops and the AAGP (preliminary global adjustment); last partitioned adjustment: \num{43079} benchmarks, with BDG calculation reference 1 December 1992 (IBGE, 2011, p. 12).} \\
    \bottomrule
  \end{tabular}
\end{frame}

\begin{frame}{\tr{A4 · Linha do tempo da altimetria brasileira (2)}{A4 · Timeline of Brazilian leveling (2)}}
  \tiny
  \renewcommand{\arraystretch}{1.25}
  \begin{tabular}{@{}lL{11.8cm}@{}}
    \toprule
    \textbf{\tr{Ano}{Year}} & \textbf{\tr{Marco}{Milestone}} \\
    \midrule
    \textbf{1996} & \tr{A ideia da RMPG é proposta pela UFRJ e pela UFPR.}{The RMPG idea is proposed by UFRJ and UFPR.} \\
    \textbf{1997} & \tr{Criação do GT III do SIRGAS (Datum Vertical) durante a Assembleia Científica da IAG no Rio de Janeiro.}{SIRGAS WG III (Vertical Datum) is established during the IAG Scientific Assembly in Rio de Janeiro.} \\
    \textbf{2000--01} & \tr{Campanha SIRGAS2000 com GPS em marégrafos; GT III consolidado em Cartagena (2001); a Coordenação de Geodésia passa a administrar a RMPG (2001).}{SIRGAS2000 campaign with GPS at tide gauges; WG III consolidated in Cartagena (2001); the Geodesy Coordination begins to run the RMPG (2001).} \\
    \textbf{2005--11} & \tr{DALTI: ajustamento simultâneo da rede com o GHOST; divulgado em junho de 2011 ($\num{69590}$ RRNN, $>\num{74000}$ observações).}{DALTI: simultaneous adjustment of the network with GHOST; released in June 2011 ($\num{69590}$ benchmarks, $>\num{74000}$ observations).} \\
    \textbf{2010} & \tr{Manual Técnico n.º 11 do IBGE: controle geodésico de estações maregráficas (CGEM).}{IBGE Technical Manual No.~11: geodetic control of tide gauge stations (CGEM).} \\
    \textbf{2015} & \tr{Resolução n.º 1 da IAG: define o IHRS ($W_0=\num{62636853.4}$\,m$^{2}$/s$^{2}$); a realização regional deve densificar o IHRF.}{IAG Resolution No.~1: defines the IHRS ($W_0=\num{62636853.4}$\,m$^{2}$/s$^{2}$); the regional realization should densify the IHRF.} \\
    \textbf{2018--19} & \tr{30/07/2018: REALT-2018 (números geopotenciais, gravidade real, altitudes normais). 30/09/2019: 2.ª edição do relatório.}{30/07/2018: REALT-2018 (geopotential numbers, real gravity, normal heights). 30/09/2019: 2nd edition of the report.} \\
    \textbf{2020} & \tr{hgeoHNOR2020: conversão de $h$ (GNSS) em $H^{*}$ (aula do quase-geoide).}{hgeoHNOR2020: converts $h$ (GNSS) to $H^{*}$ (quasigeoid class).} \\
    \textbf{2022} & \tr{1.o de abril: a RMPG passa a publicar também os dados meteorológicos.}{1 April: the RMPG also begins to publish meteorological data.} \\
    \textbf{2024} & \tr{Série consolidada da RMPG: jul/2001 a dez/2024 (\num{23.5} anos).}{Consolidated RMPG series: Jul 2001 to Dec 2024 (\num{23.5} years).} \\
    \textbf{2026} & \tr{Referência temporal do exemplo: \num{72.5} anos desde 1953,5, o meio da janela do NMM de Imbituba. A extrapolação de tendência é ilustrativa.}{
      Example time reference: \num{72.5} years since 1953.5, the midpoint of Imbituba's MSL window. Trend extrapolation is illustrative.} \\
    \bottomrule
  \end{tabular}
\end{frame}

\begin{frame}{\tr{A5 · A conta do ciclo nodal e do nível médio}{A5 · The computation of the nodal cycle and mean level}}
  \scriptsize
  \begin{columns}[onlytextwidth]
    \begin{column}{0.48\textwidth}
      \begin{tcolorbox}[colback=mTeal!8, colframe=mTeal, title={\tr{Marés e ciclo nodal}{Tides and nodal cycle}}]
        \tiny
        \begin{align*}
          T_{M2}&=\num{12.4206}\,\mathrm{h},\quad T_{S2}=\num{12.0000}\,\mathrm{h}\\
          T_{\text{viva-morta}}&=\left(\tfrac{1}{T_{S2}}-\tfrac{1}{T_{M2}}\right)^{-1}=\num{14.77}\,\mathrm{d}\\
          T_{\text{nodal}}&=\num{6798.38}\,\mathrm{d}=\num{18.61}\,\mathrm{yr}\\
          N_{M2}&=\frac{\num{6798.38}\times24}{\num{12.4206}}=\num{13136}\\
          \delta_{\max}&=\varepsilon\pm i=\num{23.44}\gr\pm\num{5.15}\gr=\num{28.6}\gr,\ \num{18.3}\gr
        \end{align*}
      \end{tcolorbox}
    \end{column}
    \begin{column}{0.48\textwidth}
      \begin{tcolorbox}[colback=mIndigo!8, colframe=mIndigo, title={\tr{Quanto tempo de média? (simulação)}{How long to average? (simulation)}}]
        \tiny
        \tr{Modelo sintético: ciclo anual \num{8}\,cm, ciclo nodal \num{2}\,cm, ruído AR(1) $\sigma=\num{6}$\,cm, $\rho=\num{0.5}$; \num{4000} séries, semente fixa. Erro rms do NMM estimado:}{Synthetic model: annual cycle \num{8}\,cm, nodal cycle \num{2}\,cm, AR(1) noise $\sigma=\num{6}$\,cm, $\rho=\num{0.5}$; \num{4000} series, fixed seed. Rms error of the estimated MSL:}
        \smallskip

        \renewcommand{\arraystretch}{1.15}
        \begin{tabular}{@{}lr@{}}
          1 \tr{mês}{month} & \num{8.3}\,cm \\
          1 \tr{ano}{year} & \num{3.1}\,cm \\
          2 \tr{anos}{years} & \num{2.5}\,cm \\
          9 \tr{anos}{years} & \num{1.4}\,cm \\
          19 \tr{anos}{years} & \num{0.7}\,cm \\
        \end{tabular}
        \smallskip

        \tr{Parâmetros \emph{ilustrativos}, não medidos.}{\emph{Illustrative} parameters, not measured.}
      \end{tcolorbox}
    \end{column}
  \end{columns}
  \medskip
  \tr{\textbf{Subida do mar}: $\num{2.24}\ \mathrm{mm/ano}\times\num{72.5}\ \mathrm{anos}=\num{162}$\,mm ($\pm\num{29}$\,mm pela taxa de $\pm\num{0.4}$\,mm/ano). Janela de Imbituba: 1949--57, meio em \num{1953.5}; até 2026: \num{72.5} anos.}{\textbf{Sea rise}: $\num{2.24}\ \mathrm{mm/yr}\times\num{72.5}\ \mathrm{yr}=\num{162}$\,mm ($\pm\num{29}$\,mm from the rate of $\pm\num{0.4}$\,mm/yr). Imbituba window: 1949--57, middle at \num{1953.5}; to 2026: \num{72.5} years.}
\end{frame}

\begin{frame}{\tr{A6 · Glossário bilíngue (1)}{A6 · Bilingual glossary (1)}}
  \tiny
  \begin{columns}[onlytextwidth]
    \begin{column}{0.48\textwidth}
      \begin{itemize}
        \setlength{\itemsep}{2pt}
        \item \textbf{altitude / height}: \tr{distância vertical contada do datum até o ponto (a ``régua'').}{vertical distance counted from the datum to the point (the ``ruler'').}
        \item \textbf{altitude normal / normal height} ($H^{*}$): $C/\bar\gamma$.
        \item \textbf{altitude ortométrica / orthometric height} ($H$): $C/\bar g$.
        \item \textbf{número geopotencial / geopotential number} ($C$): $W_0-W_P$, m$^{2}$/s$^{2}$.
        \item \textbf{datum vertical / vertical datum}: \tr{um valor de NMM, numa estação e época, que vale altitude zero.}{one MSL value, at one station and epoch, taken as zero altitude.}
        \item \textbf{nível médio do mar (NMM) / mean sea level (MSL)}: \tr{média do nível da água ao longo de um período.}{average of the water level over a period.}
        \item \textbf{marégrafo / tide gauge}: \tr{instrumento que registra o nível do mar.}{instrument that records the sea level.}
        \item \textbf{sizígia / spring tide; quadratura / neap tide}.
        \item \textbf{ciclo nodal / nodal cycle}: \tr{\num{18.6} anos, regressão dos nodos da Lua.}{\num{18.6} years, regression of the Moon's nodes.}
      \end{itemize}
    \end{column}
    \begin{column}{0.48\textwidth}
      \begin{itemize}
        \setlength{\itemsep}{2pt}
        \item \textbf{nível de redução, zero hidrográfico / reduction level, chart datum}: \tr{plano de referência das cartas náuticas.}{reference plane of nautical charts.}
        \item \textbf{nivelamento geométrico / geometric (spirit) leveling}: \tr{medida de desníveis com nível e mira.}{measurement of height differences with level and rod.}
        \item \textbf{nível / level; mira / (leveling) rod}.
        \item \textbf{ré / backsight; vante / foresight}: $\Delta h=$ \tr{ré $-$ vante}{back $-$ fore}.
        \item \textbf{lance / setup; seção / section; linha / line}.
        \item \textbf{circuito / loop; macrocircuito / macro-loop}.
        \item \textbf{fechamento / misclosure} ($w$): \tr{soma dos desníveis de um circuito.}{sum of the height differences of a loop.}
        \item \textbf{tolerância / tolerance}: \tr{limite acima do qual se refaz a seção.}{limit above which the section is redone.}
        \item \textbf{referência de nível (RN) / benchmark (BM)}.
      \end{itemize}
    \end{column}
  \end{columns}
\end{frame}

\begin{frame}{\tr{A6 · Glossário bilíngue (2)}{A6 · Bilingual glossary (2)}}
  \scriptsize
  \begin{columns}[onlytextwidth]
    \begin{column}{0.48\textwidth}
      \begin{itemize}
        \setlength{\itemsep}{2pt}
        \item \textbf{ajustamento / adjustment}: \tr{cálculo que concilia medidas redundantes.}{computation that reconciles redundant measurements.}
        \item \textbf{mínimos quadrados / least squares}: \tr{minimiza $\sum p\,v^{2}$.}{minimizes $\sum p\,v^{2}$.}
        \item \textbf{peso / weight} ($p=1/(\eta^{2}K)$); \textbf{resíduo / residual} ($v$).
        \item \textbf{injunção / constraint}: \tr{ponto (ou valor) mantido fixo no ajuste.}{point (or value) held fixed in the adjustment.}
        \item \textbf{rede livre / free network; rede injuncionada / constrained network}.
        \item \textbf{defeito de posto / rank defect}: \tr{falta de um datum.}{lack of a datum.}
        \item \textbf{redundância / redundancy} ($r=m-n$).
        \item \textbf{fator de variância / variance factor} ($\hat\sigma_0^{2}$).
      \end{itemize}
    \end{column}
    \begin{column}{0.48\textwidth}
      \begin{itemize}
        \setlength{\itemsep}{2pt}
        \item \textbf{erro grosseiro / blunder}; \textbf{\emph{data snooping}}: \tr{teste dos resíduos padronizados.}{test of standardized residuals.}
        \item \textbf{erro aleatório / random error; sistemático / systematic}.
        \item \textbf{colimação / collimation; refração / refraction; curvatura / curvature}.
        \item \textbf{gravidade / gravity; gravimetria / gravimetry}.
        \item \textbf{quase-geoide / quasigeoid} ($\zeta$); \textbf{geoide / geoid} ($N$).
        \item \textbf{GNSS}: \tr{sistemas globais de navegação por satélite.}{global navigation satellite systems.}
        \item \textbf{subsidência / subsidence}: \tr{afundamento do terreno.}{sinking of the ground.}
        \item \textbf{época / epoch}: \tr{data a que um valor se refere.}{date to which a value refers.}
      \end{itemize}
    \end{column}
  \end{columns}
\end{frame}

\begin{frame}{\tr{A7 · Siglas e instituições (1)}{A7 · Acronyms and institutions (1)}}
  \scriptsize
  \renewcommand{\arraystretch}{1.2}
  \begin{tabular}{@{}lL{11.4cm}@{}}
    \toprule
    \textbf{\tr{Sigla}{Acronym}} & \textbf{\tr{Nome e significado}{Name and meaning}} \\
    \midrule
    \textbf{CNG} & \tr{Conselho Nacional de Geografia (órgão do IBGE, iniciou o nivelamento de precisão em 1945).}{National Geography Council (an IBGE body; started precision leveling in 1945).} \\
    \textbf{IBGE} & \tr{Instituto Brasileiro de Geografia e Estatística (autoridade geodésica do país).}{Brazilian Institute of Geography and Statistics (the country's geodetic authority).} \\
    \textbf{SGB} & \tr{Sistema Geodésico Brasileiro: redes planimétrica, altimétrica e gravimétrica.}{Brazilian Geodetic System: planimetric, vertical and gravimetric networks.} \\
    \textbf{RAAP} & \tr{Rede Altimétrica de Alta Precisão: a rede de nivelamento do SGB.}{High-Precision Altimetric Network: the SGB leveling network.} \\
    \textbf{RN, RRNN} & \tr{Referência de Nível, Referências de Nível.}{Benchmark, benchmarks.} \\
    \textbf{REALT} & \tr{Reajustamento da Rede Altimétrica com Números Geopotenciais (REALT-2018; 2.ª ed. 2019).}{Readjustment of the Altimetric Network with Geopotential Numbers (REALT-2018; 2nd ed.\ 2019).} \\
    \textbf{AAGP} & \tr{Ajustamento Altimétrico Global Preliminar (1991--92).}{Preliminary Global Altimetric Adjustment (1991--92).} \\
    \textbf{DALTI} & \tr{Equipe/projeto de altimetria da Coordenação de Geodésia do IBGE (ajuste simultâneo de 2011).}{Altimetry team/project of IBGE's Geodesy Coordination (2011 simultaneous adjustment).} \\
    \textbf{GHOST} & \emph{Geodetic adjustment using Helmert blocking Of Space and Terrestrial data} \tr{(programa canadense)}{(Canadian software)}. \\
    \textbf{BDG} & \tr{Banco de Dados Geodésicos do IBGE (fichas das estações).}{IBGE Geodetic Database (station sheets).} \\
    \bottomrule
  \end{tabular}
\end{frame}

\begin{frame}{\tr{A7 · Siglas e instituições (2)}{A7 · Acronyms and institutions (2)}}
  \tiny
  \renewcommand{\arraystretch}{1.2}
  \begin{tabular}{@{}lL{10.2cm}@{}}
    \toprule
    \textbf{\tr{Sigla}{Acronym}} & \textbf{\tr{Nome e significado}{Name and meaning}} \\
    \midrule
    \textbf{RMPG} & \tr{Rede Maregráfica Permanente para Geodésia (IBGE; seis estações em 2024).}{Permanent Tide Gauge Network for Geodesy (IBGE; six stations in 2024).} \\
    \textbf{CGEM} & \tr{Controle Geodésico de Estações Maregráficas (Manual Técnico do IBGE n.º 11, 2010).}{Geodetic Control of Tide Gauge Stations (IBGE Technical Manual No.~11, 2010).} \\
    \textbf{DHN} & \tr{Diretoria de Hidrografia e Navegação (Marinha do Brasil; cartas náuticas e zeros hidrográficos).}{Directorate of Hydrography and Navigation (Brazilian Navy; nautical charts and hydrographic zeros).} \\
    \textbf{PSMSL} & \emph{Permanent Service for Mean Sea Level} \tr{(séries de nível do mar, p.\,ex.\ Imbituba 1948--68)}{(sea level series, e.g., Imbituba 1948--68)}. \\
    \textbf{NOAA / NTDE} & \emph{National Oceanic and Atmospheric Administration} / \emph{National Tidal Datum Epoch} \tr{(época de datum de \num{19} anos)}{(\num{19}-year datum epoch)}. \\
    \textbf{IAG} & \emph{International Association of Geodesy}. \\
    \textbf{IHRS / IHRF} & \emph{International Height Reference System / Frame} \tr{(IAG, 2015; $W_0=\num{62636853.4}$\,m$^{2}$/s$^{2}$)}{(IAG, 2015; $W_0=\num{62636853.4}$\,m$^{2}$/s$^{2}$)}. \\
    \textbf{SIRGAS, GT III} & \tr{Sistema de Referência Geocêntrico para as Américas; Grupo de Trabalho III (Datum Vertical).}{Geocentric Reference System for the Americas; Working Group III (Vertical Datum).} \\
    \textbf{RBMC} & \tr{Rede Brasileira de Monitoramento Contínuo dos Sistemas GNSS (p.\,ex.\ a estação IMBI em Imbituba).}{Brazilian Network for Continuous Monitoring of GNSS (e.g., station IMBI at Imbituba).} \\
    \textbf{BNDG} & \tr{Banco Nacional de Dados Gravimétricos.}{National Gravity Database.} \\
    \textbf{MAPGEO2015, hgeoHNOR2020} & \tr{Modelo de geoide do IBGE; modelo de conversão $h\to H^{*}$ (aula do quase-geoide).}{IBGE geoid model; conversion model $h\to H^{*}$ (quasigeoid class).} \\
    \bottomrule
  \end{tabular}
\end{frame}

\begin{frame}{\tr{A8 · Fontes primárias (1)}{A8 · Primary sources (1)}}
  \scriptsize
  \begin{itemize}
    \setlength{\itemsep}{6pt}
    \item \textbf{[K]} IBGE (1983).
    \href{https://geoftp.ibge.gov.br/metodos_e_outros_documentos_de_referencia/normas/legisla_resolucao_d.pdf}{\emph{Especificações e Normas Gerais para Levantamentos Geodésicos em Território Brasileiro}}.
    R.PR 22, 21/07/1983, \tr{Quadro VI, p. 30 (PDF p. 26): controles de campo e erro-padrão após ajuste}{Table VI, p. 30 (PDF p. 26): field checks and standard error after adjustment}.

    \item \textbf{[B]} IBGE (2011).
    \href{https://geoftp.ibge.gov.br/metodos_e_outros_documentos_de_referencia/outros_documentos_tecnicos/apresentacoes_artigos/relatorioajustamento.pdf}{\emph{Ajustamento Simultâneo da Rede Altimétrica de Alta Precisão do Sistema Geodésico Brasileiro}}.
    \tr{p. 2: \num{69590} RRNN; p. 8: origem e troca de datum em 1958;
    pp. 10--12: oito ajustes, macrocircuitos e AAGP, referência 01/12/1992}{
    p. 2: \num{69590} benchmarks; p. 8: origins and datum replacement in 1958;
    pp. 10--12: eight adjustments, macro-loops and AAGP, reference 1 December 1992}.

    \item \textbf{[C]} IBGE (2019).
    \href{https://biblioteca.ibge.gov.br/visualizacao/livros/liv101666.pdf}{\emph{Reajustamento da Rede Altimétrica com Números Geopotenciais}}, \tr{2.ª ed.}{2nd ed.}
    ISBN 978-85-240-4502-8.
    \tr{p. 13: adoção de Imbituba em 1959 e Santana em 1980; p. 35: bloco principal;
    p. 36, Figura 16: frequências dos resíduos; pp. 40--44: linhas individuais e Amapá}{
    p. 13: adoption of Imbituba in 1959 and Santana in 1980; p. 35: principal block;
    p. 36, Figure 16: residual frequencies; pp. 40--44: individual lines and Amapá}.
  \end{itemize}
\end{frame}

\begin{frame}{\tr{A8 · Fontes primárias (2) e leituras}{A8 · Primary sources (2) and further reading}}
  \scriptsize
  \begin{itemize}
    \setlength{\itemsep}{5pt}
    \item \textbf{[D]} Siqueira, J.; Lima, M.; Luz, R. (2018).
    \href{https://sirgas.ipgh.org/docs/Boletines/Bol23/32_Siqueira_et_al_2018_Repro_red_vertical_Brasil.pdf}{\emph{Reajustamento da Rede Altimétrica Brasileira em Números Geopotenciais}}.
    Simposio SIRGAS2018, Boletín 23.
    \tr{slides 3--4: data dos ajustes; 9: mudança de metodologia; 20--21: interpolação;
    31--36: resultados}{slides 3--4: adjustment dates; 9: methodology change;
    20--21: interpolation; 31--36: results}.

    \item \textbf{[I]} SIRGAS,
    \href{https://sirgas.ipgh.org/en/organization/working-groups/working-group-iii/}{\emph{Working Group III (Vertical datum)}}.
    \tr{Criação na Assembleia Científica da IAG de 1997; definição compatível com o IHRS;
    realização como densificação do IHRF, com época de referência}{
    Established at the 1997 IAG Scientific Assembly; IHRS-compatible definition;
    realization as IHRF densification, with a reference epoch}.

    \item \textbf{[A, E]} IBGE,
    \href{https://www.ibge.gov.br/geociencias/informacoes-sobre-posicionamento-geodesico/rede-geodesica/16283-rede-altimetrica.html}{\emph{Rede Altimétrica}};
    \emph{Controle Geodésico de Estações Maregráficas} (2010);
    \tr{relatórios da RMPG. [J] NOAA, \emph{Tidal Datums}: época de 19 anos}{
    RMPG reports. [J] NOAA, \emph{Tidal Datums}: 19-year epoch}.

    \item \textbf{[G, H]} \tr{Taxas de NMM e subsidência, e diferença Imbituba/Santana,
    são estimativas de estudos (2013--2019), não parâmetros oficiais de transformação.
    Referências e limites da conferência em}{MSL/subsidence rates and the Imbituba/Santana
    offset are study estimates (2013--2019), not official transformation parameters.
    References and verification limits are in} \texttt{sgb\_altimetria\_fontes.md}.
  \end{itemize}
\end{frame}

\begin{frame}{\tr{A8 · Grandezas diferentes na avaliação da qualidade}{A8 · Different quantities in quality assessment}}
  \scriptsize
  \renewcommand{\arraystretch}{1.35}
  \begin{tabular}{@{}L{4.3cm}L{3.0cm}L{5.0cm}@{}}
    \toprule
    \tr{Grandeza}{Quantity} & \tr{Valor}{Value} & \tr{Interpretação}{Meaning}\\
    \midrule
    \tr{Ruído de uma observação}{Observation noise} &
      $\sigma=\eta\sqrt K$ & \tr{Modelo estocástico; $\eta$ escolhido ou estimado}{Stochastic model; $\eta$ selected or estimated}\\
    \tr{Ida/volta de seção}{Section forward/backward discrepancy} &
      $3\,\mathrm{mm}\sqrt K$ & \tr{Controle de campo, item 3.1}{Field check, item 3.1}\\
    \tr{Ida/volta de linha}{Line forward/backward discrepancy} &
      $4\,\mathrm{mm}\sqrt K$ & \tr{Controle de campo, item 3.2}{Field check, item 3.2}\\
    \tr{Fechamento de circuito}{Circuit misclosure} &
      $|w|/K\le\num{0.5}$\,mm/km & \tr{Após as correções, item 3.3}{After corrections, item 3.3}\\
    \tr{Erro-padrão da linha ajustada}{Adjusted-line standard error} &
      $2\,\mathrm{mm}\sqrt K$ & \tr{Máximo aceitável após ajuste, item 4}{Maximum acceptable after adjustment, item 4}\\
    \tr{Resíduos do REALT-2018}{REALT-2018 residuals} &
      \num{92.7}--\num{99.5}\,\% & \tr{Frequências empíricas do bloco principal}{Empirical principal-block frequencies}\\
    \bottomrule
  \end{tabular}
  \fonteslide{IBGE (1983), \tr{Quadro VI, p. 30, alta precisão fundamental}{Table VI, p. 30, fundamental high precision};
  IBGE (2019), p. 36. $K$: \tr{comprimento em km, contado uma vez}{length in km, counted once}.}
  \smallskip
  \tr{A padronização $v/\sqrt K$ depende da distância. O teste estatístico
  $v_i/\sigma_{v_i}$ do apêndice A2 é adimensional e tem outra finalidade.}{
  Normalization by $v/\sqrt K$ depends on distance. The statistical test
  $v_i/\sigma_{v_i}$ in Appendix A2 is dimensionless and serves a different purpose.}
\end{frame}

\begin{frame}{\tr{A8 · Alcance dos números e ressalvas}{A8 · Scope of the figures and qualifications}}
  \footnotesize
  \begin{itemize}
    \setlength{\itemsep}{5pt}
    \item \tr{\textbf{Torres/Imbituba}: IBGE (2011), p. 8, informa 1958; IBGE (2019), p. 13,
    informa adoção em 1959. A equivalência dos atos históricos continua inconclusiva.}{
    \textbf{Torres/Imbituba}: IBGE (2011), p. 8, gives 1958; IBGE (2019), p. 13,
    gives adoption in 1959. Whether these refer to the same historical act remains unresolved.}

    \item \tr{\textbf{Contagem}: \num{66882} RRNN é o bloco principal do REALT-2018.
    O total nacional não é inferido dessa tabela. As \num{69590} RRNN e as mais de
    \num{74000} observações referem-se ao relatório de 2011.}{
    \textbf{Counts}: \num{66882} benchmarks describes the REALT-2018 principal block.
    A national total is not inferred from that table. The \num{69590} benchmarks and
    more than \num{74000} observations belong to the 2011 report.}

    \item \tr{\textbf{Estimativas}: diferença entre data, taxa do NMM e subsidência
    mantêm a ressalva de estudos. A projeção até 2026 supõe taxa constante.}{
    \textbf{Estimates}: datum offsets, MSL trend and subsidence retain their study
    qualifications. Projection to 2026 assumes a constant rate.}

    \item \tr{\textbf{Exemplos ilustrativos}: série sintética, RRNN fictícias,
    rede de quatro nós, resíduos GNSS e comprimentos de \num{230000}/\num{370000}\,km.
    Os desenhos esquemáticos não são mapas.}{
    \textbf{Illustrative examples}: synthetic series, fictitious benchmarks, four-node
    network, GNSS residuals and \num{230000}/\num{370000}\,km lengths.
    Schematic drawings are not maps.}
  \end{itemize}
\end{frame}

\end{document}

